\chapter{Modular forms and modular curves}\label{ch:modforms}

This chapter builds the modular forms and modular curves of the paper: the curve $X_0(N)$ from which the modular parametrisation starts, the space $S_2(\Gamma_0(N))$ in which the
newform $f=f_E$ lives, its $q$-expansions and their integrality, the Petersson product that defines ``orthogonal to $f$'' and the Petersson norm $\langle f,f\rangle$, and the
Atkin-Lehner involution $w_N$ of the paper's Section 4.1. It assumes complex analysis and uses \Cref{ch:groups} (group actions, $\SL_2(\Z/N)$), \Cref{ch:rings} (congruences, the
Chinese remainder theorem, Kronecker symbols), \Cref{ch:modules} (free modules and saturation), \Cref{ch:reps} (characters of finite abelian groups), \Cref{ch:curves} (Riemann
surfaces, Riemann-Hurwitz, holomorphic differentials, Riemann-Roch, the cup product) and \Cref{ch:ellcurves} (lattices, $j$, the curves \curve{11a1} and others). Hecke operators,
newforms and twists are \Cref{ch:hecke}; the Jacobian $J_0(N)$, the modular parametrisation and the modular degree are \Cref{ch:moddeg}; Galois representations of newforms are
\Cref{ch:modgalrep}.

\section*{Core results at a glance}

Throughout, $\HH$ is the upper half plane, $q=e^{2\pi iz}$, $\mu(N)=[\SL_2(\Z):\Gamma_0(N)]$, the slash operator is $(f|_k\gamma)(z)=(\det\gamma)^{k/2}(cz+d)^{-k}f(\gamma z)$, $W_N=\sm0{-1}N0$, and
the Petersson product is unnormalised.

\begin{enumerate}[label=\textbf{\thechapter.\arabic*}]
\item \textbf{Conventions.} $\HH$, $\HH^*=\HH\cup\PP^1(\Q)$, $d\mu=dx\,dy/y^2$; $q=e^{2\pi iz}$ is never a prime here; slash operator with $(\det\gamma)^{k/2}$ (equal to Diamond-Shurman's in weight
$2$); $W_N=\sm0{-1}N0$ is the paper's $w$; Petersson product unnormalised; ``coefficientwise'' means at all $n$. \Cref{conv:modforms:conventions}.
\item \textbf{$\GL_2^+(\R)$ on $\HH$.} $\operatorname{Im}\gamma z=\det\gamma\operatorname{Im}z/\abs{cz+d}^2$, $\frac{d}{dz}\gamma z=\det\gamma/(cz+d)^2$, the cocycle $j(\gamma\gamma',z)=j(\gamma,\gamma'z)j(\gamma',z)$, $\HH=\SL_2(\R)/\mathrm{SO}(2)$, $d\mu$
invariant. \emph{(proved here)} \Cref{prop:modforms:upper-half-plane}.
\item \textbf{Element types} (definition): elliptic, parabolic, hyperbolic by $\abs{\tr\gamma}<,=,>2$; elliptic elements of $\SL_2(\Z)$ have order $3$, $4$, $6$; the fixed points of parabolic
elements of $\SL_2(\Z)$ are $\PP^1(\Q)$. \emph{(proved here)} \Cref{prop:modforms:element-types}.
\item \textbf{Fundamental domain.} $\mathcal F=\{\abs{\operatorname{Re}z}\le\frac12,\abs z\ge1\}$ meets every orbit, its interior points are inequivalent, the only nontrivial stabilisers are at $i$, $\zeta_3$,
$-\bar\zeta_3$ (orders $2,3,3$ in $\PSL_2$), $\SL_2(\Z)=\langle S,T\rangle$, $\vol\mathcal F=\pi/3$. \emph{(proved here)} \Cref{thm:modforms:fundamental-domain}.
\item \textbf{$X(1)$} is a sphere, obtained from $\mathcal F$ by gluing, with elliptic points $i$, $\zeta_3$ only. \emph{(proved here)} \Cref{cor:modforms:level-one-quotient}.
\item \textbf{Congruence subgroups} (definition): $\Gamma(N)\subseteq\Gamma_1(N)\subseteq\Gamma_0(N)$. \Cref{def:modforms:congruence-subgroups}.
\item \textbf{Indices.} $\SL_2(\Z)\to\SL_2(\Z/N)$ is onto; $[\SL_2(\Z):\Gamma(N)]=N^3\prod_{p\mid N}(1-p^{-2})$, $[\SL_2(\Z):\Gamma_1(N)]=N^2\prod(1-p^{-2})$,
$\Gamma_0(N)\backslash\SL_2(\Z)\simeq\PP^1(\Z/N)$ and $\mu(N)=N\prod_{p\mid N}(1+\frac1p)$; $\Gamma_0(N)/\Gamma_1(N)\simeq(\Z/N)^\times$. \emph{(proved here)} \Cref{prop:modforms:indices}.
\item \textbf{Fundamental domains of subgroups.} $\bigcup\gamma_i\mathcal F$ over coset representatives; $\vol(\Gamma_0(N)\backslash\HH)=\pi\mu(N)/3$. \emph{(proved here)} \Cref{prop:modforms:coset-domains}.
\item \textbf{Cusps and widths} (definition). \Cref{def:modforms:cusp-width}.
\item \textbf{Cusps as orbits.} Cusps of $\Gamma\ni-I$ are the $\langle T\rangle$-orbits on $\Gamma\backslash\SL_2(\Z)$, widths are orbit lengths, $\sum h_s=[\SL_2(\Z):\Gamma]$. \emph{(proved here)}
\Cref{prop:modforms:cusps-orbits}.
\item \textbf{Cusps of $\Gamma_0(N)$.} $\gcd(\text{denominator},N)$ is an invariant; cusps $a/c$, $c\mid N$, $a\bmod\gcd(c,N/c)$, width $N/\gcd(c^2,N)$, $\nu_\infty(N)=\sum_{c\mid N}\varphi(\gcd(c,N/c))$.
\emph{(proved here for the invariant, for $N$ squarefree and for $N=p^2$; general case sketched here; \cite{DiamondShurman}, \cite{Ogg73})} \Cref{prop:modforms:cusps-gamma0}.
\item \textbf{Elliptic points of $\Gamma_0(N)$} (definition and count): $\nu_2(N)$ and $\nu_3(N)$ are the numbers of roots of $u^2+1$ and $u^2-u+1$ modulo $N$, given by Kronecker symbols; the
fibres of $Y_0(N)\to Y(1)$ over $i$ and $\zeta_3$. \emph{(proved here)} \Cref{prop:modforms:elliptic-gamma0}.
\item \textbf{$Y(\Gamma)$ is a Riemann surface.} Proper discontinuity, Hausdorff quotient, finite cyclic stabilisers; charts $t=w^e$. \emph{(proved here for (i); (ii) sketched here;
\cite{DiamondShurman}, \cite{Shimura71})} \Cref{thm:modforms:y0-riemann-surface}.
\item \textbf{$X(\Gamma)$ is a compact Riemann surface} (with charts $q_h$ at the cusps), algebraic by Riemann existence and GAGA. \emph{(proof sketched here; \cite{DiamondShurman},
\Cref{ch:curves})} \Cref{thm:modforms:x0-compact}.
\item \textbf{Genus formula.} $X_0(N)\to X(1)$ has degree $\mu(N)$, ramified over $i$, $\zeta_3$, $\infty$, and $g(X_0(N))=1+\frac\mu{12}-\frac{\nu_2}4-\frac{\nu_3}3-\frac{\nu_\infty}2$; table of
genera ($g(X_0(11))=1$, $g(X_0(121))=6$, $g(X_0(539))=49$). \emph{(proved here, given the charts of \Cref{thm:modforms:y0-riemann-surface}(ii) and \Cref{thm:modforms:x0-compact}, and Riemann-Hurwitz from \Cref{ch:curves})} \Cref{thm:modforms:genus-formula}.
\item \textbf{Moduli.} $Y_0(N)$ classifies $(E,C)$, $C$ cyclic of order $N$; $Y_1(N)$ classifies $(E,P)$; $Y(N)$ classifies $(E,(P,Q))$, $(P,Q)$ a basis of $E[N]$ with Weil pairing
$e_N(P,Q)=\zeta_N$, where $\zeta_N=e_N(\frac1N,\frac\tau N)$ on $\C/(\Z+\Z\tau)$ is $e^{2\pi i/N}$ or its inverse according to the normalisation of $e_N$; $w_N(E,C)=(E/C,E[N]/C)$; cusps are degenerate
objects; $X_0(N)$ has a model over $\Q$. \emph{(proved here for (i) to (iii), given \Cref{ch:ellcurves}; (iv), (v) stated; \cite{DeligneRapoport73}, \cite{KatzMazur85}, \cite{DiamondShurman})}
\Cref{thm:modforms:moduli}.
\item \textbf{Slash operator} (definition) $(f|_k\gamma)(z)=(\det\gamma)^{k/2}(cz+d)^{-k}f(\gamma z)$. \Cref{def:modforms:slash}.
\item \textbf{Slash properties.} Action, scalars act by $\sgn(\lambda)^k$, $(f|_2\gamma)\,dz=\gamma^*(f\,dz)$, $y^{k/2}\abs f$ invariant. \emph{(proved here)} \Cref{lem:modforms:slash}.
\item \textbf{Modular and cusp forms} (definition): $M_k(\Gamma)\supseteq S_k(\Gamma)$, holomorphy and vanishing at all cusps. \Cref{def:modforms:modular-form}.
\item \textbf{First properties.} Odd weight vanishes for $\Gamma\ni-I$; products; $f(dz)\in M_k(\Gamma_0(Md))$ for $f\in M_k(\Gamma_0(M))$. \emph{(proved here)} \Cref{prop:modforms:mf-basic}.
\item \textbf{$q$-expansions.} $f=\sum_{n\ge0}a_nq^n$ on all of $\HH$, $a_n=\int_0^1f(x+iy)e^{-2\pi in(x+iy)}dx$, expansions at the other cusps in $q_h$, cusp forms have no constant terms, and
$f\mapsto(a_n)$ is injective. \emph{(proved here)} \Cref{prop:modforms:q-expansion}.
\item \textbf{Hecke bound.} For $f\in S_k(\Gamma_0(N))$, $y^{k/2}\abs f$ is bounded and $a_n(f)=O(n^{k/2})$. \emph{(proved here)} \Cref{prop:modforms:hecke-bound}.
\item \textbf{Ramanujan-Petersson.} For a newform and $p\nmid N$, $\abs{a_p}\le2p^{(k-1)/2}$, so $\beta_p=\bar\alpha_p$ and $\abs{\alpha_p/\beta_p}=1$. \emph{(stated; \cite{Deligne}, \cite{DiamondShurman})}
\Cref{thm:modforms:ramanujan-petersson}.
\item \textbf{Nebentypus} (definition): $M_k(\Gamma_0(N),\chi)$, diamond operators. \Cref{def:modforms:nebentypus}.
\item \textbf{Decomposition by characters.} $M_k(\Gamma_1(N))=\bigoplus_\chi M_k(\Gamma_0(N),\chi)$; $M_k(\Gamma_0(N),\chi)=0$ unless $\chi(-1)=(-1)^k$. \emph{(proved here)} \Cref{prop:modforms:nebentypus}.
\item \textbf{Eisenstein series of level one.} $G_k=2\zeta(k)E_k\in M_k(\SL_2(\Z))$ for even $k\ge4$, $E_k=1-\frac{2k}{B_k}\sum\sigma_{k-1}(n)q^n$; $E_2(-1/z)=z^2E_2(z)+\frac{6z}{\pi i}$ and $E_2^*=E_2-\frac3{\pi y}$
is invariant. \emph{(proved here for $k\ge4$; for $E_2$, proved here given the transformation law of $\eta$)} \Cref{prop:modforms:eisenstein-level-one}.
\item \textbf{$\Delta$ and $j$.} $\Delta=(E_4^3-E_6^2)/1728=\sum\tau(n)q^n\in S_{12}$ with integer coefficients, $\Delta(\Lambda_z)=(2\pi)^{12}\Delta(z)$; $j=E_4^3/\Delta=q^{-1}+744+196884q+\cdots$, $j(z)=j(\Lambda_z)$;
$\Delta=\eta^{24}$. \emph{(proved here; $\Delta=\eta^{24}$ given the transformation law of $\eta$)} \Cref{prop:modforms:delta-j}.
\item \textbf{Valence formula, level one.} $v_\infty+\frac12v_i+\frac13v_{\zeta_3}+\sum v_p=\frac k{12}$. \emph{(proved here)} \Cref{thm:modforms:valence-level-one}.
\item \textbf{Level one structure.} $M_2=0$, $M_k=\C E_k$ for $k=4,\dots,10,14$, $S_k=\Delta M_{k-12}$, dimension formula, $M_*=\C[E_4,E_6]$; zeros of $\Delta$, $E_4$, $E_6$. \emph{(proved here)}
\Cref{cor:modforms:level-one-structure}.
\item \textbf{$j$ is a bijection} $\SL_2(\Z)\backslash\HH\to\C$ and an isomorphism $X(1)\simeq\PP^1(\C)$; every $j_0$ is $j(\Lambda)$. \emph{(proved here; the isomorphism $X(1)\simeq\PP^1(\C)$ given the charts of \Cref{thm:modforms:y0-riemann-surface}(ii) and \Cref{thm:modforms:x0-compact}, and \Cref{ch:curves})} \Cref{thm:modforms:j-bijective}.
\item \textbf{The function field of $X_0(N)$.} $j_N(z)=j(Nz)$ is a meromorphic function on $X_0(N)$; $\C(X(1))=\C(j)$; the conjugates of $j_N$ over $\C(j)$ are the $\mu(N)$ distinct functions
$j(N\alpha_iz)$, $\SL_2(\Z)=\bigsqcup\Gamma_0(N)\alpha_i$, so $j_N$ has degree $\mu(N)$ over $\C(j)$ and $\C(X_0(N))=\C(j,j_N)$. \emph{(proved here, given the charts of \Cref{thm:modforms:y0-riemann-surface}(ii) and \Cref{thm:modforms:x0-compact},
and \Cref{ch:curves} for $[\C(X_0(N)):\C(j)]=\deg\pi$)} \Cref{prop:modforms:function-field}.
\item \textbf{Valence formula for $\Gamma_0(N)$.} $\sum_Pv_P(f)/e_P=k\mu(N)/12$. \emph{(proof sketched here; \cite{DiamondShurman}, \cite{Shimura71})} \Cref{thm:modforms:valence-gamma0}.
\item \textbf{Weight two forms are differentials.} $S_2(\Gamma_0(N))\simeq H^0(X_0(N),\Omega^1)$ by $f\mapsto f(z)\,dz$; $M_2$ corresponds to differentials with simple poles at the cusps, residue
$\frac h{2\pi i}b_0$; $\dim S_2=g$, $\dim M_2=g+\nu_\infty-1$. \emph{(proved here, given the charts of \Cref{thm:modforms:y0-riemann-surface}(ii) and \Cref{thm:modforms:x0-compact}, and \Cref{ch:curves})} \Cref{thm:modforms:weight-two-differentials}.
\item \textbf{Dimension formulas} for $S_k(\Gamma_0(N))$, $M_k(\Gamma_0(N))$. \emph{(proved here for $k\le2$; sketched here for $k\ge4$; \cite{DiamondShurman})} \Cref{thm:modforms:dimension-formulas}.
\item \textbf{Weight two Eisenstein series.} $E_2(z)-NE_2(Nz)\in M_2(\Gamma_0(N))\smallsetminus S_2$, with $a_0=1-N$; $M_2(\Gamma_0(p))=\C E_2^{(p)}\oplus S_2(\Gamma_0(p))$. \emph{(proved here, given the
transformation law of $\eta$)} \Cref{prop:modforms:eisenstein-weight-two}.
\item \textbf{The $\eta$-function.} Holomorphic without zeros, $\eta(z+1)=e^{\pi i/12}\eta(z)$, $\eta(-1/z)=\sqrt{z/i}\,\eta(z)$, $\eta'/\eta=\frac{\pi i}{12}E_2$. \emph{(transformation law stated;
\cite{Apostol90}; the rest proved here)} \Cref{prop:modforms:eta}.
\item \textbf{Level $11$.} $S_2(\Gamma_0(11))=\C\,\eta(z)^2\eta(11z)^2$, zeros only at the two cusps, $f_{11}|_2W_{11}=-f_{11}$; the coefficients are those of \curve{11a1}. \emph{(proved here given
the transformation law of $\eta$, except (iv): stated, \Cref{ch:moddeg}, checked for $n\le2000$)} \Cref{thm:modforms:level-eleven}.
\item \textbf{Petersson product} (definition): the unnormalised integral $\langle f,g\rangle=\int_{\Gamma_0(N)\backslash\HH}f\bar gy^k\,d\mu$. \Cref{def:modforms:petersson}.
\item \textbf{Petersson properties.} Well defined, convergent, positive definite; in weight $2$, $\langle f,g\rangle=\frac i2\int_{X_0(N)}\omega_f\wedge\overline{\omega_g}$. \emph{(proved here)} \Cref{prop:modforms:petersson}.
\item \textbf{Normalisations.} PARI's \texttt{mfpetersson} is the book's product divided by $\mu(N)$; other books divide by $\pi\mu(N)/3$; $\degphi\cdot\vol(\Lambda_E)=4\pi^2c_E^2\langle f,f\rangle$.
\Cref{conv:modforms:petersson-normalisations}, checked in \Cref{ex:modforms:petersson-numerics}.
\item \textbf{Change of group.} $\langle f,g\rangle_\Gamma=[\bar\Gamma_0(N):\bar\Gamma]\langle f,g\rangle_{\Gamma_0(N)}$; invariance under normalising $\alpha$; for $\Gamma\subset\SL_2(\Q)$ containing some $\Gamma(M)$ with
finite index, $\langle f,g\rangle_\Gamma=[\bar\Gamma:\bar\Gamma_1]^{-1}\langle f,g\rangle_{\Gamma_1}$ ($\Gamma_1\subseteq\Gamma\cap\SL_2(\Z)$ a congruence subgroup of finite index, $\bar{\ }$ the image in $\PSL_2(\R)$) is
well defined, and $\langle f|_k\alpha,g|_k\alpha\rangle_{\alpha^{-1}\Gamma\alpha}=\langle f,g\rangle_\Gamma$ for every $\alpha\in\GL_2^+(\Q)$. \emph{(proved here)} \Cref{prop:modforms:petersson-change}.
\item \textbf{Orthogonal complement.} $S_k=\C f\oplus f^\perp$ with projection $g\mapsto\frac{\langle g,f\rangle}{\langle f,f\rangle}f$. \emph{(proved here)} \Cref{prop:modforms:orthogonal-complement}.
\item \textbf{Integral forms} (definition): $S_k(\Gamma_0(N);A)$ for $A\subseteq\C$, and $S_k(\Gamma_0(N);R)=S_k(\Gamma_0(N);\Z)\otimes R$ for any ring $R$. \Cref{def:modforms:integral-forms}.
\item \textbf{Integral bases.} $S_k(\Gamma_0(N))$ and $M_k(\Gamma_0(N))$ have bases of forms with integer coefficients. \emph{(stated; \cite{Shimura71}, \cite{DiamondShurman})} \Cref{thm:modforms:integral-basis}.
\item \textbf{Sturm's bound.} For $f\in M_k(\Gamma_0(N);\mathcal O)$, $a_n(f)\in\lambda$ for $n\le k\mu(N)/12$ implies it for all $n$. \emph{(proof sketched here; \cite{Sturm87}, \cite{Stein})}
\Cref{thm:modforms:sturm}.
\item \textbf{The integral lattice.} $S_2(\Gamma_0(N);\Z)$ is free of rank $\dim S_2$ and spans; truncation at $B_2(N)$ is a split injection; $q$-expansion is injective over every ring, and over $\Z_\ell$ the
image is the $\ell$-integral forms; congruences modulo $m$ or $\lambda$ are decided by $a_n$, $n\le B_2(N)$. \emph{(proved here, given \Cref{thm:modforms:integral-basis,thm:modforms:sturm})}
\Cref{cor:modforms:integral-lattice}.
\item \textbf{Atkin-Lehner matrices} (definition): $W_d=\sm{dx}y{Nz}{dw}$, $\det W_d=d$, for $d\,\|\,N$. \Cref{def:modforms:atkin-lehner}.
\item \textbf{Atkin-Lehner involutions.} The matrix $W_d$ normalises $\Gamma_0(N)$, $W_d^2\in d\Gamma_0(N)$, $W_dW_{d'}\in\gcd(d,d')\mathcal W_{d*d'}$, where $\mathcal W_e$ is the set of Atkin-Lehner matrices for $e\,\|\,N$ and $d*d'=dd'/\gcd(d,d')^2$; commuting holomorphic involutions $w_d$ of $X_0(N)$ and of
$M_k(\Gamma_0(N))$ ($k$ even); $w_N\colon z\mapsto-1/(Nz)$; $w_d^*\omega_f=\omega_{f|_2W_d}$. \emph{(proved here; the statement on $X_0(N)$ given the charts of \Cref{thm:modforms:y0-riemann-surface}(ii) and \Cref{thm:modforms:x0-compact})} \Cref{prop:modforms:atkin-lehner}.
\item \textbf{Atkin-Lehner and pairings.} $w_d$ is a self-adjoint isometric involution of $S_k(\Gamma_0(N))$; $\dim S_2^+=g(X_0(N)/w_N)$; $w_N^*$ preserves the cup product and acts by $\varepsilon_N(f)$ on
$\omega_f$, $\overline{\omega_f}$ when $f|_2W_N=\varepsilon_N(f)f$. \emph{(proved here, given \Cref{ch:curves})} \Cref{prop:modforms:atkin-lehner-petersson}.
\item \textbf{Atkin-Lehner and cusps.} $w_N$ maps cusps of denominator invariant $c$ to $N/c$; for $N$ squarefree the $w_e$ act simply transitively on the cusps. \emph{(proved here)}
\Cref{prop:modforms:atkin-lehner-cusps}.
\end{enumerate}

\section*{Where the paper uses this chapter}

\begin{itemize}
\item \textbf{Section 1} (``let $f=f_E\in S_2(\Gamma_0(N))$ be the newform attached to it''; $\varphi_E\colon X_0(N)\to E$; the Petersson norm): \Cref{def:modforms:modular-form}, \Cref{prop:modforms:q-expansion},
\Cref{thm:modforms:x0-compact}, \Cref{thm:modforms:moduli}, \Cref{thm:modforms:genus-formula}, \Cref{thm:modforms:weight-two-differentials}, \Cref{def:modforms:petersson}, \Cref{pl:modforms:x0},
\Cref{pl:modforms:petersson}.
\item \textbf{Section 3, opening paragraph} (the congruence number: ``coefficientwise'', ``integer coefficients'', ``orthogonal to $f$''; the Hecke algebra as a subring of
$\End S_2(\Gamma_0(N))$): \Cref{prop:modforms:q-expansion}, \Cref{def:modforms:integral-forms}, \Cref{prop:modforms:orthogonal-complement}, \Cref{pl:modforms:coefficientwise}.
\item \textbf{Lemma lem:crit} (congruence criterion: $f^\perp$, $S_2(\Gamma_0(N);\Z)$, the version over $\Z_\ell$): \Cref{prop:modforms:orthogonal-complement}, \Cref{thm:modforms:integral-basis},
\Cref{cor:modforms:integral-lattice}, \Cref{pl:modforms:crit}; the perfect pairing with $\T$ is \Cref{ch:hecke}'s, the lemma is proved in \Cref{ch:modgalrep}.
\item \textbf{Lemma lem:old} (oldforms): the forms $h(dz)$, \Cref{prop:modforms:mf-basic}(e); orthogonality, \Cref{def:modforms:petersson}.
\item \textbf{The optimal-level theorem in Section 3} (``nebentypus of order prime to $\ell$''): \Cref{def:modforms:nebentypus}, \Cref{prop:modforms:nebentypus}, \Cref{pl:modforms:nebentypus}.
\item \textbf{Example ex:qual} (the four congruences, proved for all $n$ by Sturm's bound): \Cref{thm:modforms:sturm}, \Cref{cor:modforms:integral-lattice}, \Cref{ex:modforms:paper-congruences},
\Cref{pl:modforms:example}.
\item \textbf{Section 4.1} (the operator $w=W_{M^\Sigma}$, ``the Atkin-Lehner involution $w_{M^\Sigma}$ with no scalar factor''; $H^1(X_0(M^\Sigma),\Z_\ell)$; the step on $w_N$ in Proposition prop:pullback):
\Cref{def:modforms:atkin-lehner}, \Cref{prop:modforms:atkin-lehner}, \Cref{prop:modforms:atkin-lehner-petersson}, \Cref{pl:modforms:wN}.
\item \textbf{Remark rem:degree} ($\degphi\cdot\vol(\Lambda_E)=4\pi^2c_E^2\langle f,f\rangle$, unnormalised product): \Cref{conv:modforms:petersson-normalisations}, \Cref{ex:modforms:petersson-numerics}.
\item \textbf{Theorem thm:twist and Step 1 of the proof of Theorem thm:quant} (integrality of twists, $S_2(\Gamma_0(N);\Z_\ell)$, ``the congruence survives''): \Cref{def:modforms:integral-forms},
\Cref{cor:modforms:integral-lattice}, \Cref{pl:modforms:crit}, \Cref{pl:modforms:petersson}.
\item \textbf{Lemma lem:euler and Theorem thm:dfg14} ($\abs{\alpha/\beta}=1$ at good primes): \Cref{thm:modforms:ramanujan-petersson}, \Cref{pl:modforms:ramanujan}.
\item \textbf{Conjecture conj:eis and Section sec:eis} (the Eisenstein case): \Cref{prop:modforms:eisenstein-weight-two}, \Cref{ex:modforms:eisenstein-congruence}, \Cref{pl:modforms:eisenstein}.
\item \textbf{Definition def:terms} ($v_q(j_E)$): \Cref{prop:modforms:delta-j}, \Cref{thm:modforms:j-bijective}, \Cref{pl:modforms:ch9}.
\item \textbf{The running example \curve{11a1}}: \Cref{thm:modforms:level-eleven}, \Cref{ex:modforms:eisenstein-congruence}, \Cref{ex:modforms:petersson-numerics}.
\end{itemize}

\section{The upper half plane}\label{sec:modforms:upper-half-plane}

This section fixes the stage: the upper half plane, the action of $\GL_2^+(\R)$ by M\"obius transformations, and the invariant
measure. Nothing here is specific to the paper, but everything later rests on it: the curve $X_0(N)$ of the paper's Section 1 is a
quotient of $\HH$, the newform $f$ is a function on $\HH$, and the Petersson norm $\langle f,f\rangle$ is an integral against the measure
$d\mu$ defined here. We also fix the conventions of the chapter. The reader with a physics background will meet three old friends
on the way: $q$-expansions are Fourier series (\Cref{prop:modforms:q-expansion}), the valence formula is the argument principle
(\Cref{thm:modforms:valence-level-one}), and the Petersson product is a weighted $L^2$ inner product (\Cref{def:modforms:petersson}).

\begin{convention}[notation of this chapter]\label{conv:modforms:conventions}
\leavevmode
\begin{enumerate}[label=(\alph*)]
\item \emph{The upper half plane.} $\HH=\{z=x+iy\in\C:y>0\}$, as in \Cref{ch:curves,ch:ellcurves}; $\HH^*=\HH\cup\PP^1(\Q)$ with
$\PP^1(\Q)=\Q\cup\{\infty\}$; $d\mu=dx\,dy/y^2$. The variable of $q$-expansions is $q=e^{2\pi iz}$, and $q_h=e^{2\pi iz/h}$ for $h\ge1$. In
this chapter primes are called $p$ or $\ell$ and $q$ is always $e^{2\pi iz}$, except inside the environments ``In the paper'', which say
so when they use the paper's $q$ for a bad prime (\Cref{ch:notation}, B-C1).
\item \emph{Matrices.} $S=\sm0{-1}10$, $T=\sm1101$, and $\zeta_3=e^{2\pi i/3}$. The letter $T$ for this matrix is local to
\Cref{sec:modforms:upper-half-plane,sec:modforms:modular-group} and the exercises; elsewhere in the book $T$ is $T_\ell E$ (\Cref{ch:notation},
B-C15). For $\gamma=\sm abcd$ we write $j(\gamma,z)=cz+d$ (not to be confused with the $j$-invariant, which always appears as $j(z)$
or $j(\Lambda)$).
\item \emph{Indices.} For a subgroup $\Gamma\subseteq\SL_2(\Z)$ of finite index, $\bar\Gamma$ is its image in $\PSL_2(\Z)=\SL_2(\Z)/\{\pm I\}$. The degree
of $X(\Gamma)\to X(1)$ (\Cref{thm:modforms:genus-formula}) is $[\PSL_2(\Z):\bar\Gamma]$, which equals $[\SL_2(\Z):\Gamma]$ when $-I\in\Gamma$ (always the case for
$\Gamma_0(N)$) and half of it otherwise.
\item \emph{The slash operator.} For $\gamma\in\GL_2^+(\R)$, $k\in\Z$ and $f\colon\HH\to\C$,
\[
(f|_k\gamma)(z)=(\det\gamma)^{k/2}(cz+d)^{-k}f(\gamma z).
\]
Diamond and Shurman \cite{DiamondShurman} use $(\det\gamma)^{k-1}$ instead; the two agree in weight $2$, the only weight the paper
uses, and differ by a positive scalar otherwise. The choice $(\det\gamma)^{k/2}$ makes every Atkin-Lehner operator an exact
involution and an isometry in every even weight (\Cref{prop:modforms:atkin-lehner,prop:modforms:atkin-lehner-petersson}).
\item \emph{Atkin-Lehner.} $W_N=\sm0{-1}N0$. This is exactly the paper's matrix $w$ of Section 4.1, with $M^\Sigma$ in place of $N$. The
symbol $w_N$ denotes both the involution $z\mapsto-1/(Nz)$ of $X_0(N)$ and the operator $f\mapsto f|_kW_N$ on forms. Some books use
$\sm01{-N}0=-W_N$; since $-I$ acts trivially on $\HH$ and on forms of even weight, nothing changes.
\item \emph{The Petersson product} $\langle f,g\rangle=\int_{\Gamma_0(N)\backslash\HH}f\bar g\,y^k\,d\mu$ is \emph{unnormalised} and linear in the first
variable (\Cref{def:modforms:petersson}, \Cref{conv:modforms:petersson-normalisations}, \Cref{ch:notation}, B-C8).
\item \emph{Coefficients.} $a_n(f)$ is the $n$-th coefficient of the expansion of $f$ at $\infty$. ``$f\equiv g\pmod{\mathfrak a}$ coefficientwise'' means
$a_n(f)-a_n(g)\in\mathfrak a$ for \emph{all} $n\ge0$; a congruence ``at all primes $p\notin\Sigma$'' is a different and weaker statement,
and the two are always distinguished, as the paper distinguishes them in Lemma lem:crit and Lemma lem:old.
\item \emph{Numbers attached to $N$.} $\mu(N)=[\SL_2(\Z):\Gamma_0(N)]$; $\nu_2(N)$, $\nu_3(N)$ the numbers of elliptic points of order $2$ and
$3$ of $\Gamma_0(N)$; $\nu_\infty(N)$ its number of cusps; $g(N)$ the genus of $X_0(N)$; $B_k(N)=\lfloor k\mu(N)/12\rfloor$ the Sturm bound, so
$B_2(N)=\lfloor\mu(N)/6\rfloor$.
\item \emph{What is not here.} Hecke operators, eigenforms, newforms, oldforms and twists of forms are \Cref{ch:hecke}'s; this chapter
uses them only in examples and in the environments ``In the paper'', with a forward reference. When $\dim S_2(\Gamma_0(N))=1$ the
unique form with $a_1=1$ is called \emph{the normalised generator}. The Jacobian $J_0(N)$, the modular parametrisation and the
Manin constant are \Cref{ch:moddeg}'s.
\end{enumerate}
\end{convention}

\begin{proposition}[the action of $\GL_2^+(\R)$ on $\HH$]\label{prop:modforms:upper-half-plane}
The group $\GL_2^+(\R)$ of real $2\times2$ matrices of positive determinant acts on $\HH$ by $\gamma z=\frac{az+b}{cz+d}$, and the scalar matrices act
trivially. For $\gamma=\sm abcd\in\GL_2^+(\R)$ and $z=x+iy\in\HH$:
\begin{enumerate}[label=(\roman*)]
\item $\operatorname{Im}(\gamma z)=\det\gamma\cdot\operatorname{Im}z/\abs{cz+d}^2$; in particular $\gamma z\in\HH$;
\item $\frac{d}{dz}(\gamma z)=\det\gamma/(cz+d)^2$;
\item $j(\gamma\gamma',z)=j(\gamma,\gamma'z)\,j(\gamma',z)$ for $\gamma,\gamma'\in\GL_2^+(\R)$;
\item $\PSL_2(\R)=\SL_2(\R)/\{\pm I\}$ acts faithfully and transitively on $\HH$, the stabiliser of $i$ in $\SL_2(\R)$ is $\mathrm{SO}(2)$, and
$\SL_2(\R)/\mathrm{SO}(2)\to\HH$, $g\mapsto gi$, is a bijection;
\item the measure $d\mu=dx\,dy/y^2$ is $\GL_2^+(\R)$-invariant: $\int_{\gamma U}\phi\,d\mu=\int_U\phi(\gamma z)\,d\mu(z)$ for measurable $U\subseteq\HH$ and $\phi\ge0$.
\end{enumerate}
\end{proposition}

\begin{proof}
(i) $\gamma z=\frac{(az+b)(c\bar z+d)}{\abs{cz+d}^2}$, and the imaginary part of the numerator $ac\abs z^2+adz+bc\bar z+bd$ is $(ad-bc)y$.
(ii) $\frac{d}{dz}\frac{az+b}{cz+d}=\frac{a(cz+d)-c(az+b)}{(cz+d)^2}=\frac{\det\gamma}{(cz+d)^2}$.
(iii) If $\gamma'=\sm{a'}{b'}{c'}{d'}$, the bottom row of $\gamma\gamma'$ is $(ca'+dc',\,cb'+dd')$, and
$j(\gamma,\gamma'z)j(\gamma',z)=\bigl(c\frac{a'z+b'}{c'z+d'}+d\bigr)(c'z+d')=(ca'+dc')z+(cb'+dd')$. The same computation for the top row shows
$(\gamma\gamma')z=\gamma(\gamma'z)$, so we have an action; $\lambda I$ acts as the identity.
(iv) $\sm{\sqrt y}{x/\sqrt y}0{1/\sqrt y}\in\SL_2(\R)$ sends $i$ to $x+iy$, so the action is transitive. If $\gamma i=i$ with $\gamma\in\SL_2(\R)$ then $ai+b=-c+di$, so
$a=d$, $b=-c$ and $a^2+c^2=1$: $\gamma$ is a rotation, and conversely rotations fix $i$. If $\gamma z=z$ for all $z$, then $cz^2+(d-a)z-b=0$ identically, so
$c=b=0$, $a=d$, and $\gamma=\pm I$. The bijection is orbit-stabiliser (\Cref{ch:groups}).
(v) $z\mapsto\gamma z$ is holomorphic, so its real Jacobian is $\abs{d(\gamma z)/dz}^2=(\det\gamma)^2/\abs{cz+d}^4$ by (ii), while
$\operatorname{Im}(\gamma z)^{-2}=\abs{cz+d}^4/((\det\gamma)^2y^2)$ by (i). The product is $y^{-2}$, which is the change of variables formula for $d\mu$.
\end{proof}

The metric $ds^2=(dx^2+dy^2)/y^2$ is the hyperbolic metric, $d\mu$ is its area element, and the M\"obius maps of $\SL_2(\R)$ are its
orientation-preserving isometries. We will not need the metric, only the measure.

\begin{proposition}[elliptic, parabolic and hyperbolic elements]\label{prop:modforms:element-types}
Let $\gamma\in\SL_2(\R)$, $\gamma\ne\pm I$. Call $\gamma$ \emph{elliptic} if $\abs{\tr\gamma}<2$, \emph{parabolic} if $\abs{\tr\gamma}=2$ and \emph{hyperbolic} if
$\abs{\tr\gamma}>2$. Acting on $\HH\cup\PP^1(\R)$:
\begin{enumerate}[label=(\roman*)]
\item an elliptic $\gamma$ has exactly one fixed point in $\HH$ (and its conjugate in the lower half plane); a parabolic $\gamma$ has exactly one fixed point,
which lies in $\PP^1(\R)$; a hyperbolic $\gamma$ has two fixed points, both in $\PP^1(\R)$;
\item if $\gamma\in\SL_2(\Z)$ is elliptic, then $\tr\gamma\in\{0,\pm1\}$, and $\gamma$ has order $4$ (if $\tr\gamma=0$), $3$ (if $\tr\gamma=-1$) or $6$ (if $\tr\gamma=1$);
its image in $\PSL_2(\Z)$ has order $2$ or $3$;
\item the fixed points of the parabolic elements of $\SL_2(\Z)$ are exactly the points of $\PP^1(\Q)$.
\end{enumerate}
\end{proposition}

\begin{proof}
(i) If $c\ne0$, the fixed points in $\C$ are the roots of $cz^2+(d-a)z-b=0$, whose discriminant is $(d-a)^2+4bc=(a+d)^2-4$; so there are two complex
conjugate roots, one in $\HH$, if $\abs{\tr\gamma}<2$, one real double root if $\abs{\tr\gamma}=2$, and two real roots if $\abs{\tr\gamma}>2$, and $\infty$ is not
fixed. If $c=0$ then $\gamma=\sm a b0{a^{-1}}$ fixes $\infty$; $\abs{\tr\gamma}=\abs{a+a^{-1}}\ge2$, with equality iff $a=\pm1$, in which case $b\ne0$ and $\infty$ is the
only fixed point; otherwise $b/(a^{-1}-a)$ is a second, real, fixed point.
(ii) The trace is an integer of absolute value less than $2$. The characteristic polynomial is $X^2+1$, $X^2+X+1$ or $X^2-X+1$, which divide $X^4-1$,
$X^3-1$ and $X^6-1$ (more precisely $X^3+1$); by Cayley-Hamilton $\gamma^2=-I$, $\gamma^3=I$, $\gamma^3=-I$ respectively, and no smaller power is $\pm I$
since $\gamma\ne\pm I$ and the eigenvalues $\pm i$, $\zeta_3^{\pm1}$, $-\zeta_3^{\pm1}$ are not real.
(iii) A parabolic $\gamma\in\SL_2(\Z)$ fixes $\infty$ (if $c=0$) or the double root $(a-d)/(2c)\in\Q$. Conversely, $a/c\in\Q$ in lowest terms is $\sigma\infty$ for
some $\sigma=\sm abcd\in\SL_2(\Z)$ (solve $ad-bc=1$), and $\infty=\sigma\infty$ for $\sigma=I$; then $\sigma T\sigma^{-1}$ has trace $2$, is not $I$, and fixes $\sigma\infty$.
\end{proof}

\section{The modular group and its fundamental domain}\label{sec:modforms:modular-group}

The modular group $\SL_2(\Z)$ acts on $\HH$ with a very simple fundamental domain. This section proves the classical facts: the
domain $\mathcal F$, the identifications on its boundary, the generation of $\SL_2(\Z)$ by $S$ and $T$, and the area $\pi/3$. They are used three
times later: the quotient $X(1)$ is a sphere (the base of the genus formula), the valence formula is a contour integral around
$\partial\mathcal F$, and the fundamental domains of $\Gamma_0(N)$ over which the Petersson product is integrated are unions of translates of $\mathcal F$.
\Cref{ch:ellcurves} uses this section through the statement that $\SL_2(\Z)\backslash\HH$ classifies elliptic curves over $\C$.

\begin{theorem}[the fundamental domain of $\SL_2(\Z)$]\label{thm:modforms:fundamental-domain}
Let $\mathcal F=\{z\in\HH:\abs{\operatorname{Re}z}\le\frac12,\ \abs z\ge1\}$ and let $G=\langle S,T\rangle\subseteq\SL_2(\Z)$.
\begin{enumerate}[label=(\roman*)]
\item Every $z\in\HH$ is $G$-equivalent to a point of $\mathcal F$.
\item If two distinct points $z,z'$ of $\mathcal F$ are $\SL_2(\Z)$-equivalent, then either $\operatorname{Re}z=\pm\frac12$ and $z'=z\mp1$, or $\abs z=1$ and $z'=-1/z$.
\item The stabiliser of $z\in\mathcal F$ in $\PSL_2(\Z)$ is trivial, except for $z=i$ (generated by $S$, of order $2$), $z=\zeta_3$ (generated by $ST$, of order
$3$) and $z=-\bar\zeta_3=e^{\pi i/3}$ (generated by $TS$, of order $3$).
\item $G=\SL_2(\Z)$, and $S^2=(ST)^3=-I$.
\item $\int_{\mathcal F}d\mu=\pi/3$.
\end{enumerate}
\end{theorem}

\begin{proof}
(i) For $\gamma\in\SL_2(\Z)$, $\operatorname{Im}(\gamma z)=y/\abs{cz+d}^2$ by \Cref{prop:modforms:upper-half-plane}(i). Since $\abs{cz+d}^2=(cx+d)^2+c^2y^2$, there are only
finitely many pairs $(c,d)\in\Z^2$ with $\abs{cz+d}\le1$; so $\operatorname{Im}(\gamma z)$, $\gamma\in G$, takes a largest value. Choose $\gamma\in G$ realising it and $n\in\Z$
with $\abs{\operatorname{Re}(T^n\gamma z)}\le\frac12$, and put $z'=T^n\gamma z$; then $\operatorname{Im}z'=\operatorname{Im}\gamma z$ is still maximal. If $\abs{z'}<1$, then
$\operatorname{Im}(Sz')=\operatorname{Im}(-1/z')=\operatorname{Im}z'/\abs{z'}^2>\operatorname{Im}z'$, contradicting maximality, as $ST^n\gamma\in G$. So $z'\in\mathcal F$.

(ii) and (iii). Let $z\in\mathcal F$ and $\gamma=\sm abcd\in\SL_2(\Z)$, $\gamma\ne\pm I$, with $\gamma z\in\mathcal F$. Exchanging $(z,\gamma)$ with $(\gamma z,\gamma^{-1})$ if
necessary we may assume $\operatorname{Im}\gamma z\ge\operatorname{Im}z$, that is, $\abs{cz+d}\le1$. As $\operatorname{Im}z\ge\sqrt3/2$, $\abs c\sqrt3/2\le\abs{cz+d}\le1$, so
$c\in\{0,\pm1\}$.
If $c=0$ then $a=d=\pm1$ and $\gamma z=z+b'$ with $b'=\pm b\ne0$; both $z$ and $z+b'$ have real part in $[-\frac12,\frac12]$, so $b'=\pm1$ and
$\operatorname{Re}z=-b'/2$. This is the first case of (ii), and $\gamma$ fixes no point of $\mathcal F$.
If $c=\pm1$, replace $\gamma$ by $-\gamma$ (the same transformation) to get $c=1$, so $\abs{z+d}\le1$. If $d\ne0$ then $\abs{x+d}\ge\frac12$ and
$\abs{z+d}^2=(x+d)^2+y^2\ge\frac14+\frac34=1$, with equality only for $d=1$, $z=\zeta_3$ or $d=-1$, $z=-\bar\zeta_3$. If $d=0$ then $\abs z\le1$, so $\abs z=1$.
\begin{itemize}
\item $c=1$, $d=0$: then $b=-1$ and $\gamma z=a-1/z=a-\bar z$ (as $\abs z=1$). For $\gamma z\in\mathcal F$ we need $\abs{a-x}\le\frac12$ with $\abs x\le\frac12$: either
$a=0$ and $\gamma z=-\bar z=-1/z$ (the second case of (ii); a fixed point only for $z=i$, and then $\gamma=S$), or $a=1$, $x=\frac12$, $z=-\bar\zeta_3$ and
$\gamma z=z$ ($\gamma=\sm1{-1}10=TS$), or $a=-1$, $x=-\frac12$, $z=\zeta_3$ and $\gamma z=z$ ($\gamma=\sm{-1}{-1}10=(ST)^2$).
\item $c=d=1$, $z=\zeta_3$: then $a-b=1$ and $\gamma z=a-\frac1{\zeta_3+1}=a-e^{-\pi i/3}=a-\frac12+\frac{\sqrt3}2i$, which lies in $\mathcal F$ for $a=0$ ($\gamma=\sm0{-1}11=ST$,
fixing $\zeta_3$) and $a=1$ ($\gamma=\sm1011$, sending $\zeta_3$ to $-\bar\zeta_3=\zeta_3+1$, an instance of the first case of (ii)).
\item $c=1$, $d=-1$, $z=-\bar\zeta_3$: symmetric to the previous case (conjugate by $z\mapsto-\bar z$), giving $TS$ and the element sending $-\bar\zeta_3$
to $\zeta_3$.
\end{itemize}
Collecting the cases gives (ii) and the list (iii); the orders follow from $S^2=-I$ and $(ST)^3=-I$, checked by multiplying out:
$ST=\sm0{-1}11$, $(ST)^2=\sm{-1}{-1}10$, $(ST)^3=-I$, and $TS=T(ST)T^{-1}$.

(iv) Let $\gamma\in\SL_2(\Z)$ and $z_0=2i$, an interior point of $\mathcal F$ with trivial stabiliser in $\PSL_2(\Z)$ by (iii). By (i) there is $g\in G$ with
$g\gamma z_0\in\mathcal F$. By (ii), the equivalent points $z_0$ and $g\gamma z_0$ of $\mathcal F$ coincide (an interior point is equivalent to no other point of $\mathcal F$), so
$g\gamma=\pm I$ by (iii), and $\gamma=\pm g^{-1}\in G$ since $-I=S^2\in G$.

(v) $\int_{\mathcal F}d\mu=\int_{-1/2}^{1/2}\int_{\sqrt{1-x^2}}^\infty\frac{dy}{y^2}\,dx=\int_{-1/2}^{1/2}\frac{dx}{\sqrt{1-x^2}}=2\arcsin\frac12=\frac\pi3$.
\end{proof}

For instance, $z=(17+i)/40$ is moved into $\mathcal F$ by $T^{-3}ST^2S$: $Sz=-40/(17+i)=-\frac{68}{29}+\frac4{29}i$, $T^2Sz=-\frac{10}{29}+\frac4{29}i$, which has absolute
value less than $1$, $ST^2Sz=\frac52+i$, and $T^{-3}ST^2Sz=-\frac12+i$, a point of the left edge of $\mathcal F$.

\begin{corollary}[the level one quotient]\label{cor:modforms:level-one-quotient}
The quotient $\SL_2(\Z)\backslash\HH$ is obtained from $\mathcal F$ by gluing the two vertical sides by $T$ and folding the arc $\{\abs z=1\}\cap\mathcal F$ at $i$ by $S$.
Adding the single orbit $\PP^1(\Q)=\SL_2(\Z)\infty$ gives $X(1)=\SL_2(\Z)\backslash\HH^*$, which is homeomorphic to a sphere. The only points of $\SL_2(\Z)\backslash\HH$ with
nontrivial stabiliser are the orbits of $i$ and of $\zeta_3$ (which contains $-\bar\zeta_3$), with stabilisers of orders $2$ and $3$ in $\PSL_2(\Z)$.
\end{corollary}

\begin{proof}
The description of the identifications is \Cref{thm:modforms:fundamental-domain}(i), (ii); the stabilisers are (iii), and $-\bar\zeta_3=T\zeta_3$. That
$\PP^1(\Q)$ is a single orbit is the transitivity of $\SL_2(\Z)$ on fractions in lowest terms (\Cref{prop:modforms:element-types}(iii)). For the topology, compactify
$\mathcal F$ by the point $\infty$, with the neighbourhoods $\{\operatorname{Im}z>C\}\cup\{\infty\}$ (\Cref{thm:modforms:x0-compact}); the result is a closed disc, a curvilinear
triangle with vertices $\zeta_3$, $-\bar\zeta_3$, $\infty$. The gluings give a cell decomposition of $X(1)$ with three vertices (the class of $\zeta_3$ and $-\bar\zeta_3$,
the point $i$, the point $\infty$), two edges (the class of the vertical sides, the class of the two halves of the arc) and one face. Its Euler characteristic
is $3-2+1=2$, and a closed orientable surface of Euler characteristic $2$ is a sphere.
\end{proof}

So $X(1)$ has genus $0$; this is used in the genus formula (\Cref{thm:modforms:genus-formula}) and does not depend on the function $j$, which will give an
explicit isomorphism $X(1)\simeq\PP^1(\C)$ in \Cref{thm:modforms:j-bijective}.

\section{Congruence subgroups and their indices}\label{sec:modforms:congruence-subgroups}

The paper works with $\Gamma_0(N)$ throughout: $S_2(\Gamma_0(N))$, $X_0(N)$, $J_0(N)$. This section defines the three standard congruence subgroups,
computes their indices, and builds a fundamental domain for $\Gamma_0(N)$ from the translates of $\mathcal F$. The index $\mu(N)=[\SL_2(\Z):\Gamma_0(N)]$ is the
degree of $X_0(N)\to X(1)$, enters the genus formula, the dimension formulas and the Sturm bound, and is the factor by which PARI's Petersson product differs
from the book's (\Cref{conv:modforms:petersson-normalisations}).

\begin{definition}[congruence subgroups]\label{def:modforms:congruence-subgroups}
For $N\ge1$ let
\[
\Gamma(N)=\ker\bigl(\SL_2(\Z)\to\SL_2(\Z/N)\bigr),\qquad\Gamma_1(N)=\Bigl\{\gamma\equiv\sm1*01\bmod N\Bigr\},
\]
\[
\Gamma_0(N)=\Bigl\{\gamma\equiv\sm**0*\bmod N\Bigr\},
\]
so that $\Gamma(N)\subseteq\Gamma_1(N)\subseteq\Gamma_0(N)\subseteq\SL_2(\Z)$. A \emph{congruence subgroup} is a subgroup of $\SL_2(\Z)$ containing $\Gamma(N)$ for some $N$, its
\emph{level} being the least such $N$.
\end{definition}

Why these three groups? Each is the stabiliser of a piece of level-$N$ structure on the lattice $\Lambda_\tau=\Z+\Z\tau$, as \Cref{thm:modforms:moduli} makes precise:
$\Gamma(N)$ fixes the basis $(\frac1N,\frac\tau N)$ of $\frac1N\Lambda_\tau/\Lambda_\tau$ (a full level structure), $\Gamma_1(N)$ fixes the point $\frac1N$, and $\Gamma_0(N)$ only the
cyclic subgroup $\langle\frac1N\rangle$ it generates. The weakest structure gives the largest group and the smallest curve, and it is the one the paper needs. A small example:
for $N=2$ the units modulo $2$ are trivial, so $\Gamma_0(2)=\Gamma_1(2)$. At $N=3$ the matrix $\sm1031$ lies in $\Gamma_1(3)$, while $\sm{-1}03{-1}$ lies in
$\Gamma_0(3)$ but not in $\Gamma_1(3)$ ($d=-1\not\equiv1\pmod3$); neither lies in $\Gamma(3)$, whose elements are $\equiv I$ entrywise.

\begin{proposition}[indices]\label{prop:modforms:indices}
Let $N\ge1$.
\begin{enumerate}[label=(\roman*)]
\item Reduction $\SL_2(\Z)\to\SL_2(\Z/N)$ is surjective. Hence $\Gamma(N)$ is normal of index $\#\SL_2(\Z/N)=N^3\prod_{p\mid N}(1-p^{-2})$.
\item $[\SL_2(\Z):\Gamma_1(N)]=N^2\prod_{p\mid N}(1-p^{-2})$, and $\Gamma_1(N)$ is normal in $\Gamma_0(N)$ with $\Gamma_0(N)/\Gamma_1(N)\simeq(\Z/N)^\times$ via $\gamma\mapsto d\bmod N$.
\item The map $\Gamma_0(N)\gamma\mapsto(c:d)$, $(c,d)$ the bottom row of $\gamma$, is a bijection $\Gamma_0(N)\backslash\SL_2(\Z)\to\PP^1(\Z/N)$, the set of pairs $(c,d)\in(\Z/N)^2$ with
$\gcd(c,d,N)=1$ modulo multiplication by $(\Z/N)^\times$. Hence
\[
\mu(N)=[\SL_2(\Z):\Gamma_0(N)]=\#\PP^1(\Z/N)=N\prod_{p\mid N}\Bigl(1+\frac1p\Bigr).
\]
\item $-I\in\Gamma_0(N)$ for every $N$, and $-I\in\Gamma_1(N)$ if and only if $N\le2$.
\end{enumerate}
\end{proposition}

\begin{proof}
(i) Let $\bar\gamma=\sm{\bar a}{\bar b}{\bar c}{\bar d}\in\SL_2(\Z/N)$. Lift $\bar c,\bar d$ to integers $c,d$ with $c\ne0$ (take $c=N$ if $\bar c=0$). Since $\gcd(c,d,N)=1$, we can
replace $d$ by $d+tN$ with $\gcd(c,d+tN)=1$: let $t$ be the product of the primes dividing $c$ but not $d$; a prime $p\mid c$ that divides $d$ does not divide $N$
or $t$, so $p\nmid d+tN$, and a prime $p\mid c$ not dividing $d$ divides $t$, so again $p\nmid d+tN$. Now choose $a',b'\in\Z$ with $a'd-b'c=1$. The solutions
$(x,y)\in(\Z/N)^2$ of $xd-yc=1$ are $(a',b')+k(c,d)$, $k\in\Z/N$. Indeed the matrix $g=\sm{a'}{b'}cd$ has determinant $1$, so its rows $(a',b')$ and $(c,d)$ form a basis
of $(\Z/N)^2$, and every $(x,y)$ is $s(a',b')+k(c,d)$ with $s,k\in\Z/N$; then $xd-yc=s(a'd-b'c)+k(cd-dc)=s$, so $xd-yc=1$ exactly when $s=1$. So $(\bar a,\bar b)=(a'+kc,b'+kd)\bmod N$ for some $k\in\Z$, and $\sm{a'+kc}{b'+kd}cd\in\SL_2(\Z)$
lifts $\bar\gamma$. For the order: by the Chinese remainder theorem $\SL_2(\Z/N)=\prod\SL_2(\Z/p^e)$; reduction $\GL_2(\Z/p^e)\to\GL_2(\F_p)$ is onto with kernel
$I+pM_2(\Z/p^e)$ of order $p^{4(e-1)}$, $\#\GL_2(\F_p)=(p^2-1)(p^2-p)$, and $\det\colon\GL_2(\Z/p^e)\to(\Z/p^e)^\times$ is onto, so
$\#\SL_2(\Z/p^e)=p^{4(e-1)}(p^2-1)(p^2-p)/(p^{e-1}(p-1))=p^{3e}(1-p^{-2})$.
(ii) $\Gamma_1(N)$ is the preimage of the subgroup $\{\sm1b01\}$ of order $N$, so its index is $\#\SL_2(\Z/N)/N$. The map $\Gamma_0(N)\to(\Z/N)^\times$, $\gamma\mapsto d$, is a
homomorphism (the lower right entry of a product of upper triangular matrices is the product of the lower right entries), onto by (i) applied to
$\sm{d^{-1}}00d$, with kernel $\Gamma_1(N)$ (if $c\equiv0$ and $d\equiv1$ then $a\equiv1$).
(iii) $\Gamma_0(N)$ is the preimage of the upper triangular group $B\subseteq\SL_2(\Z/N)$, so $\Gamma_0(N)\backslash\SL_2(\Z)=B\backslash\SL_2(\Z/N)$ by (i). Left multiplication by
$\sm{u}{b}0{u^{-1}}$ multiplies the bottom row by $u^{-1}$; conversely if $g,g'\in\SL_2(\Z/N)$ have bottom rows $(c,d)$ and $u(c,d)$ then $g'g^{-1}$ has bottom row
$(0,u)$ and lies in $B$. Every row with $\gcd(c,d,N)=1$ is the bottom row of an element of $\SL_2(\Z/N)$ (complete it as in (i)). So
$\Gamma_0(N)\backslash\SL_2(\Z)\simeq\PP^1(\Z/N)$, and $\#\PP^1(\Z/N)=\prod_{p^e\|N}\#\PP^1(\Z/p^e)$ by the Chinese remainder theorem. To count $\PP^1(\Z/p^e)$: in a pair
$(c,d)$ with $\gcd(c,d,p)=1$ at least one entry is a unit. If $c$ is a unit, scaling by $c^{-1}$ gives a unique representative $(1:d)$, $d\in\Z/p^e$: $p^e$ points. If $c$ is not a
unit, then $d$ is, and scaling by $d^{-1}$ gives a unique representative $(c:1)$ with $p\mid c$: $p^{e-1}$ points. So $\#\PP^1(\Z/p^e)=p^e+p^{e-1}=p^e(1+\frac1p)$, and the product
over $p^e\,\|\,N$ is the stated formula. For example $\mu(20)=\#\PP^1(\Z/4)\cdot\#\PP^1(\Z/5)=6\cdot6=36$.
(iv) $-I$ is upper triangular, and $-I\equiv I\bmod N$ iff $N\mid2$.
\end{proof}

The values of $\mu(N)$ at the levels of the paper's examples are:
\begin{center}
\begin{tabular}{@{}lccccccccccc@{}}
\toprule
$N$ & $11$ & $20$ & $37$ & $52$ & $77$ & $99$ & $121$ & $140$ & $539$ & $726$ & $980$\\
$\mu(N)$ & $12$ & $36$ & $38$ & $84$ & $96$ & $144$ & $132$ & $288$ & $672$ & $1584$ & $2016$\\
\bottomrule
\end{tabular}
\end{center}
(computed from (iii); for $N\in\{11,20,37,77,99,121,140\}$ also by a direct count of $\PP^1(\Z/N)$).

\begin{proposition}[fundamental domains from cosets]\label{prop:modforms:coset-domains}
Let $\Gamma\subseteq\SL_2(\Z)$ be a subgroup of finite index containing $-I$, and $\SL_2(\Z)=\bigsqcup_{i=1}^m\Gamma\gamma_i$. Then $\mathcal F_\Gamma=\bigcup_i\gamma_i\mathcal F$ is a
fundamental domain for $\Gamma$: every point of $\HH$ is $\Gamma$-equivalent to a point of $\mathcal F_\Gamma$, and no two distinct interior points of $\mathcal F_\Gamma$ that lie in
the interiors of the translates $\gamma_i\mathcal F$ are $\Gamma$-equivalent. Its area is $\int_{\mathcal F_\Gamma}d\mu=m\pi/3$; for $\Gamma_0(N)$, $\operatorname{vol}(\Gamma_0(N)\backslash\HH)=\pi\mu(N)/3$.
For $N=p$ prime one may take $\gamma_0=I$ and $\gamma_{j+1}=ST^j$, $0\le j\le p-1$.
\end{proposition}

\begin{proof}
For $z\in\HH$ write $z=\gamma w$ with $w\in\mathcal F$ and $\gamma\in\SL_2(\Z)$ (\Cref{thm:modforms:fundamental-domain}(i)), and $\gamma=\delta\gamma_i$ with $\delta\in\Gamma$; then
$\delta^{-1}z=\gamma_iw\in\gamma_i\mathcal F$. If $\gamma_iw=\delta\gamma_jw'$ with $w,w'$ interior points of $\mathcal F$ and $\delta\in\Gamma$, then $w=\gamma_i^{-1}\delta\gamma_jw'$, so $w=w'$
and $\gamma_i^{-1}\delta\gamma_j=\pm I$ by \Cref{thm:modforms:fundamental-domain}(ii), (iii); as $-I\in\Gamma$, $\Gamma\gamma_i=\Gamma\gamma_j$, so $i=j$ and the two points coincide. The
translates meet only in boundaries, which have measure zero, so the area is $m$ times that of $\mathcal F$ by the invariance of $d\mu$. For $N=p$ the bottom
rows of $I$ and $ST^j=\sm0{-1}1j$ are $(0:1)$ and $(1:j)$, the $p+1$ points of $\PP^1(\F_p)$.
\end{proof}

For $\Gamma_0(11)$ the domain consists of $\mathcal F$ and the eleven translates $ST^j\mathcal F$, $0\le j\le10$. The translate $ST^j\mathcal F$ is the image of the strip piece
$T^j\mathcal F$ under $z\mapsto-1/z$, a curvilinear triangle bounded by three geodesics (two of them semicircles through $0$) with one ideal vertex at the cusp $0$. So eleven of the twelve
triangles cluster at $0$, and the domain has two ideal vertices, $\infty$ and $0$, which are the two cusps of $\Gamma_0(11)$ (\Cref{prop:modforms:cusps-gamma0}).

\section{Cusps, widths and elliptic points}\label{sec:modforms:cusps}

The compact curve $X_0(N)$ is $Y_0(N)$ together with finitely many cusps, and the map $X_0(N)\to X(1)$ is ramified only at the cusps and at the
elliptic points. This section counts both. The counts feed the genus formula (\Cref{thm:modforms:genus-formula}), the dimension formulas
(\Cref{thm:modforms:dimension-formulas}) and the action of the Atkin-Lehner involutions on cusps (\Cref{prop:modforms:atkin-lehner-cusps}). In the paper
the cusp $\infty$ is the base point of the modular parametrisation $\varphi_E\colon X_0(N)\to E$ (\Cref{ch:moddeg}), and the expansion of $f$ at $\infty$ is the
$q$-expansion whose coefficients are the $a_n(f)$ of every congruence.

\begin{definition}[cusps and widths]\label{def:modforms:cusp-width}
Let $\Gamma$ be a congruence subgroup. The \emph{cusps} of $\Gamma$ are the $\Gamma$-orbits on $\PP^1(\Q)$. For a cusp $s=\sigma\infty$ ($\sigma\in\SL_2(\Z)$) the
\emph{width} of $s$ is the least $h\ge1$ such that $\sigma T^h\sigma^{-1}\in\pm\Gamma$. If $-I\notin\Gamma$, the cusp is \emph{irregular} when $-\sigma T^h\sigma^{-1}\in\Gamma$ but
$\sigma T^h\sigma^{-1}\notin\Gamma$, and \emph{regular} otherwise.
\end{definition}

The width exists because $\sigma T^N\sigma^{-1}\in\Gamma(N)\subseteq\Gamma$ when $\Gamma\supseteq\Gamma(N)$ ($\Gamma(N)$ is normal). It does not depend on the choices: $\sigma'\infty=\sigma\infty$ iff
$\sigma'\in\sigma\cdot\{\pm T^m\}$, which commutes with $T^h$, and replacing $s$ by $\gamma s$ ($\gamma\in\Gamma$) conjugates $\sigma T^h\sigma^{-1}$ by $\gamma$.

\begin{proposition}[cusps as orbits]\label{prop:modforms:cusps-orbits}
Let $\Gamma$ be a congruence subgroup containing $-I$. The map $\Gamma\sigma\mapsto\Gamma\sigma\infty$ induces a bijection between the orbits of $\langle T\rangle$ acting on the right on
$\Gamma\backslash\SL_2(\Z)$ and the cusps of $\Gamma$, and the width of the cusp $\Gamma\sigma\infty$ is the length of the orbit of $\Gamma\sigma$. In particular $\Gamma$ has finitely many cusps
and $\sum_s h_s=[\SL_2(\Z):\Gamma]$. For $\Gamma_0(N)$, via \Cref{prop:modforms:indices}(iii), the cusps are the orbits of $(c:d)\mapsto(c:c+d)$ on $\PP^1(\Z/N)$.
\end{proposition}

\begin{proof}
The map $\Gamma\backslash\SL_2(\Z)\to\Gamma\backslash\PP^1(\Q)$ is onto because $\SL_2(\Z)$ is transitive on $\PP^1(\Q)$ (\Cref{prop:modforms:element-types}(iii)). Its fibres:
$\Gamma\sigma\infty=\Gamma\sigma'\infty$ iff $\sigma'\infty=\gamma\sigma\infty$ for some $\gamma\in\Gamma$ iff $\sigma'\in\Gamma\sigma\operatorname{Stab}(\infty)=\Gamma\sigma\{\pm T^m\}$, that is, iff
$\Gamma\sigma'$ lies in the $\langle T\rangle$-orbit of $\Gamma\sigma$ ($-I\in\Gamma$). The orbit of $\Gamma\sigma$ has length the least $h\ge1$ with $\Gamma\sigma T^h=\Gamma\sigma$, that is,
$\sigma T^h\sigma^{-1}\in\Gamma$. The right action of $T$ on bottom rows is $(c,d)\mapsto(c,c+d)$.
\end{proof}

The simplest case shows how the bookkeeping works. For $\Gamma_0(p)$, $p$ prime, $\PP^1(\F_p)$ has the $p+1$ points $(0:1)$ and $(1:d)$, $d\in\F_p$. The map $(c:d)\mapsto(c:c+d)$
fixes $(0:1)$, which is the coset of $I$, so the cusp $\infty=I\infty$ has width $1$; and it sends $(1:d)$ to $(1:d+1)$, so the $p$ points $(1:d)$ form a single orbit of length $p$,
containing $(1:0)$, the coset of $S$, whose cusp is $S\infty=0$. Hence $\Gamma_0(p)$ has the two cusps $\infty$ and $0$, of widths $1$ and $p$, and $1+p=\mu(p)$. Geometrically: of the
$p+1$ translates of $\mathcal F$ making up the fundamental domain of \Cref{prop:modforms:coset-domains}, one touches $\infty$ and $p$ touch $0$.

\begin{warning}[irregular cusps]\label{warn:modforms:irregular}
If $-I\notin\Gamma$, \Cref{prop:modforms:cusps-orbits} holds with $\PSL_2$ in place of $\SL_2$, and some cusps can be irregular: for $\Gamma_1(4)$ and $\sigma=\sm1021$
($s=\frac12$), $\sigma T\sigma^{-1}=\sm{-1}1{-4}3$ is congruent to $\sm{-1}10{-1}$ modulo $4$, so $-\sigma T\sigma^{-1}\in\Gamma_1(4)$ but $\sigma T\sigma^{-1}\notin\Gamma_1(4)$. At an irregular
cusp a form of odd weight has an expansion in half-integral powers of $q_h$, and ``$b_0=0$'' is not the right vanishing condition. None of this happens
for $\Gamma_0(N)$, which contains $-I$; the paper only needs $\Gamma_0(N)$.
\end{warning}

\begin{proposition}[the cusps of $\Gamma_0(N)$]\label{prop:modforms:cusps-gamma0}
Let $N\ge1$. Write a point of $\PP^1(\Q)$ as $x/y$ with $\gcd(x,y)=1$, $y\ge0$ (and $\infty=1/0$).
\begin{enumerate}[label=(\roman*)]
\item The integer $\gcd(y,N)$ depends only on the $\Gamma_0(N)$-class of $x/y$. The cusp $\infty\sim1/N$ has width $1$ and the cusp $0=0/1$ has width $N$.
\item Every cusp is equivalent to some $a/c$ with $c\mid N$, $c\ge1$, $\gcd(a,c)=1$, and $a/c$, $a'/c'$ of this form are equivalent iff $c=c'$ and
$a\equiv a'\pmod{\gcd(c,N/c)}$. Hence
\[
\nu_\infty(N)=\sum_{c\mid N}\varphi\bigl(\gcd(c,N/c)\bigr),
\]
and the cusp $a/c$ has width $N/\gcd(c^2,N)$.
\item If $N$ is squarefree, the cusps are the $2^{\omega(N)}$ classes of $1/c$, $c\mid N$, and $1/c$ has width $N/c$.
\item If $N=p^2$ with $p$ prime, the cusps are $\infty$, $0$ and $u/p$ for $1\le u\le p-1$, of widths $1$, $p^2$ and $1$; so $\nu_\infty(p^2)=p+1$.
\end{enumerate}
\end{proposition}

\begin{proof}
(i) For $\gamma=\sm abcd\in\Gamma_0(N)$, $\gamma(x/y)=(ax+by)/(cx+dy)$, and $(ax+by,cx+dy)$ is primitive because $\gamma\in\SL_2(\Z)$; since $N\mid c$ and $\gcd(d,N)=1$,
$\gcd(cx+dy,N)=\gcd(dy,N)=\gcd(y,N)$. For the widths: $T\in\Gamma_0(N)$, and $ST^hS^{-1}=\sm10{-h}1\in\Gamma_0(N)$ iff $N\mid h$. Finally $1/N=\sm10N1\infty$ with $\sm10N1\in\Gamma_0(N)$.

(iii) By the Chinese remainder theorem $\PP^1(\Z/N)=\prod_{p\mid N}\PP^1(\F_p)$, and $T$ acts on each factor by $(c:d)\mapsto(c:c+d)$, which fixes $(0:1)$ and permutes the $p$
points $(1:d)$ cyclically. For a point $P$ let $A(P)$ be the set of primes $p$ at which its component is not $(0:1)$. Then $T^jP=P$ iff $p\mid j$ for all $p\in A(P)$, and
by the Chinese remainder theorem $j\bmod\prod_{p\in A}p$ moves the components at $A$ simultaneously through all of $\prod_{p\in A}\F_p$. So the $\langle T\rangle$-orbits are the sets
$\{P:A(P)=A\}$, one for each subset $A$, of length $\prod_{p\in A}p$. For $c=\prod_{p\notin A}p$ the coset of $\sigma=\sm10c1$ has bottom row $(c:1)$, which has $A(P)=A$, and
$\sigma\infty=1/c$; its width is $\prod_{p\in A}p=N/c$ by \Cref{prop:modforms:cusps-orbits}.

(iv) $\PP^1(\Z/p^2)$ consists of the $p^2$ points $(1:d)$ and the $p$ points $(pu:1)$, $u\in\Z/p$. The points $(1:d)$ form one $T$-orbit, which contains $(1:0)$, the coset
of $S$; its cusp is $S\infty=0$, of width $p^2$. Each $(pu:1)$ is fixed by $T$, since $(pu:1+pu)=((1+pu)^{-1}pu:1)=(pu:1)$. The point $(0:1)$ is the coset of $I$
(cusp $\infty$). For $1\le u\le p-1$ choose $s,t$ with $us-tp=1$ and $\sigma=\sm utps\in\SL_2(\Z)$; its bottom row $(p:s)$ equals $(ps^{-1}:1)=(pu:1)$ because $s^{-1}\equiv u
\pmod p$; so the cusp of $(pu:1)$ is $\sigma\infty=u/p$, of width $1$.

(ii) We do not prove the general case; see \cite[Section 3.8]{DiamondShurman} (where the statement is a proposition; number not verified) and
\cite{Ogg73}. Sketch: by (i) the class of $x/y$ determines $c=\gcd(y,N)$; one moves $x/y$ by an element of $\Gamma_0(N)$ to a fraction $a/c$ with the same
$c$, and then shows that the residue of $a$ modulo $\gcd(c,N/c)$ is the remaining invariant, by the same orbit computation on $\PP^1(\Z/N)$ as in (iii) and (iv),
prime by prime. The counts and widths in (ii) agree with (iii), (iv), and with a direct computation of the $T$-orbits on $\PP^1(\Z/N)$ for $N\in\{8,12,18,20,36,72,99,121\}$.
\end{proof}

\begin{example}[cusps at the paper's levels]\label{ex:modforms:cusp-lists}
The following lists come from (ii), and each was checked against the $T$-orbits on $\PP^1(\Z/N)$:
\begin{itemize}
\item $N=11$: $\infty$, $0$, of widths $1$, $11$ (sum $12=\mu(11)$).
\item $N=20$: $\infty$, $\frac1{10}$, $\frac15$, $\frac14$, $\frac12$, $0$, of widths $1,1,4,5,5,20$ (sum $36$).
\item $N=99$: the cusps $a/c$ for $c\in\{1,3,9,11,33,99\}$, with $a$ taken modulo $\gcd(c,99/c)$, which is $1,3,1,1,3,1$ respectively. There are eight cusps, of widths $99$ ($c=1$), $11,11$ ($c=3$: the cusps
$\frac13,\frac23$), $11$ ($c=9$), $9$ ($c=11$), $1,1$ ($c=33$), $1$ ($c=99$); sum $144$.
\item $N=121$: twelve cusps, $\infty$, $0$ and $u/11$ for $1\le u\le10$, of widths $1$, $121$ and ten of width $1$ (sum $132$).
\item $\nu_\infty(539)=16$, $\nu_\infty(726)=\nu_\infty(980)=48$.
\end{itemize}
\end{example}

\begin{proposition}[elliptic points of $\Gamma_0(N)$]\label{prop:modforms:elliptic-gamma0}
A point $z\in\HH$ is an \emph{elliptic point} of $\Gamma_0(N)$ if its stabiliser in $\bar\Gamma_0(N)$ is nontrivial; its order is the order of that stabiliser. The
elliptic points of order $2$ of $Y_0(N)=\Gamma_0(N)\backslash\HH$ lie over the orbit of $i$ and are in bijection with the solutions $u\bmod N$ of $u^2+1\equiv0$; those of order $3$
lie over the orbit of $\zeta_3$ and are in bijection with the solutions of $u^2-u+1\equiv0\pmod N$. Hence
\[
\nu_2(N)=\begin{cases}0,&4\mid N,\\ \prod_{p\mid N}\bigl(1+\bigl(\frac{-4}p\bigr)\bigr),&\text{otherwise},\end{cases}\qquad
\nu_3(N)=\begin{cases}0,&9\mid N,\\ \prod_{p\mid N}\bigl(1+\bigl(\frac{-3}p\bigr)\bigr),&\text{otherwise},\end{cases}
\]
with Kronecker symbols (\Cref{ch:rings}), so $\bigl(\frac{-4}2\bigr)=\bigl(\frac{-3}3\bigr)=0$. Moreover the fibre of $Y_0(N)\to Y(1)$ over the orbit of $i$ consists of the $\nu_2(N)$
elliptic points and $(\mu(N)-\nu_2(N))/2$ points with trivial stabiliser, and over the orbit of $\zeta_3$ of the $\nu_3(N)$ elliptic points and $(\mu(N)-\nu_3(N))/3$
points with trivial stabiliser.
\end{proposition}

\begin{proof}
By \Cref{cor:modforms:level-one-quotient} every elliptic point is $\Gamma_0(N)$-equivalent to $\gamma i$ or $\gamma\zeta_3$ for some $\gamma\in\SL_2(\Z)$, with stabiliser
$\gamma\langle S\rangle\gamma^{-1}$, respectively $\gamma\langle ST\rangle\gamma^{-1}$, in $\PSL_2(\Z)$. The points of $Y_0(N)$ over the orbit of $i$ are the classes $\Gamma_0(N)\gamma i$, and
$\Gamma_0(N)\gamma i=\Gamma_0(N)\gamma'i$ iff $\gamma'\in\Gamma_0(N)\gamma\langle S\rangle$. So they correspond to the orbits of $\langle S\rangle$ (of order $2$ in $\PSL_2$) acting on the right on
$\Gamma_0(N)\backslash\SL_2(\Z)=\PP^1(\Z/N)$. The point $\gamma i$ is elliptic for $\Gamma_0(N)$ iff $\gamma S\gamma^{-1}\in\Gamma_0(N)$ iff $\Gamma_0(N)\gamma S=\Gamma_0(N)\gamma$: the elliptic points
are the fixed points of $S$, and the other points form orbits of length $2$; this gives the count of the fibre. Now $(c,d)S=(d,-c)$, and $(d:-c)=(c:d)$ means
$d=uc$, $-c=ud$ for a unit $u$; then $\gcd(c,N)=\gcd(c,d,N)=1$, so $c$ is a unit and we may take $c=1$: the fixed points are the $(1:u)$ with $u^2+1\equiv0$,
distinct for distinct $u$. (Explicitly: $d=uc$ and $-c=ud=u^2c$ give $c(u^2+1)=0$ with $c$ a unit, so $u^2+1=0$; conversely $(1,u)S=(u,-1)=u\cdot(1,u)$ when $u^2=-1$.)
Likewise $\langle ST\rangle$ has order $3$ in $\PSL_2(\Z)$, $ST=\sm0{-1}11$ and $(c,d)ST=(d,d-c)$. The equation $(d:d-c)=(c:d)$ means $d=uc$ and $d-c=ud$ for a unit $u$, so
$uc-c=u^2c$, that is, $c(u^2-u+1)=0$; as before $c$ is a unit, we may take $c=1$, and the fixed points are the $(1:u)$ with $u^2-u+1\equiv0$ (conversely
$(1,u)ST=(u,u-1)=(u,u^2)=u\cdot(1,u)$ for such $u$). The points that are not fixed fall into orbits of length $3$, which gives the count of the fibre over $\zeta_3$. For instance, modulo $13$
the roots of $u^2+1$ are $5,8$ and those of $u^2-u+1$ are $4,10$, so the elliptic points of $\Gamma_0(13)$ are the classes of $(1:5)$, $(1:8)$ (order $2$) and $(1:4)$, $(1:10)$
(order $3$), and $\nu_2(13)=\nu_3(13)=2$.

The number of roots is multiplicative in $N$ (Chinese remainder theorem). For $p$ odd, $u^2+1$ has $1+(\frac{-1}p)=1+(\frac{-4}p)$ roots modulo $p$, simple, which lift
uniquely modulo $p^e$ by Hensel's lemma (\Cref{ch:rings}); modulo $2$ there is one root, and modulo $4$ none. For $p\ne3$, $u^2-u+1$ has discriminant $-3$ and
$1+(\frac{-3}p)$ roots modulo $p$, simple if $p\ne3$, lifting uniquely (for $p=2$: none, $(\frac{-3}2)=-1$); modulo $3$ it is $(u+1)^2$, with one root, and modulo
$9$ it has none ($u\equiv2,5,8$ give $3,21,57\equiv3\pmod9$).
\end{proof}

For example $\nu_2(11)=\nu_3(11)=0$; $\nu_2(13)=\nu_3(13)=2$; $\nu_2(37)=\nu_3(37)=2$; $\nu_2(5)=\nu_2(10)=2$, $\nu_3(5)=\nu_3(10)=0$; and $\nu_2=\nu_3=0$ at every level
$N\in\{20,52,77,99,121,140,539,726,980\}$ of the paper's examples (for $\nu_2$: $4\mid N$ or a prime $p\equiv3\pmod4$ divides $N$; for $\nu_3$: $9\mid N$ or a prime
$p\equiv2\pmod3$ divides $N$).

\section{The modular curves as compact Riemann surfaces}\label{sec:modforms:modular-curves}

This section builds the curve of the paper's title page: $X_0(N)$, the compact Riemann surface on which the newform $f$ is a differential and from which
the modular parametrisation $\varphi_E\colon X_0(N)\to E$ starts. We construct $Y_0(N)$ and $X_0(N)$ with their charts, prove the genus formula from the
Riemann-Hurwitz formula of \Cref{ch:curves}, describe $Y_0(N)$ as the moduli space of elliptic curves with a cyclic subgroup of order $N$, and state the
existence of the model over $\Q$. In Section 4.1 of the paper the cohomology $H^1(X_0(M^\Sigma),\Z_\ell)$ is the cohomology of this compact Riemann surface.

\begin{theorem}[$Y(\Gamma)$ is a Riemann surface]\label{thm:modforms:y0-riemann-surface}
Let $\Gamma$ be a congruence subgroup.
\begin{enumerate}[label=(\roman*)]
\item $\Gamma$ acts properly discontinuously on $\HH$: for every compact $K\subset\HH$ the set $\{\gamma\in\Gamma:\gamma K\cap K\ne\emptyset\}$ is finite. Every $z\in\HH$ has a
neighbourhood $U$ such that $\gamma U\cap U\ne\emptyset$ only for $\gamma$ in the stabiliser $\Gamma_z$, which is finite and cyclic modulo $\pm I$. The quotient $Y(\Gamma)=\Gamma\backslash\HH$ is a
connected Hausdorff space.
\item $Y(\Gamma)$ is a Riemann surface: near the image of a point $z_0$ with $\#\bar\Gamma_{z_0}=e$, the function $t=\bigl(\frac{z-z_0}{z-\bar z_0}\bigr)^e$ is a chart, and the
transition maps are holomorphic. A $\Gamma$-invariant holomorphic function on an open $\Gamma$-stable subset of $\HH$ is a holomorphic function on its image in
$Y(\Gamma)$.
\end{enumerate}
\end{theorem}

\begin{proof}
(i) Choose $0<y_1<y_2$ and $R$ with $K\subseteq\{y_1\le\operatorname{Im}z\le y_2,\ \abs z\le R\}$. If $z\in K$ and $\gamma z\in K$ then
$\abs{cz+d}^2=\operatorname{Im}z/\operatorname{Im}\gamma z\le y_2/y_1$, so $\abs c\le\sqrt{y_2/y_1}/y_1$ and $\abs d\le\sqrt{y_2/y_1}+\abs cR$. Applying this to $\gamma^{-1}=\sm d{-b}{-c}a$
and the point $\gamma z\in K$ bounds $\abs a$. If $c\ne0$ then $b=(ad-1)/c$; if $c=0$ then $\gamma=\pm\sm1b01$ with $\abs b\le2R$. So finitely many $\gamma$. The
stabiliser $\Gamma_z$ is a discrete subgroup of the compact group $g\,\mathrm{SO}(2)g^{-1}$ (\Cref{prop:modforms:upper-half-plane}(iv)), hence finite, and modulo $\pm I$ it is a
finite subgroup of the circle group, hence cyclic. Given $z$, take a closed disc $K$ around $z$; of the finitely many $\gamma$ with $\gamma K\cap K\ne\emptyset$, those not
fixing $z$ move $z$ off itself, and a small enough open disc $U\ni z$ satisfies $\gamma U\cap U=\emptyset$ for them. Likewise, if $z_1,z_2$ are not $\Gamma$-equivalent,
take a closed disc $K$ containing both; for the finitely many $\gamma$ with $\gamma K\cap K\ne\emptyset$ we have $\gamma z_1\ne z_2$, so small discs $U_1\ni z_1$, $U_2\ni z_2$
satisfy $\gamma U_1\cap U_2=\emptyset$ for all $\gamma\in\Gamma$; their images in $Y(\Gamma)$ are disjoint open sets. $Y(\Gamma)$ is connected as the image of $\HH$.
(ii) We sketch this; see \cite[Sections 2.1, 2.2]{DiamondShurman} and \cite[Section 1.5]{Shimura71}. The Cayley map $w=\frac{z-z_0}{z-\bar z_0}$ sends $z_0$ to $0$ and
conjugates $\bar\Gamma_{z_0}$ to the rotations $w\mapsto\zeta w$, $\zeta^e=1$, so $t=w^e$ identifies a small disc around $z_0$ modulo $\bar\Gamma_{z_0}$ with a disc around $0$.
Where two such charts overlap, at least one of the points involved has trivial stabiliser and the transition is a composite of M\"obius maps and an
$e$-th power or root, holomorphic away from $t=0$. A $\Gamma_{z_0}$-invariant holomorphic function of $w$ is a holomorphic function of $w^e$ (expand in powers of $w$:
only exponents divisible by $e$ survive).
\end{proof}

\begin{theorem}[$X(\Gamma)$ is a compact Riemann surface]\label{thm:modforms:x0-compact}
Give $\HH^*=\HH\cup\PP^1(\Q)$ the topology in which the neighbourhoods of $\infty$ are the sets $\{\operatorname{Im}z>C\}\cup\{\infty\}$ and those of $s=\sigma\infty$ are their images
under $\sigma\in\SL_2(\Z)$ (discs tangent to $\R$ at $s$, together with $s$). For a congruence subgroup $\Gamma$, $X(\Gamma)=\Gamma\backslash\HH^*$ is a compact connected Hausdorff
space, and with the chart $q_h=e^{2\pi i\sigma^{-1}(z)/h}$ near a cusp $s=\sigma\infty$ of width $h$ (\Cref{def:modforms:cusp-width}) it is a compact Riemann surface. We write
$X_0(N)$, $X_1(N)$, $X(N)$ for $\Gamma_0(N)$, $\Gamma_1(N)$, $\Gamma(N)$, and $X(1)$ for $\SL_2(\Z)$. By the Riemann existence theorem and GAGA (\Cref{ch:curves}), $X_0(N)$ is
the Riemann surface of a smooth projective curve over $\C$, and holomorphic maps between these curves are morphisms.
\end{theorem}

\begin{proof}[Sketch of proof]
Compactness: $\HH^*=\SL_2(\Z)\cdot(\mathcal F\cup\{\infty\})$ by \Cref{thm:modforms:fundamental-domain}(i), so $X(\Gamma)$ is the image of the finitely many compact sets
$\gamma_i(\mathcal F\cup\{\infty\})$ of \Cref{prop:modforms:coset-domains}. Hausdorff: the horoball neighbourhoods of inequivalent cusps, and of a cusp and a point of $\HH$,
can be chosen with disjoint $\Gamma$-orbits. Chart at a cusp: the stabiliser of $\infty$ in $\sigma^{-1}\Gamma\sigma$ is generated by $\pm T^h$, and $z\mapsto e^{2\pi iz/h}$ identifies
$\{\operatorname{Im}z>C\}$ modulo $T^h$ with a punctured disc; adding the cusp fills the puncture. See \cite[Section 2.4]{DiamondShurman}, \cite[Section 1.5]{Shimura71}. For
Riemann existence and GAGA see \Cref{ch:curves}.
\end{proof}

\begin{theorem}[the genus of $X_0(N)$]\label{thm:modforms:genus-formula}
The map $\pi\colon X_0(N)\to X(1)$, $\Gamma_0(N)z\mapsto\SL_2(\Z)z$, is a holomorphic map of degree $\mu(N)$, unramified outside the images of $i$, $\zeta_3$ and $\infty$. Over $i$ it
has $\nu_2(N)$ points with $e=1$ and $(\mu(N)-\nu_2(N))/2$ points with $e=2$; over $\zeta_3$, $\nu_3(N)$ points with $e=1$ and $(\mu(N)-\nu_3(N))/3$ with $e=3$; over $\infty$ the
$\nu_\infty(N)$ cusps, with $e$ equal to the width. Consequently
\[
g(X_0(N))=1+\frac{\mu(N)}{12}-\frac{\nu_2(N)}4-\frac{\nu_3(N)}3-\frac{\nu_\infty(N)}2.
\]
\end{theorem}

\begin{proof}
$\pi$ is continuous and, in the charts of \Cref{thm:modforms:y0-riemann-surface,thm:modforms:x0-compact}, holomorphic. Ramification: let $z\in\HH$ have stabilisers of orders $e_1$ in
$\bar\Gamma_0(N)$ and $e_2$ in $\PSL_2(\Z)$, so $e_1\mid e_2$. With $w$ the Cayley coordinate at $z$, the chart of $X_0(N)$ is $t=w^{e_1}$ and that of $X(1)$ is $s=w^{e_2}$, so
$s=t^{e_2/e_1}$ and the ramification index is $e_2/e_1$. At a cusp $s=\sigma\infty$ of width $h$, the chart of $X_0(N)$ is $q_h$ and that of $X(1)$ at $\infty$ is
$q=e^{2\pi i\sigma^{-1}(z)}=q_h^h$, so the ramification index is $h$. By \Cref{cor:modforms:level-one-quotient} only the orbits of $i$ and $\zeta_3$ have $e_2>1$, so $\pi$ is
unramified over the rest of $Y(1)$, and the degree is the number of points over a generic point, which is $[\PSL_2(\Z):\bar\Gamma_0(N)]=\mu(N)$ (cosets of
$\Gamma_0(N)$ act simply on the orbit of a point with trivial stabiliser). The fibres over $i$ and $\zeta_3$ were counted in \Cref{prop:modforms:elliptic-gamma0}, and
the widths add up to $\mu(N)$ (\Cref{prop:modforms:cusps-orbits}). Now $g(X(1))=0$ by \Cref{cor:modforms:level-one-quotient}, and the Riemann-Hurwitz formula of
\Cref{ch:curves}, $2g-2=\deg\pi\cdot(2g(X(1))-2)+\sum_P(e_P-1)$, gives the following. Over $i$ the $(\mu-\nu_2)/2$ points with $e=2$ contribute $1$ each; over $\zeta_3$ the
$(\mu-\nu_3)/3$ points with $e=3$ contribute $2$ each; over $\infty$ the cusps contribute $\sum_s(h_s-1)=\mu-\nu_\infty$; all other points have $e=1$. So
\[
2g-2=\mu\cdot(-2)+\frac{\mu-\nu_2}2+\frac{2(\mu-\nu_3)}3+(\mu-\nu_\infty)=\frac\mu6-\frac{\nu_2}2-\frac{2\nu_3}3-\nu_\infty,
\]
since $-2+\frac12+\frac23+1=\frac16$; dividing by $2$ and adding $1$ gives the formula.
\end{proof}

Two instances. For $N=11$: $g=1+\frac{12}{12}-0-0-\frac22=1$. For $N=37$, where $u^2+1$ and $u^2-u+1$ both have two roots modulo $37$ ($37\equiv1\pmod{12}$):
$g=1+\frac{38}{12}-\frac24-\frac23-\frac22=\frac{12+38-6-8-12}{12}=2$.

\begin{center}
\begin{tabular}{@{}rrrrrr@{\qquad}rrrrrr@{}}
\toprule
$N$ & $\mu$ & $\nu_2$ & $\nu_3$ & $\nu_\infty$ & $g$ & $N$ & $\mu$ & $\nu_2$ & $\nu_3$ & $\nu_\infty$ & $g$\\
\midrule
$1$ & $1$ & $1$ & $1$ & $1$ & $0$ & $37$ & $38$ & $2$ & $2$ & $2$ & $2$\\
$11$ & $12$ & $0$ & $0$ & $2$ & $1$ & $43$ & $44$ & $0$ & $2$ & $2$ & $3$\\
$13$ & $14$ & $2$ & $2$ & $2$ & $0$ & $77$ & $96$ & $0$ & $0$ & $4$ & $7$\\
$14$ & $24$ & $0$ & $0$ & $4$ & $1$ & $99$ & $144$ & $0$ & $0$ & $8$ & $9$\\
$20$ & $36$ & $0$ & $0$ & $6$ & $1$ & $121$ & $132$ & $0$ & $0$ & $12$ & $6$\\
$23$ & $24$ & $0$ & $0$ & $2$ & $2$ & $140$ & $288$ & $0$ & $0$ & $12$ & $19$\\
$25$ & $30$ & $2$ & $0$ & $6$ & $0$ & $539$ & $672$ & $0$ & $0$ & $16$ & $49$\\
$27$ & $36$ & $0$ & $0$ & $6$ & $1$ & $726$ & $1584$ & $0$ & $0$ & $48$ & $109$\\
$36$ & $72$ & $0$ & $0$ & $12$ & $1$ & $980$ & $2016$ & $0$ & $0$ & $48$ & $145$\\
\bottomrule
\end{tabular}
\end{center}
Every entry was computed from the formulas of \Cref{prop:modforms:indices,prop:modforms:cusps-gamma0,prop:modforms:elliptic-gamma0}; for $N\le140$ the counts
$\mu,\nu_2,\nu_3,\nu_\infty$ were also obtained by brute force on $\PP^1(\Z/N)$ (fixed points of $S$ and $ST$, orbits of $T$), and for every tabulated $N$ the genus agrees
with PARI's \texttt{mfdim([N,2],1)}, which is $\dim S_2(\Gamma_0(N))=g$ by \Cref{thm:modforms:weight-two-differentials}. By the formula, among $N\le2000$ the genus is $0$
exactly for $N\in\{1,\dots,10,12,13,16,18,25\}$ and $1$ exactly for $N\in\{11,14,15,17,19,20,21,24,27,32,36,49\}$; Exercise \ref{ex:modforms:ex-genus-bound} shows that
$g(X_0(N))\ge2$ for all larger $N$, so these lists are complete. In particular $X_0(11)$ has genus $1$; it is the elliptic curve \curve{11a1} (\Cref{ch:moddeg}). These
values agree with those quoted in \Cref{ch:curves}.

\begin{theorem}[the moduli interpretation]\label{thm:modforms:moduli}
For $\tau\in\HH$ let $\Lambda_\tau=\Z+\Z\tau$ and $E_\tau=\C/\Lambda_\tau$.
\begin{enumerate}[label=(\roman*)]
\item Every pair $(E,C)$ of an elliptic curve $E/\C$ and a cyclic subgroup $C\subset E$ of order $N$ is isomorphic to $(E_\tau,\langle\frac1N+\Lambda_\tau\rangle)$ for some $\tau\in\HH$, and
$\Gamma_0(N)\tau\mapsto(E_\tau,\langle\frac1N\rangle)$ is a bijection from $Y_0(N)$ onto the set of isomorphism classes of such pairs.
\item Likewise $\Gamma_1(N)\tau\mapsto(E_\tau,\frac1N+\Lambda_\tau)$ is a bijection from $Y_1(N)$ onto the isomorphism classes of pairs $(E,P)$ with $P\in E$ of exact order $N$. Let
$\zeta_N=e_N(\frac1N+\Lambda_\tau,\frac\tau N+\Lambda_\tau)$, with $e_N$ the Weil pairing of \Cref{ch:ellcurves}; it is the same primitive $N$-th root of unity for every $\tau\in\HH$, namely
$e^{2\pi i/N}$ or $e^{-2\pi i/N}$ according to the normalisation of $e_N$ (\Cref{ch:ellcurves}). Then $\Gamma(N)\tau\mapsto(E_\tau,(\frac1N+\Lambda_\tau,\frac\tau N+\Lambda_\tau))$ is a bijection from
$Y(N)$ onto the isomorphism classes of pairs $(E,(P,Q))$ with $(P,Q)$ a basis of $E[N]$ and $e_N(P,Q)=\zeta_N$.
\item Under (i) the map $\pi$ of \Cref{thm:modforms:genus-formula} is $(E,C)\mapsto E$, and $(E,C)\mapsto(E/C,E[N]/C)$ is $\Gamma_0(N)\tau\mapsto\Gamma_0(N)\bigl(-1/(N\tau)\bigr)$, the
Atkin-Lehner involution $w_N$ (\Cref{prop:modforms:atkin-lehner}).
\item (Cusps.) The cusps correspond to degenerate objects: generalised elliptic curves (N\'eron polygons) with a cyclic subgroup of order $N$; over the Tate
curve $\mathbb G_m/\mathsf q^\Z$ the cusp $\infty$ corresponds to the subgroup $\mu_N$ and the cusp $0$ to a subgroup generated by an $N$-th root of $\mathsf q$.
\item ($\Q$-structure.) There is a smooth projective curve $X_0(N)_\Q$ over $\Q$ whose complex points are $X_0(N)$, such that the bijection (i) is compatible with the
action of $\Gal(\Qbar/\Q)$ on pairs $(E,C)$ defined over $\Qbar$. The cusps $\infty$ and $0$ are rational, and a cusp $a/c$, $c\mid N$, is defined over $\Q(\zeta_{\gcd(c,N/c)})$.
\end{enumerate}
\end{theorem}

\begin{proof}
(i) By \Cref{ch:ellcurves} every $E/\C$ is isomorphic to some $E_\tau$, every isomorphism $E_\tau\to E_{\tau'}$ is $z\mapsto\alpha z$ with $\alpha\Lambda_\tau=\Lambda_{\tau'}$, and for every such $\alpha$
there is $\gamma=\sm abcd\in\SL_2(\Z)$ with $\tau'=\gamma\tau$ and $\alpha=\pm(c\tau+d)^{-1}$. A cyclic subgroup of order $N$ of $\frac1N\Lambda_\tau/\Lambda_\tau\simeq(\Z/N)^2$ is generated by
$(c\tau+d)/N$ with $\gcd(c,d,N)=1$; by \Cref{prop:modforms:indices}(i) we may take $\sm abcd\in\SL_2(\Z)$, and then $z\mapsto z/(c\tau+d)$ maps $(E_\tau,\langle\frac{c\tau+d}N\rangle)$
to $(E_{\gamma\tau},\langle\frac1N\rangle)$. So every pair has the stated form. For injectivity, let $\tau'=\gamma\tau$ as above. From $\tau=\gamma^{-1}\tau'=\frac{d\tau'-b}{-c\tau'+a}$,
\[
c\tau+d=\frac{c(d\tau'-b)+d(a-c\tau')}{a-c\tau'}=\frac{ad-bc}{a-c\tau'}=(a-c\tau')^{-1},
\]
so the isomorphism $z\mapsto z/(c\tau+d)$ sends $\frac1N$ to $\frac{a-c\tau'}N$, which generates $\langle\frac1N+\Lambda_{\tau'}\rangle$ iff $N\mid c$ (and then $a$ is a
unit modulo $N$). So $(E_\tau,\langle\frac1N\rangle)\simeq(E_{\tau'},\langle\frac1N\rangle)$ iff $\tau'\in\Gamma_0(N)\tau$; the sign $\pm$ does not matter since $-I\in\Gamma_0(N)$.
(ii) The same computation: $\frac{a-c\tau'}N\equiv\pm\frac1N\pmod{\Lambda_{\tau'}}$ iff $N\mid c$ and $a\equiv\pm1\pmod N$, that is, $\gamma\in\pm\Gamma_1(N)$, which acts on $\HH$ as $\Gamma_1(N)$.
For $Y(N)$: if $(P,Q)$ is a basis of $E_\tau[N]$, write $P=\frac{a+b\tau}N$, $Q=\frac{c+d\tau}N$ modulo $\Lambda_\tau$; bilinearity gives $e_N(P,Q)=\zeta_N^{ad-bc}$, so $e_N(P,Q)=\zeta_N$ iff
$ad-bc\equiv1\pmod N$, and by \Cref{prop:modforms:indices}(i) we may then take $\sm abcd\in\SL_2(\Z)$. Then $\omega_1=a+b\tau$, $\omega_2=c+d\tau$ is a basis of $\Lambda_\tau$ with
$\tau'=\omega_2/\omega_1=\frac{d\tau+c}{b\tau+a}\in\HH$, and $z\mapsto z/\omega_1$ maps $(E_\tau,(P,Q))$ to $(E_{\tau'},(\frac1N,\frac{\tau'}N))$: this is surjectivity. Conversely an isomorphism
$z\mapsto\alpha z$ from $(E_\tau,(\frac1N,\frac\tau N))$ to $(E_{\tau'},(\frac1N,\frac{\tau'}N))$ satisfies $\alpha=m_{11}+m_{12}\tau'$, $\alpha\tau=m_{21}+m_{22}\tau'$ with $\sm{m_{11}}{m_{12}}{m_{21}}{m_{22}}\equiv I\pmod N$ of
determinant $1$ (it maps the oriented basis $(1,\tau)$ to an oriented basis of $\Lambda_{\tau'}$), so $\tau=\frac{m_{22}\tau'+m_{21}}{m_{12}\tau'+m_{11}}$ with $\sm{m_{22}}{m_{21}}{m_{12}}{m_{11}}\in\Gamma(N)$: this is injectivity.
That $\zeta_N$ does not depend on $\tau$ is part (iv) of the proposition on isogenies of complex tori in \Cref{ch:ellcurves} ($e_N=\zeta^{\det}$ on $\frac1N\Lambda/\Lambda$ for an
oriented basis).
(iii) The point $\pi(\Gamma_0(N)\tau)=\SL_2(\Z)\tau$ of $Y(1)$ corresponds to the isomorphism class of $E_\tau$ alone (the case $N=1$ of (i)), so $\pi$ forgets $C$. For the second statement, $E_\tau/\langle\frac1N\rangle=\C/(\frac1N\Z+\Z\tau)$, and multiplication by $N$ identifies it with $\C/\Lambda_{N\tau}$, sending $E_\tau[N]/C$ to the
subgroup generated by $\tau=\frac{N\tau}N$. Multiplication by $(N\tau)^{-1}$ identifies $(\C/\Lambda_{N\tau},\langle\frac{N\tau}N\rangle)$ with $(\C/\Lambda_{-1/(N\tau)},\langle\frac1N\rangle)$,
since $(N\tau)^{-1}\Lambda_{N\tau}=\Z(N\tau)^{-1}+\Z=\Lambda_{-1/(N\tau)}$.
(iv), (v) We do not prove these; see \cite{DeligneRapoport73} and \cite{KatzMazur85} for the moduli problem over $\Z$ and its compactification,
\cite[Sections 7.7 and 1.5]{DiamondShurman} and \cite[Section 6.7]{Shimura71} for the model over $\Q$, and \cite{Ogg73} for the fields of definition of the cusps (all
section numbers not verified).
\end{proof}

\begin{paperlink}\label{pl:modforms:x0}
The paper's Section 1 speaks of ``the $X_0(N)$-optimal curve, for which the modular parametrisation $\varphi_E\colon X_0(N)\to E$ has minimal degree''. Here $X_0(N)$ is
the curve of \Cref{thm:modforms:x0-compact}, defined over $\Q$ by \Cref{thm:modforms:moduli}(v), and $\varphi_E$ sends the rational cusp $\infty$ to the origin
(\Cref{ch:moddeg}). For $N=11$ the curve itself has genus $1$ (the table above), and is the curve \curve{11a1}, so that $\deg\varphi_{\curve{11a1}}=1$. In Section 4.1
the paper uses $H^1(X_0(M^\Sigma),\Z_\ell)$, the cohomology of the compact Riemann surface $X_0(M^\Sigma)$ (\Cref{ch:curves}); its rank is $2g(X_0(M^\Sigma))$, for example
$2g(X_0(121))=12$.
\end{paperlink}

\section{Modular forms, \texorpdfstring{$q$}{q}-expansions and nebentypus}\label{sec:modforms:forms}

This section defines the objects of the paper's first sentence about modular forms, ``let $f=f_E\in S_2(\Gamma_0(N))$ be the newform attached to it'': the
spaces $M_k(\Gamma_0(N))$ and $S_k(\Gamma_0(N))$, holomorphy and vanishing at the cusps, and the $q$-expansion $\sum a_n(f)q^n$, whose coefficients are the numbers
compared in every congruence of the paper. It also proves the elementary growth bound for the $a_n$, states the Ramanujan-Petersson bound, and
introduces forms with a character (nebentypus), which occur in the optimal-level theorem used in the paper's Section 3.

\begin{definition}[the slash operator]\label{def:modforms:slash}
For $k\in\Z$, $\gamma=\sm abcd\in\GL_2^+(\R)$ and $f\colon\HH\to\C$, put $(f|_k\gamma)(z)=(\det\gamma)^{k/2}(cz+d)^{-k}f(\gamma z)$ (\Cref{conv:modforms:conventions}(d)).
\end{definition}

\begin{lemma}[properties of the slash operator]\label{lem:modforms:slash}
For $\gamma,\gamma'\in\GL_2^+(\R)$ and $f\colon\HH\to\C$:
\begin{enumerate}[label=(\roman*)]
\item $(f|_k\gamma)|_k\gamma'=f|_k(\gamma\gamma')$;
\item $f|_k(\lambda I)=\sgn(\lambda)^kf$ for $\lambda\in\R^\times$;
\item in weight $2$, $(f|_2\gamma)(z)\,dz=\gamma^*\bigl(f(z)\,dz\bigr)$, the pullback of the differential $f(z)\,dz$ under $z\mapsto\gamma z$;
\item $\operatorname{Im}(\gamma z)^{k/2}\abs{f(\gamma z)}=\operatorname{Im}(z)^{k/2}\abs{(f|_k\gamma)(z)}$; so if $f|_k\gamma=f$ for all $\gamma$ in a group $\Gamma$, the function $y^{k/2}\abs{f}$ is $\Gamma$-invariant.
\end{enumerate}
\end{lemma}

\begin{proof}
(i) $(f|_k\gamma)|_k\gamma'(z)=(\det\gamma')^{k/2}j(\gamma',z)^{-k}(\det\gamma)^{k/2}j(\gamma,\gamma'z)^{-k}f(\gamma\gamma'z)$, which is $f|_k(\gamma\gamma')(z)$ by the cocycle identity
\Cref{prop:modforms:upper-half-plane}(iii) and $\det(\gamma\gamma')=\det\gamma\det\gamma'$. (ii) $(\lambda^2)^{k/2}\lambda^{-k}=\abs\lambda^k\lambda^{-k}$. (iii) $\gamma^*(f(z)\,dz)=f(\gamma z)\,d(\gamma z)=
f(\gamma z)\det\gamma\,(cz+d)^{-2}\,dz$ by \Cref{prop:modforms:upper-half-plane}(ii). (iv) Take absolute values and use \Cref{prop:modforms:upper-half-plane}(i):
$\operatorname{Im}(\gamma z)^{k/2}=(\det\gamma)^{k/2}y^{k/2}\abs{cz+d}^{-k}$.
\end{proof}

With Diamond-Shurman's $(\det\gamma)^{k-1}$, (i), (ii) up to a power of $\abs\lambda$, and (iii) hold, but (iv) and the isometry property of
\Cref{prop:modforms:atkin-lehner-petersson} need the exponent $k/2$. In weight $2$ the two normalisations coincide, which is why the paper's operator $w$ of
Section 4.1 is the Atkin-Lehner involution ``with no scalar factor'' in either convention.

\begin{definition}[modular forms and cusp forms]\label{def:modforms:modular-form}
Let $\Gamma$ be a congruence subgroup and $k\in\Z$. A \emph{modular form of weight $k$ for $\Gamma$} is a holomorphic function $f\colon\HH\to\C$ such that
\begin{enumerate}[label=(\alph*)]
\item $f|_k\gamma=f$ for all $\gamma\in\Gamma$, and
\item $f$ is \emph{holomorphic at the cusps}: for every $\sigma\in\SL_2(\Z)$ the function $f|_k\sigma$ is bounded on $\{\operatorname{Im}z\ge1\}$.
\end{enumerate}
It is a \emph{cusp form} if moreover $(f|_k\sigma)(z)\to0$ as $\operatorname{Im}z\to\infty$, for every $\sigma\in\SL_2(\Z)$. The modular forms form a complex vector space $M_k(\Gamma)$,
the cusp forms a subspace $S_k(\Gamma)$. The \emph{level} of a form for $\Gamma_0(N)$ is $N$; the paper's ``cusp form of level $N$'' is an element of $S_2(\Gamma_0(N))$.
\end{definition}

Condition (a) says that $f(z)\,(dz)^{k/2}$ is invariant under $\Gamma$ (for $k=2$ this is \Cref{lem:modforms:slash}(iii) below), so that $f$ lives on the quotient $Y(\Gamma)$ in this twisted
sense. Condition (b) is the analogue of ``holomorphic'' at the points that $Y(\Gamma)$ is missing. A cusp $s=\sigma\infty$ is moved to $\infty$ by $\sigma^{-1}$, and the behaviour of $f$
near $s$ is the behaviour of $f|_k\sigma$ near $\infty$; so the condition must be imposed for every $\sigma$, not only for $\sigma=I$. It is not automatic: after
\Cref{thm:modforms:level-eleven} we exhibit a holomorphic function of weight $12$ for $\Gamma_0(11)$ that vanishes at $\infty$ to order $21$ but has a pole at the cusp $0$. Since
$\SL_2(\Z)$ has the single cusp $\infty$, at level one only $\sigma=I$ matters.

\begin{proposition}[first properties]\label{prop:modforms:mf-basic}
Let $\Gamma\subseteq\SL_2(\Z)$ be a congruence subgroup.
\begin{enumerate}[label=(\alph*)]
\item Condition (b) of \Cref{def:modforms:modular-form}, and the cusp condition, at $\sigma$ depend only on the cusp $\Gamma\sigma\infty$.
\item If $-I\in\Gamma$ and $k$ is odd, $M_k(\Gamma)=0$. In particular only even weights occur for $\Gamma_0(N)$.
\item If $\Gamma'\subseteq\Gamma$ then $M_k(\Gamma)\subseteq M_k(\Gamma')$ and $S_k(\Gamma)\subseteq S_k(\Gamma')$.
\item $M_k(\Gamma)M_l(\Gamma)\subseteq M_{k+l}(\Gamma)$.
\item Let $f\in M_k(\Gamma_0(M))$ and $d\ge1$. The function $f(dz)=d^{-k/2}\bigl(f|_k\sm d001\bigr)(z)$ lies in $M_k(\Gamma_0(Md))$, and in $S_k(\Gamma_0(Md))$ if $f$ is a cusp form.
\end{enumerate}
\end{proposition}

\begin{proof}
(a) If $\sigma'\infty=\gamma\sigma\infty$ with $\gamma\in\Gamma$, then $\sigma'=\gamma\sigma(\pm T^m)$, and $f|_k\sigma'=\pm f|_k\sigma|_kT^m$ is a translate of $\pm f|_k\sigma$, which preserves
$\{\operatorname{Im}z\ge1\}$. (b) $f=f|_k(-I)=(-1)^kf$ by \Cref{lem:modforms:slash}(ii). (c) is clear. (d) $(fg)|_{k+l}\gamma=(f|_k\gamma)(g|_l\gamma)$, and products of bounded
functions are bounded. (e) Let $\alpha=\sm d001$. For $\gamma=\sm abce\in\Gamma_0(Md)$, $\alpha\gamma=\gamma'\alpha$ with $\gamma'=\sm a{db}{c/d}e\in\Gamma_0(M)$, so
$f|_k\alpha|_k\gamma=f|_k\gamma'|_k\alpha=f|_k\alpha$. For $\sigma\in\SL_2(\Z)$ write $\alpha\sigma=\sigma'\beta$ with $\sigma'\in\SL_2(\Z)$ and $\beta=\sm{a'}{b'}0{d'}$ upper triangular with
integer entries. To find them, write the first column of the integer matrix $\alpha\sigma$ as $a'v$ with $a'\ge1$ and $v\in\Z^2$ primitive, complete $v$ to a matrix
$\sigma'\in\SL_2(\Z)$ with first column $v$ (possible because $v$ is primitive), and put $\beta=\sigma'^{-1}\alpha\sigma$. It has integer entries, its first column is
$\sigma'^{-1}(a'v)=(a',0)^{\mathsf T}$, and $\det\beta=\det\alpha=d$, so $\beta=\sm{a'}{b'}0{d'}$ with $a'd'=d$ and $d'\ge1$. Then $f|_k\alpha\sigma=(f|_k\sigma')|_k\beta$, and
$f|_k\sigma'=\sum_{n\ge0}b_nq_h^n$ (\Cref{prop:modforms:q-expansion}(iii) below) is bounded on every half plane $\{\operatorname{Im}z\ge C\}$, $C>0$, and tends to $b_0$; since $\beta$ maps
$\{\operatorname{Im}z\ge1\}$ to $\{\operatorname{Im}z\ge a'/d'\}$, $f|_k\alpha\sigma$ is bounded there, and tends to $0$ if $b_0=0$.
\end{proof}

The forms $f(dz)$ of (e) are the raw material of \Cref{ch:hecke}'s oldforms; the paper's Lemma lem:old works with ``the forms $h(dz)$, $d\mid N/M$'' for a newform
$h$ of level $M$, which lie in $S_2(\Gamma_0(N))$ by (e).

\begin{proposition}[$q$-expansions]\label{prop:modforms:q-expansion}
Let $f\in M_k(\Gamma_0(N))$.
\begin{enumerate}[label=(\roman*)]
\item $f(z+1)=f(z)$, and $f(z)=\sum_{n\ge0}a_n(f)q^n$, $q=e^{2\pi iz}$, for all $z\in\HH$, the series converging absolutely and uniformly on $\{\operatorname{Im}z\ge y_0\}$ for every $y_0>0$.
For every $y>0$,
\[
a_n(f)=\int_0^1f(x+iy)\,e^{-2\pi in(x+iy)}\,dx .
\]
\item More generally, a holomorphic $g\colon\HH\to\C$ of period $1$ is $g(z)=G(q)$ for a holomorphic $G$ on the punctured disc $0<\abs q<1$, with a Laurent expansion
$G=\sum_{n\in\Z}a_nq^n$; $g$ is bounded on $\{\operatorname{Im}z\ge1\}$ iff $G$ has a removable singularity at $0$ iff $a_n=0$ for all $n<0$. Then we write $g(\infty)=a_0$.
\item For $\sigma\in\SL_2(\Z)$ with $s=\sigma\infty$ of width $h$, $f|_k\sigma$ has period $h$ and an expansion $f|_k\sigma=\sum_{n\ge0}b_n(\sigma)q_h^n$, the \emph{expansion of $f$ at $s$}. The form
$f$ is a cusp form iff $b_0(\sigma)=0$ for a set of $\sigma$ representing all cusps.
\item ($q$-expansion principle, analytic form.) The map $f\mapsto(a_n(f))_{n\ge0}$ is injective on $M_k(\Gamma_0(N))$.
\end{enumerate}
\end{proposition}

\begin{proof}
(ii) Since $z\mapsto q=e^{2\pi iz}$ maps $\HH$ onto the punctured disc and identifies exactly the points that differ by integers, $G(q)=g(z)$ is well defined and holomorphic (locally $z=\frac1{2\pi i}\log q$).
Its Laurent series converges absolutely and locally uniformly on the punctured disc. The set $\{\operatorname{Im}z\ge1\}$ corresponds to $0<\abs q\le e^{-2\pi}$, and Riemann's
theorem on removable singularities gives the equivalences.
(i) $T\in\Gamma_0(N)$ gives the period, and (ii) the expansion with $n\ge0$ (taking $\sigma=I$ in \Cref{def:modforms:modular-form}(b)). A power series with radius of
convergence at least $1$ converges absolutely and uniformly on $\abs q\le e^{-2\pi y_0}<1$. The coefficient formula is Cauchy's formula
$a_n=\frac1{2\pi i}\oint_{\abs q=e^{-2\pi y}}G(q)q^{-n-1}\,dq$ with $q=e^{2\pi i(x+iy)}$, $dq=2\pi iq\,dx$.
(iii) $\sigma T^h\sigma^{-1}\in\pm\Gamma_0(N)=\Gamma_0(N)$, so $f|_k\sigma|_kT^h=f|_k\sigma$, that is, $f|_k\sigma$ has period $h$; apply (ii) to $z\mapsto(f|_k\sigma)(hz)$. The cusp condition at $\sigma$ is
$b_0(\sigma)=0$, and depends only on the cusp by \Cref{prop:modforms:mf-basic}(a).
(iv) By (i) the coefficients determine $f$ on all of $\HH$.
\end{proof}

A much stronger statement holds: $f$ is determined by $a_0,\dots,a_{B_k(N)}$ (\Cref{thm:modforms:valence-gamma0}, and \Cref{thm:modforms:sturm} for congruences).

The expansion is a Fourier series. A function on $\HH$ of period $1$ in $x$ has a Fourier series $\sum_nc_n(y)e^{2\pi inx}$; holomorphy (the Cauchy-Riemann equation
$\partial_yf=i\partial_xf$) forces $c_n'(y)=-2\pi nc_n(y)$, so $c_n(y)=a_ne^{-2\pi ny}$, and the $n$-th mode is $a_ne^{2\pi inx}e^{-2\pi ny}=a_nq^n$. These are the separated solutions of
Laplace's equation on a half-strip that stay bounded as $y\to\infty$ exactly when $n\ge0$; the condition ``holomorphic at $\infty$'' is the boundary condition that kills the
growing modes. So a modular form is a Fourier series on the quotient $\HH/\langle T\rangle$, a half-infinite cylinder, with one more symmetry, $z\mapsto-1/(Nz)$ or the rest
of $\Gamma_0(N)$, that the coefficients must respect.

\begin{paperlink}\label{pl:modforms:coefficientwise}
Every congruence in the paper is a statement about the coefficients $a_n$ of \Cref{prop:modforms:q-expansion}(i). In the opening paragraph of Section 3 the
congruence number is ``the largest integer $r$ such that $f\equiv g\pmod r$ coefficientwise for some $g\in S_2(\Gamma_0(N))$ with integer coefficients orthogonal to
$f$'': coefficientwise means $a_n(f)\equiv a_n(g)\pmod r$ for all $n\ge0$ (\Cref{conv:modforms:conventions}(g)); integer coefficients means $g\in S_2(\Gamma_0(N);\Z)$
(\Cref{def:modforms:integral-forms}); orthogonality is for the Petersson product (\Cref{def:modforms:petersson}). Lemma lem:crit asks for $a_n(g)\equiv a_n(f)\pmod\lambda$
``for every $n\ge1$'', while Lemma lem:old starts from a congruence at the primes $p\nmid\ell N$ only and passes to all $n$ through the Hecke recursion
(\Cref{ch:hecke}).
\end{paperlink}

\begin{proposition}[the Hecke bound]\label{prop:modforms:hecke-bound}
Let $f\in S_k(\Gamma_0(N))$. Then $y^{k/2}\abs{f(z)}$ is bounded on $\HH$, and there is a constant $C_f$ with $\abs{a_n(f)}\le C_fn^{k/2}$ for all $n\ge1$. Consequently the
Dirichlet series $\sum_{n\ge1}a_n(f)n^{-s}$ converges absolutely for $\operatorname{Re}s>\frac k2+1$.
\end{proposition}

\begin{proof}
The function $\phi=y^{k/2}\abs f$ is continuous and $\Gamma_0(N)$-invariant (\Cref{lem:modforms:slash}(iv)), so it suffices to bound it on the fundamental domain
$\bigcup_i\gamma_i\mathcal F$ of \Cref{prop:modforms:coset-domains}. For $w\in\mathcal F$, $\phi(\gamma_iw)=\operatorname{Im}(w)^{k/2}\abs{(f|_k\gamma_i)(w)}$ by \Cref{lem:modforms:slash}(iv). On $\mathcal F$,
$\operatorname{Im}w\ge\sqrt3/2$, and $f|_k\gamma_i=\sum_{n\ge1}b_nq_h^n$ without constant term, so $\abs{(f|_k\gamma_i)(w)}\le\abs{q_h}\sum_{n\ge1}\abs{b_n}r_0^{n-1}=C_ie^{-2\pi\operatorname{Im}w/h}$ with
$r_0=e^{-\pi\sqrt3/h}<1$. Since $y^{k/2}e^{-2\pi y/h}$ is bounded, $\phi$ is bounded, say by $C$. Then by \Cref{prop:modforms:q-expansion}(i) at height $y=1/n$:
$\abs{a_n(f)}\le\int_0^1\abs{f(x+i/n)}e^{2\pi}\,dx\le Cn^{k/2}e^{2\pi}$. Here $\abs{e^{-2\pi in(x+i/n)}}=e^{2\pi n\cdot(1/n)}=e^{2\pi}$, and $\abs{f(x+iy)}\le Cy^{-k/2}=Cn^{k/2}$ at
$y=1/n$. The choice $y=1/n$ balances the two factors: a larger $y$ makes $e^{2\pi ny}$ grow exponentially in $n$, a smaller one makes $y^{-k/2}$ worse. The last statement
follows by comparison with $\sum n^{k/2-\operatorname{Re}s}$, which converges for $\operatorname{Re}s-\frac k2>1$.
\end{proof}

For forms that are not cuspidal the bound fails: the Eisenstein series $E_k$ of \Cref{prop:modforms:eisenstein-level-one} has $\abs{a_n}=\frac{2k}{\abs{B_k}}\sigma_{k-1}(n)\ge\frac{2k}{\abs{B_k}}n^{k-1}$.
The true size of the coefficients of an eigenform is much smaller.

\begin{theorem}[Ramanujan-Petersson; Eichler-Shimura, Igusa, Deligne]\label{thm:modforms:ramanujan-petersson}
Let $f\in S_k(\Gamma_0(N))$ be a normalised newform (\Cref{ch:hecke}) and $p\nmid N$ a prime. Then $\abs{a_p(f)}\le2p^{(k-1)/2}$; equivalently the roots $\alpha_p,\beta_p$ of
$X^2-a_p(f)X+p^{k-1}$ satisfy $\abs{\alpha_p}=\abs{\beta_p}=p^{(k-1)/2}$, so that $\beta_p=\bar\alpha_p$ and $\abs{\alpha_p/\beta_p}=1$. Moreover $a_n(f)=O(n^{(k-1)/2+\varepsilon})$ for every $\varepsilon>0$.
\end{theorem}

We do not prove this; for $k=2$ it follows from the Eichler-Shimura relation and Weil's Riemann hypothesis for curves over finite fields (\Cref{ch:curves},
\Cref{ch:moddeg}), and for $f=f_E$ it is Hasse's bound $\abs{a_p(E)}\le2\sqrt p$ (\Cref{ch:ellcurves}); in general it is Deligne's theorem \cite{Deligne}. The statement and
references are in \cite[Section 8.7]{DiamondShurman} (section number not verified). Checks: $\abs{a_p(\curve{11a1})}\le2\sqrt p$ for all $p<10^5$ (the largest ratio is
$0.9940$), and Ramanujan's $\abs{\tau(p)}\le2p^{11/2}$ for all $p<2000$ (largest ratio $0.9594$). The equivalence: $\alpha_p\beta_p=p^{k-1}$ and $\alpha_p+\beta_p=a_p$ real, so the roots
are complex conjugate of absolute value $p^{(k-1)/2}$ iff the discriminant $a_p^2-4p^{k-1}$ is $\le0$.

\begin{paperlink}\label{pl:modforms:ramanujan}
Lemma lem:euler and Theorem thm:dfg14 use the roots $\alpha,\beta$ of $X^2-a_pX+p$ at a good prime $p$ ($k=2$). Wiles's
level-raising factor $(1-\frac\alpha\beta p^{-1})(1-p^{-1})(1-\frac\beta\alpha p^{-1})$ of the paragraph after Theorem thm:dfg14 is a nonzero real number because $\beta=\bar\alpha$ and
$\abs{\alpha/\beta}=1$: then $\frac\beta\alpha=(\frac\alpha\beta)^{-1}=\overline{\alpha/\beta}$, so the third factor is the complex conjugate of the first, and the product equals
$(1-p^{-1})\abs{1-\frac\alpha\beta p^{-1}}^2$, which is positive because $\abs{\frac\alpha\beta p^{-1}}=p^{-1}<1$.
\end{paperlink}

\begin{definition}[nebentypus]\label{def:modforms:nebentypus}
Let $\chi$ be a Dirichlet character modulo $N$. The space of modular forms of weight $k$, level $N$ and \emph{character} (or \emph{nebentypus}) $\chi$ is
\[
M_k(\Gamma_0(N),\chi)=\Bigl\{f\in M_k(\Gamma_1(N)):f|_k\gamma=\chi(d)f\ \text{for all }\gamma=\sm abcd\in\Gamma_0(N)\Bigr\},
\]
and $S_k(\Gamma_0(N),\chi)$ is its subspace of cusp forms. For $d\in(\Z/N)^\times$ the \emph{diamond operator} $\langle d\rangle$ on $M_k(\Gamma_1(N))$ is $f\mapsto f|_k\gamma$ for any
$\gamma\in\Gamma_0(N)$ with lower right entry $\equiv d\pmod N$.
\end{definition}

The reason for this definition is that $\Gamma_1(N)$, not $\Gamma_0(N)$, is the natural group for many constructions, and $\Gamma_0(N)$ acts on $M_k(\Gamma_1(N))$ through the finite
abelian quotient $\Gamma_0(N)/\Gamma_1(N)\simeq(\Z/N)^\times$. Decomposing under a finite abelian group means decomposing by characters, and $M_k(\Gamma_0(N),\chi)$ is the
$\chi$-isotypic piece; the trivial character gives back $M_k(\Gamma_0(N))$. Nontrivial characters genuinely occur: PARI gives $\dim S_2(\Gamma_1(13))=2$ while
$S_2(\Gamma_0(13))=0$ (the genus of $X_0(13)$ is $0$), the two dimensions coming from a character of order $6$ modulo $13$ and its complex conjugate, one dimension each.

\begin{proposition}[decomposition by characters]\label{prop:modforms:nebentypus}
\begin{enumerate}[label=(\roman*)]
\item $M_k(\Gamma_0(N),\mathbf 1)=M_k(\Gamma_0(N))$ for the trivial character $\mathbf 1$.
\item $M_k(\Gamma_0(N),\chi)=0$ unless $\chi(-1)=(-1)^k$.
\item The diamond operators are well defined and give an action of $(\Z/N)^\times$ on $M_k(\Gamma_1(N))$, and
$M_k(\Gamma_1(N))=\bigoplus_\chi M_k(\Gamma_0(N),\chi)$, the sum over all Dirichlet characters modulo $N$; likewise for $S_k$.
\end{enumerate}
\end{proposition}

\begin{proof}
(i) $\chi(d)=1$. (ii) Apply the definition to $\gamma=-I$: $(-1)^kf=\chi(-1)f$. (iii) $\Gamma_1(N)$ is normal in $\Gamma_0(N)$ with quotient $(\Z/N)^\times$ via $\gamma\mapsto d$
(\Cref{prop:modforms:indices}(ii)). For $f\in M_k(\Gamma_1(N))$ and $\gamma\in\Gamma_0(N)$, $f|_k\gamma$ is again invariant under $\Gamma_1(N)=\gamma\Gamma_1(N)\gamma^{-1}$ and holomorphic at the cusps
($f|_k\gamma|_k\sigma=f|_k(\gamma\sigma)$ with $\gamma\sigma\in\SL_2(\Z)$), and depends only on $\gamma\Gamma_1(N)$; so the finite abelian group $(\Z/N)^\times$ acts linearly on
$M_k(\Gamma_1(N))$. The operators $e_\chi=\frac1{\varphi(N)}\sum_{d\in(\Z/N)^\times}\overline{\chi(d)}\langle d\rangle$ are orthogonal idempotents with $\sum_\chi e_\chi=1$ and
$\langle d\rangle e_\chi=\chi(d)e_\chi$ (orthogonality of characters, \Cref{ch:reps}). In detail: substituting $d'\mapsto d^{-1}d'$,
\[
\langle d\rangle e_\chi=\frac1{\varphi(N)}\sum_{d'}\overline{\chi(d')}\langle dd'\rangle=\frac{\chi(d)}{\varphi(N)}\sum_{d'}\overline{\chi(dd')}\langle dd'\rangle=\chi(d)e_\chi ;
\]
hence $e_\psi e_\chi=\frac1{\varphi(N)}\sum_d\overline{\psi(d)}\chi(d)e_\chi$, which is $e_\chi$ if $\psi=\chi$ and $0$ otherwise; and $\sum_\chi e_\chi=\frac1{\varphi(N)}\sum_d\bigl(\sum_\chi\overline{\chi(d)}\bigr)\langle d\rangle=\langle1\rangle=1$.
So $M_k(\Gamma_1(N))$ is the direct sum of the images $e_\chi M_k(\Gamma_1(N))$. Each image is $M_k(\Gamma_0(N),\chi)$: the image consists of $\chi$-eigenvectors of every $\langle d\rangle$, which is
the defining condition $f|_k\gamma=\chi(d)f$, and conversely $e_\chi f=f$ for such $f$. The operators preserve cusp forms.
\end{proof}

The twist of $f\in S_k(\Gamma_0(M))$ by a primitive Dirichlet character $\chi$ of conductor $r$ is a form of character $\chi^2$ and level dividing $\operatorname{lcm}(M,r^2)$
(\Cref{ch:hecke}); for the quadratic characters $\chi_{q^*}$ of the paper $\chi^2=\mathbf 1$, so twists stay in the spaces $S_2(\Gamma_0(\cdot))$ with trivial character.

\begin{paperlink}\label{pl:modforms:nebentypus}
The optimal-level theorem quoted in the paper's Section 3 produces ``a newform of weight $2$ and level $N(\bar\sigma)\ell^{\delta(\bar\sigma)}$ with nebentypus of order prime
to $\ell$'', an element of $S_2(\Gamma_0(M),\chi)$ in the sense of \Cref{def:modforms:nebentypus}. The paper then observes that ``$\det\bar\sigma=\bar\chi_\ell$ forces the
nebentypus, which is trivial modulo $\lambda$ and of order prime to $\ell$, to be trivial'', so that the newform lies in $S_2(\Gamma_0(M))$; the Galois-theoretic content of
that sentence (the determinant of $\rho_f$ is $\chi\chi_\ell$) is \Cref{ch:modgalrep}'s.
\end{paperlink}

\section{Level one: Eisenstein series, \texorpdfstring{$\Delta$}{Delta}, \texorpdfstring{$j$}{j} and the valence formula}\label{sec:modforms:level-one}

Level one is where every structural fact can be proved with a contour integral. This section constructs the Eisenstein series $E_k$, the cusp form $\Delta$ and
the function $j$, proves the valence formula, and deduces the structure of $M_*(\SL_2(\Z))$ and the bijectivity of $j$. It discharges two debts of
\Cref{ch:ellcurves}: the $q$-expansion of $j$ and the bijection $j\colon\SL_2(\Z)\backslash\HH\to\C$. The paper uses level one only through these objects: the $j$-invariant $j_E$ of
\Cref{ch:ellcurves}, whose valuations enter Definition def:terms, and the Eisenstein series behind the Eisenstein case of Section sec:eis (\Cref{prop:modforms:eisenstein-weight-two}).

\begin{proposition}[Eisenstein series of level one]\label{prop:modforms:eisenstein-level-one}
For even $k\ge4$ the series $G_k(z)=\sum_{(m,n)\in\Z^2\smallsetminus\{0\}}(mz+n)^{-k}$ converges absolutely and locally uniformly on $\HH$, $G_k\in M_k(\SL_2(\Z))$, and $G_k=2\zeta(k)E_k$ with
\[
E_k(z)=1-\frac{2k}{B_k}\sum_{n\ge1}\sigma_{k-1}(n)q^n,\qquad E_4=1+240\sum_{n\ge1}\sigma_3(n)q^n,\qquad E_6=1-504\sum_{n\ge1}\sigma_5(n)q^n,
\]
where $\sigma_r(n)=\sum_{d\mid n}d^r$ and $B_k$ is the Bernoulli number. In the notation of \Cref{ch:ellcurves}, $G_k(z)=G_k(\Lambda_z)$. The series
$E_2(z)=1-24\sum_{n\ge1}\sigma_1(n)q^n$ converges on $\HH$ but is not modular: it satisfies
\[
E_2(-1/z)=z^2E_2(z)+\frac{6z}{\pi i},
\]
equivalently $E_2^*(z)=E_2(z)-\frac3{\pi y}$ satisfies $E_2^*|_2\gamma=E_2^*$ for all $\gamma\in\SL_2(\Z)$.
\end{proposition}

\begin{proof}
Absolute and locally uniform convergence of $\sum\abs{mz+n}^{-k}$ for $k>2$ is the lattice-sum lemma of \Cref{ch:ellcurves}, so $G_k$ is holomorphic. For
$\gamma=\sm abcd\in\SL_2(\Z)$,
\[
G_k(\gamma z)=\sum\Bigl(\frac{m(az+b)+n(cz+d)}{cz+d}\Bigr)^{-k}=(cz+d)^k\sum\bigl((ma+nc)z+(mb+nd)\bigr)^{-k},
\]
which is $(cz+d)^kG_k(z)$ since $(m,n)\mapsto(ma+nc,mb+nd)$ permutes $\Z^2\smallsetminus\{0\}$. For the expansion we use the Lipschitz formula, valid for $k\ge2$ and $z\in\HH$:
\[
\sum_{n\in\Z}(z+n)^{-k}=\frac{(-2\pi i)^k}{(k-1)!}\sum_{r\ge1}r^{k-1}q^r .
\]
It follows from the partial fraction expansion $\pi\cot\pi z=\frac1z+\sum_{n\ne0}\bigl(\frac1{z+n}-\frac1n\bigr)$ and $\pi\cot\pi z=\pi i\frac{q+1}{q-1}=-\pi i-2\pi i\sum_{r\ge1}q^r$ (for $\abs q<1$), by
differentiating both $k-1$ times. In detail: $\pi\cot\pi z=\pi i\frac{e^{\pi iz}+e^{-\pi iz}}{e^{\pi iz}-e^{-\pi iz}}=\pi i\frac{q+1}{q-1}$ (multiply numerator and denominator by $e^{\pi iz}$), and
$\frac{q+1}{q-1}=-(1+q)\sum_{r\ge0}q^r=-1-2\sum_{r\ge1}q^r$. Both expansions converge locally uniformly, so they may be differentiated termwise. The $(k-1)$-st derivative of
$\frac1{z+n}$ is $(-1)^{k-1}(k-1)!\,(z+n)^{-k}$, the constants $-\frac1n$ disappear, and that of $q^r=e^{2\pi irz}$ is $(2\pi ir)^{k-1}q^r$. Hence
$(-1)^{k-1}(k-1)!\sum_n(z+n)^{-k}=-2\pi i\sum_r(2\pi ir)^{k-1}q^r=-(2\pi i)^k\sum_rr^{k-1}q^r$, which is the Lipschitz formula. Separate the terms $m=0$, which give $2\zeta(k)$, and use the symmetry $(m,n)\mapsto(-m,-n)$ ($k$ even):
\[
\begin{aligned}
G_k(z)&=2\zeta(k)+2\sum_{m\ge1}\sum_{n\in\Z}(mz+n)^{-k}=2\zeta(k)+\frac{2(2\pi i)^k}{(k-1)!}\sum_{m\ge1}\sum_{r\ge1}r^{k-1}q^{mr}\\
&=2\zeta(k)+\frac{2(2\pi i)^k}{(k-1)!}\sum_{n\ge1}\sigma_{k-1}(n)q^n.
\end{aligned}
\]
(In the first line, the inner sum for fixed $m$ is the Lipschitz formula at the point $mz$, with $e^{2\pi irmz}=q^{mr}$ and $(-2\pi i)^k=(2\pi i)^k$ as $k$ is even; in the second, the
pairs $(m,r)$ with $mr=n$ contribute $\sum_{r\mid n}r^{k-1}=\sigma_{k-1}(n)$.) With Euler's $\zeta(k)=-\frac{(2\pi i)^kB_k}{2\,k!}$ we get
$2\zeta(k)\cdot\bigl(-\frac{2k}{B_k}\bigr)=\frac{2k(2\pi i)^k}{k!}=\frac{2(2\pi i)^k}{(k-1)!}$, so the displayed expansion is $2\zeta(k)E_k$. So $G_k$ is bounded on $\{\operatorname{Im}z\ge1\}$, and since $G_k|_k\sigma=G_k$ for every $\sigma\in\SL_2(\Z)$, it is holomorphic at the cusp;
$B_4=-\frac1{30}$ and $B_6=\frac1{42}$ give $240$ and $-504$. The transformation law of $E_2$ is \Cref{prop:modforms:eta} below, combined with
$\frac{\eta'}\eta=\frac{\pi i}{12}E_2$: differentiating $\log\eta(-1/z)=\frac12\log\frac zi+\log\eta(z)$ gives $z^{-2}\frac{\eta'}\eta(-1/z)=\frac1{2z}+\frac{\eta'}\eta(z)$, which is the stated
identity. Then $\operatorname{Im}(-1/z)=y/\abs z^2$ gives $E_2^*(-1/z)=z^2E_2^*(z)+\frac{6z}{\pi i}+\frac{3(z^2-\abs z^2)}{\pi y}=z^2E_2^*(z)$, because $z^2-\abs z^2=z(z-\bar z)=2iyz$; and
$E_2^*$ is invariant under $T$. As $S$ and $T$ generate $\SL_2(\Z)$ (\Cref{thm:modforms:fundamental-domain}(iv)) and $\gamma\mapsto(\phi\mapsto\phi|_2\gamma)$ is an action (\Cref{lem:modforms:slash}(i),
valid for any function), $E_2^*|_2\gamma=E_2^*$ for all $\gamma$.
\end{proof}

Since $E_2$ rests on the transformation law of $\eta$, which we do not prove, the statements about $E_2$ carry the tag ``given the transformation law of
$\eta$'' on the card.

\begin{proposition}[$\Delta$ and $j$]\label{prop:modforms:delta-j}
\begin{enumerate}[label=(\roman*)]
\item $\Delta=(E_4^3-E_6^2)/1728$ is a nonzero cusp form of weight $12$ for $\SL_2(\Z)$ with integer coefficients,
\[
\Delta=\sum_{n\ge1}\tau(n)q^n=q-24q^2+252q^3-1472q^4+4830q^5-6048q^6+\cdots,
\]
and $\Delta(\Lambda_z)=(2\pi)^{12}\Delta(z)$, where $\Delta(\Lambda)=g_2^3-27g_3^2$ is the discriminant of \Cref{ch:ellcurves}.
\item $j=E_4^3/\Delta=1728E_4^3/(E_4^3-E_6^2)$ is holomorphic on $\HH$, $\SL_2(\Z)$-invariant, equal to $j(\Lambda_z)$ of \Cref{ch:ellcurves}, and
\[
j(z)=q^{-1}+744+196884q+21493760q^2+864299970q^3+20245856256q^4+\cdots,
\]
with all coefficients integers; $j(i)=1728$ and $j(\zeta_3)=0$.
\item (Jacobi) $\Delta(z)=q\prod_{n\ge1}(1-q^n)^{24}=\eta(z)^{24}$.
\end{enumerate}
\end{proposition}

\begin{proof}
(i) $\Delta\in M_{12}$ by \Cref{prop:modforms:mf-basic}(d). Write $E_4=1+240A$, $E_6=1-504B$ with $A=\sum\sigma_3(n)q^n$, $B=\sum\sigma_5(n)q^n$. Then
$E_4^3-E_6^2=144(5A+7B)+1728(100A^2+8000A^3-147B^2)$, since $3\cdot240^2=1728\cdot100$, $240^3=1728\cdot8000$, $504^2=1728\cdot147$ and $3\cdot240=720$, $2\cdot504=1008$.
Moreover $d^5-d^3=d^3(d-1)(d+1)$ is divisible by $12$ for every $d$, so $\sigma_5(n)\equiv\sigma_3(n)\pmod{12}$ and $5A+7B\equiv12A\equiv0\pmod{12}$. So $\Delta$ has integer
coefficients, its constant term is $0$ and $a_1=(720+1008)/1728=1$; it is a cusp form because the only cusp of $\SL_2(\Z)$ is $\infty$. The listed coefficients are computed
from the expansions. In \Cref{ch:ellcurves}, $g_2=60G_4(\Lambda)$ and $g_3=140G_6(\Lambda)$; with $G_4=2\zeta(4)E_4=\frac{\pi^4}{45}E_4$ and $G_6=2\zeta(6)E_6=\frac{2\pi^6}{945}E_6$ this gives
$g_2=\frac43\pi^4E_4$, $g_3=\frac8{27}\pi^6E_6$ and $g_2^3-27g_3^2=\frac{64}{27}\pi^{12}(E_4^3-E_6^2)=\frac{64\cdot1728}{27}\pi^{12}\Delta=(2\pi)^{12}\Delta$.
(ii) $\Delta$ has no zero on $\HH$ (\Cref{cor:modforms:level-one-structure}(ii), which uses only the valence formula), so $j$ is holomorphic and, as a quotient of two forms of
weight $12$, invariant. $j(\Lambda_z)=1728g_2^3/\Delta(\Lambda_z)=1728\cdot\frac{64}{27}\pi^{12}E_4^3/((2\pi)^{12}\Delta)=E_4^3/\Delta$. Since $\Delta=q(1-24q+\cdots)$ has integer coefficients and leading
coefficient $1$, $1/\Delta\in q^{-1}\Z[[q]]$, so $j\in q^{-1}\Z[[q]]$; the coefficients shown were computed with PARI. $E_6(i)=0$ because $E_6(-1/z)=z^6E_6(z)$ at $z=i$ gives $E_6(i)=-E_6(i)$;
so $j(i)=1728E_4(i)^3/E_4(i)^3=1728$, as $E_4(i)\ne0$ (\Cref{cor:modforms:level-one-structure}(ii)). Likewise $E_4(\zeta_3)=0$: the element $ST=\sm0{-1}11$ fixes $\zeta_3$, so
$E_4(\zeta_3)=(\zeta_3+1)^4E_4(\zeta_3)$, and $\zeta_3+1=e^{\pi i/3}$ has $(\zeta_3+1)^4=e^{4\pi i/3}\ne1$. As $\Delta(\zeta_3)\ne0$, $j(\zeta_3)=E_4(\zeta_3)^3/\Delta(\zeta_3)=0$.
(iii) Given the transformation law of $\eta$ (\Cref{prop:modforms:eta}), $\eta^{24}(z+1)=\eta^{24}(z)$ and $\eta^{24}(-1/z)=(z/i)^{12}\eta^{24}(z)=z^{12}\eta^{24}(z)$, so $\eta^{24}\in M_{12}(\SL_2(\Z))$ (as $S,T$
generate); its expansion $q\prod(1-q^n)^{24}=q-24q^2+\cdots$ has no constant term, so $\eta^{24}\in S_{12}(\SL_2(\Z))=\C\Delta$ (\Cref{cor:modforms:level-one-structure}(iii)), and comparing $a_1$ gives
$\eta^{24}=\Delta$.
\end{proof}

\begin{theorem}[the valence formula for $\SL_2(\Z)$]\label{thm:modforms:valence-level-one}
Let $f\ne0$ be a meromorphic function on $\HH$ with $f|_k\gamma=f$ for all $\gamma\in\SL_2(\Z)$, meromorphic at $\infty$ in the sense that $f(z)=\sum_{n\ge n_0}a_nq^n$ for $\operatorname{Im}z$
large. For $p\in\HH$ let $v_p(f)$ be the order of $f$ at $p$ (it depends only on the orbit of $p$), and $v_\infty(f)$ the order in $q$ at $\infty$. Then
\[
v_\infty(f)+\frac12v_i(f)+\frac13v_{\zeta_3}(f)+\sum_{p\in\SL_2(\Z)\backslash\HH,\ p\ne i,\zeta_3}v_p(f)=\frac k{12}.
\]
\end{theorem}

\begin{proof}
The order is constant on orbits because $f(\gamma z)=(cz+d)^kf(z)$ and $cz+d$ has no zeros on $\HH$. Since $f=q^{v_\infty}\cdot(\text{a nonzero power series})$ for $\operatorname{Im}z$ large, the zeros and
poles of $f$ in $\mathcal F$ lie in a compact part $\operatorname{Im}z\le Y_0$, so there are finitely many. Assume first that $f$ has no zeros or poles on the boundary of $\mathcal F$
except possibly at $i$, $\zeta_3$, $-\bar\zeta_3$. Let $Y>Y_0$ and let $\mathcal C$ be the boundary of the region obtained from $\{z\in\mathcal F:\operatorname{Im}z\le Y\}$ by removing small discs
of radius $r$ around $i$, $\zeta_3$ and $-\bar\zeta_3$, oriented counterclockwise. By the argument principle,
\[
\frac1{2\pi i}\oint_{\mathcal C}\frac{f'}f\,dz=\sum_{p\ne i,\zeta_3}v_p(f),
\]
the sum over the orbits with a representative in the interior of $\mathcal F$, one for each. We compute the left side piece by piece.
\begin{itemize}
\item The top segment, from $\frac12+iY$ to $-\frac12+iY$: in the variable $q$ it is a circle around $0$ traversed clockwise, and $\frac{f'}f\,dz=\frac{dF/dq}F\,dq$ with $F(q)=f(z)$; the
integral is $-v_\infty(f)$.
\item The two vertical sides correspond under $T$ and are traversed in opposite directions, and $f(z+1)=f(z)$, so their contributions cancel.
\item The two arcs of the unit circle, from $\zeta_3$ to $i$ (call it $A$) and from $i$ to $-\bar\zeta_3$ (call it $A'$), with small pieces removed at the ends: $S$ maps $A'$ onto $A$
with the orientation reversed. Differentiating $f(Sz)=z^kf(z)$ gives $\frac{f'}f(Sz)\,d(Sz)=\frac{f'}f(z)\,dz+k\frac{dz}z$, so
$\int_{A'}\frac{f'}f\,dz=-\int_A\bigl(\frac{f'}f\,dz+k\frac{dz}z\bigr)$ and the two arcs contribute $-k\int_A\frac{dz}z$. Along the unit circle from the argument $2\pi/3$ to $\pi/2$,
$\int\frac{dz}z=i(\frac\pi2-\frac{2\pi}3)=-\frac{\pi i}6$; dividing by $2\pi i$ gives $\frac k{12}$ in the limit $r\to0$.
\item The small arc around $i$ is a half circle traversed clockwise; as $r\to0$ it contributes $-\frac12v_i(f)$. The small arcs around $\zeta_3$ and $-\bar\zeta_3$ each subtend the angle
$\pi/3$ (the angle of $\mathcal F$ at its corners) and are traversed clockwise; they contribute $-\frac16v_{\zeta_3}(f)-\frac16v_{-\bar\zeta_3}(f)=-\frac13v_{\zeta_3}(f)$.
\end{itemize}
Adding up and letting $r\to0$ gives the formula. If $f$ has a zero or pole $p$ on a vertical side, it has one at the $T$-equivalent point on the other side: deform the
contour by a small arc going into $\mathcal F$ around one of them and out of $\mathcal F$ around the other, so that exactly one of the two is counted and the $T$-symmetry of the
contour is kept; likewise for an $S$-symmetric pair of points on the arc.
\end{proof}

Physically this is Gauss's law for the ``charge density'' $\frac1{2\pi i}\frac{f'}f$: the residue theorem counts the zeros inside the contour, and the symmetries $T$ and $S$
make the flux through the sides cancel, except for the defect $k\frac{dz}z$ produced by the factor of automorphy, which contributes the constant $\frac k{12}$.

\begin{corollary}[the structure of $M_*(\SL_2(\Z))$]\label{cor:modforms:level-one-structure}
\begin{enumerate}[label=(\roman*)]
\item $M_k(\SL_2(\Z))=0$ for $k<0$, for $k$ odd, and for $k=2$; $M_0(\SL_2(\Z))=\C$; $M_k(\SL_2(\Z))=\C E_k$ for $k=4,6,8,10,14$.
\item $\Delta$ has no zero on $\HH$ and a simple zero at $\infty$; $E_4$ vanishes only on the orbit of $\zeta_3$, simply, and $E_6$ only on the orbit of $i$, simply.
\item Multiplication by $\Delta$ is an isomorphism $M_{k-12}(\SL_2(\Z))\to S_k(\SL_2(\Z))$, and for even $k\ge0$, $\dim M_k(\SL_2(\Z))=\lfloor k/12\rfloor$ if $k\equiv2\pmod{12}$ and
$\lfloor k/12\rfloor+1$ otherwise.
\item $\bigoplus_kM_k(\SL_2(\Z))=\C[E_4,E_6]$, with $E_4$ and $E_6$ algebraically independent.
\end{enumerate}
\end{corollary}

\begin{proof}
For $f\in M_k$ nonzero all orders are $\ge0$, so the valence formula reads $n+\frac a2+\frac b3=\frac k{12}$ with integers $n,a,b\ge0$ ($n$ collecting $v_\infty$ and the orders at
the other orbits, $a=v_i$, $b=v_{\zeta_3}$). (i) For $k<0$ there is no solution, and for $k=2$ there is none ($\frac a2+\frac b3=\frac16$ is impossible). Odd $k$ is
\Cref{prop:modforms:mf-basic}(b). For $k=0$, $f-f(z_0)$ has a zero, so it is $0$: $f$ is constant. For $k\in\{4,6,8,10,14\}$ every solution has $v_\infty=0$ (for $k<12$
because $v_\infty\le k/12$, for $k=14$ because $\frac a2+\frac b3=\frac16$ is impossible), so $f\mapsto a_0(f)$ is injective on $M_k$, and $E_k\ne0$.
(ii) For $\Delta$, $v_\infty\ge1$ and the formula with $k=12$ forces $v_\infty=1$ and no other zero. For $E_4$ the only solution of $n+\frac a2+\frac b3=\frac13$ is $b=1$; for $E_6$
the only solution of $=\frac12$ is $a=1$.
(iii) If $f\in S_k$, then $f/\Delta$ is holomorphic on $\HH$ by (ii), of weight $k-12$, and its order at $\infty$ is $v_\infty(f)-1\ge0$; so $f/\Delta\in M_{k-12}$. For $k\ge4$ even,
$M_k=\C E_k\oplus S_k$ since $a_0(E_k)=1$. The dimension formula follows by induction on $k$ from (i) and $\dim M_k=1+\dim M_{k-12}$ for $k\ge4$ (with $M_{k-12}$ for $k=12$ being
$M_0=\C$).
(iv) For even $k\ge4$ choose $a,b\ge0$ with $4a+6b=k$; then $E_4^aE_6^b\in M_k$ has constant term $1$, so $f-a_0(f)E_4^aE_6^b\in S_k=\Delta M_{k-12}$, and
$\Delta=(E_4^3-E_6^2)/1728$; induction on $k$ shows $M_k\subseteq\C[E_4,E_6]$. If a nontrivial polynomial relation held, a homogeneous component of it of some weight $k$ would
be a relation $\sum_{4a+6b=k}c_{ab}E_4^aE_6^b=0$; dividing by a suitable power of $E_6$ (or $E_4$) gives a nonzero polynomial satisfied by the function $E_4^3/E_6^2$, which is
then constant; but $E_4^3/E_6^2$ vanishes at $\zeta_3$ and is $1$ at $\infty$.
\end{proof}

These dimensions agree with PARI's \texttt{mfdim([1,k],4)} for all even $k\le26$.

\begin{theorem}[$j$ is a bijection]\label{thm:modforms:j-bijective}
The function $j$ induces a bijection $\SL_2(\Z)\backslash\HH\to\C$, and extends to an isomorphism of Riemann surfaces $X(1)\to\PP^1(\C)$ with $j(\infty)=\infty$. Consequently, for every
$j_0\in\C$ there is a lattice $\Lambda\subset\C$ with $j(\Lambda)=j_0$.
\end{theorem}

\begin{proof}
For $\lambda\in\C$ the form $f_\lambda=E_4^3-\lambda\Delta\in M_{12}$ has constant term $1$, so $v_\infty(f_\lambda)=0$, and the valence formula reads $n+\frac a2+\frac b3=1$. Its solutions
in nonnegative integers are $(1,0,0)$, $(0,2,0)$, $(0,0,3)$, so $f_\lambda$ vanishes on exactly one orbit. As $\Delta$ has no zero on $\HH$, $j(z)=\lambda$ iff $f_\lambda(z)=0$; so $j$ takes
every value exactly once on $\SL_2(\Z)\backslash\HH$. By \Cref{thm:modforms:y0-riemann-surface}(ii) the invariant holomorphic function $j$ is holomorphic on $Y(1)$, and at $\infty$ it has a
simple pole in the chart $q$. So $j\colon X(1)\to\PP^1(\C)$ is a bijective holomorphic map of compact Riemann surfaces, hence of degree $1$ and an isomorphism (\Cref{ch:curves}).
Finally $j(z)=j(\Lambda_z)$ (\Cref{prop:modforms:delta-j}(ii)).
\end{proof}

\begin{paperlink}\label{pl:modforms:ch9}
\Cref{ch:ellcurves} proves that lattices up to homothety correspond to $\SL_2(\Z)\backslash\HH$ and to elliptic curves over $\C$, and left the bijectivity of $j$ and the
$q$-expansion of $j$ to this chapter. They are \Cref{thm:modforms:j-bijective} and \Cref{prop:modforms:delta-j}(ii); nothing here uses the uniformisation theorem, so
\Cref{ch:ellcurves} may use \Cref{thm:modforms:j-bijective} as a second proof that every elliptic curve over $\C$ is a complex torus. The paper uses $j$ only through
$v_q(j_E)$ (Definition def:terms, with the paper's $q$ a prime there): the algebraic $j$-invariant of \Cref{ch:ellcurves}, which equals $j(\tau)$ at any $\tau$ with
$E(\C)\simeq\C/\Lambda_\tau$. The Tate curve of \Cref{ch:reduction} is the $\mathsf q$-adic analogue of the expansion of $j$ in $q$.
\end{paperlink}

The bijectivity of $j$ also describes the function field of $X_0(N)$, which \Cref{ch:moddeg} uses to put $X_0(N)$ over $\Q$. Meromorphic functions on $X_0(N)$ are identified with
$\Gamma_0(N)$-invariant meromorphic functions on $\HH$ that are meromorphic in the charts $q_h$ at the cusps, and by Riemann existence and GAGA (\Cref{ch:curves}) they form the
function field $\C(X_0(N))$ of the curve. Write $\delta_N=\sm N001$, so that $j(Nz)=j(\delta_Nz)$.

\begin{proposition}[the function field of $X_0(N)$]\label{prop:modforms:function-field}
Let $N\ge1$, $j_N(z)=j(Nz)$, and $\SL_2(\Z)=\bigsqcup_{i=1}^{\mu(N)}\Gamma_0(N)\alpha_i$ with $\alpha_1=I$.
\begin{enumerate}[label=(\roman*)]
\item $j_N$ is a meromorphic function on $X_0(N)$, holomorphic on $Y_0(N)$.
\item Every meromorphic function on $X(1)$ is a rational function of $j$: $\C(X(1))=\C(j)$; through $\pi$ of \Cref{thm:modforms:genus-formula}, $\C(j)\subseteq\C(X_0(N))$.
\item The functions $f_i(z)=j(\delta_N\alpha_iz)$, $1\le i\le\mu(N)$, are pairwise distinct and are permuted by $z\mapsto\gamma z$, $\gamma\in\SL_2(\Z)$; the polynomial
$\Phi(X)=\prod_i(X-f_i)$ has coefficients in $\C(j)$ and is the minimal polynomial of $j_N=f_1$ over $\C(j)$. So $[\C(j,j_N):\C(j)]=\mu(N)$.
\item $\C(X_0(N))=\C(j,j_N)$.
\end{enumerate}
\end{proposition}

\begin{proof}
(i) For $\gamma=\sm abcd\in\Gamma_0(N)$, $\delta_N\gamma=\gamma'\delta_N$ with $\gamma'=\sm a{Nb}{c/N}d\in\SL_2(\Z)$, so $j_N(\gamma z)=j(\gamma'\delta_Nz)=j_N(z)$. So $j_N$ is a
$\Gamma_0(N)$-invariant holomorphic function on $\HH$, holomorphic on $Y_0(N)$ by \Cref{thm:modforms:y0-riemann-surface}(ii). At a cusp $\sigma\infty$ of width $h$, write $\delta_N\sigma=\sigma'\beta$ with
$\sigma'\in\SL_2(\Z)$ and $\beta=\sm ab0d$, $ad=N$ (as in the proof of \Cref{prop:modforms:mf-basic}(e)). Then $j_N(\sigma z)=j(\beta z)=\sum_{n\ge-1}c(n)e^{2\pi inb/d}e^{2\pi inaz/d}$, with $c(n)$ the
coefficients of $j$ (\Cref{prop:modforms:delta-j}); this is invariant under $z\mapsto z+h$ (as $\pm\sigma T^h\sigma^{-1}\in\Gamma_0(N)$), so only exponents in $\frac1h\Z$ occur, and it is a Laurent
series in $q_h$ with finitely many negative terms: $j_N$ is meromorphic at the cusp.
(ii) By \Cref{thm:modforms:j-bijective}, $j\colon X(1)\to\PP^1(\C)$ is an isomorphism, and the meromorphic functions on $\PP^1(\C)$ are the rational functions (\Cref{ch:curves}).
Composition with $\pi$ embeds $\C(X(1))$ in $\C(X_0(N))$.
(iii) Replacing $\alpha_i$ by $\gamma\alpha_i$, $\gamma\in\Gamma_0(N)$, does not change $f_i$, by (i). For $\gamma\in\SL_2(\Z)$, $\Gamma_0(N)\alpha_i\gamma=\Gamma_0(N)\alpha_{s(i)}$ for a permutation $s$
of the cosets, and $f_i(\gamma z)=f_{s(i)}(z)$; the action of $\SL_2(\Z)$ on $\Gamma_0(N)\backslash\SL_2(\Z)$ by right multiplication is transitive.
\emph{Distinctness.} Suppose $f_i=f_l$. For each $z\in\HH$, $j(\delta_N\alpha_iz)=j(\delta_N\alpha_lz)$, so by \Cref{thm:modforms:j-bijective} there is $g\in\SL_2(\Z)$ with
$\delta_N\alpha_iz=g\delta_N\alpha_lz$. For fixed $g$ the set of such $z$ is closed, and $\SL_2(\Z)$ is countable, so by Baire's theorem one of these sets has an interior point. Two
M\"obius transformations that agree on an open set are equal, and the matrices of $\GL_2^+(\R)$ inducing a given one differ by a real scalar; comparing determinants, $\delta_N\alpha_i=\pm g\delta_N\alpha_l$.
Then $\alpha_i\alpha_l^{-1}=\pm\delta_N^{-1}g\delta_N=\pm\sm{g_{11}}{g_{12}/N}{Ng_{21}}{g_{22}}$, which lies in $\SL_2(\Z)$ and has lower left entry divisible by $N$, so $\alpha_i\alpha_l^{-1}\in\Gamma_0(N)$
and $i=l$.
\emph{Coefficients.} The coefficients of $\Phi$ are the elementary symmetric functions of the $f_i$: holomorphic on $\HH$, invariant under $\SL_2(\Z)$ (which permutes the $f_i$), and
Laurent series in $e^{2\pi iz/N}$ with finitely many negative terms (by the computation in (i), $f_i(z)=j(\beta_iz)$ with $\beta_i=\sm{a_i}{b_i}0{d_i}$, $a_id_i=N$, and $a_i/d_i=a_i^2/N$), hence,
being invariant under $z\mapsto z+1$, Laurent series in $q$.
So they are meromorphic functions on $X(1)$, and lie in $\C(j)$ by (ii).
\emph{Irreducibility.} Let $P\in\C(j)[X]$ with $P(f_1)=0$. Composing with $\gamma\in\SL_2(\Z)$ fixes the coefficients of $P$, so $P(f_{s(1)})=0$; by transitivity every $f_i$ is a root, and
as the $f_i$ are distinct, $\deg P\ge\mu(N)$. Hence $\Phi$ is the minimal polynomial of $f_1=j_N$ and $[\C(j,j_N):\C(j)]=\mu(N)$.
(iv) The holomorphic map $j\circ\pi\colon X_0(N)\to\PP^1(\C)$ has degree $\deg\pi\cdot\deg j=\mu(N)\cdot1$ (\Cref{thm:modforms:genus-formula,thm:modforms:j-bijective}), and the degree of a
nonconstant function $F$ on a curve $C$, viewed as a map to $\PP^1$, is $[\C(C):\C(F)]$ (\Cref{ch:curves}). So $[\C(X_0(N)):\C(j)]=\mu(N)$, and $\C(j,j_N)\subseteq\C(X_0(N))$
has the same degree over $\C(j)$ by (iii); the two fields are equal.
\end{proof}

The polynomial $\Phi$, written as a polynomial $\Phi_N(j,X)$ in two variables, is the classical modular polynomial; for $N=2,3,5$ PARI's \texttt{polmodular} has degree $3,4,6=\mu(N)$ in
each variable. That $\Phi_N\in\Z[X,Y]$ (Kronecker), and the resulting model of $X_0(N)$ over $\Q$ with function field $\Q(j,j_N)$, are in \Cref{ch:moddeg}.

\section{Forms on \texorpdfstring{$\Gamma_0(N)$}{Gamma0(N)}: valence, differentials and dimensions}\label{sec:modforms:gamma0}

This section passes from level one to $\Gamma_0(N)$. Its central statement is that weight-two cusp forms are exactly the holomorphic differentials on $X_0(N)$,
so that $\dim S_2(\Gamma_0(N))$ is the genus. This is how the paper's newform $f$ becomes the differential $f(z)\,dz$ that the modular parametrisation pulls back
from $E$ (\Cref{ch:moddeg}), and how $S_2(\Gamma_0(N))\oplus\overline{S_2(\Gamma_0(N))}$ becomes $H^1(X_0(N),\C)$ in the proof of Proposition prop:pullback. We also state the valence
formula and the dimension formulas in every weight, and construct the weight-two Eisenstein series, which gives the Eisenstein congruence for \curve{11a1}.

\begin{theorem}[the valence formula for $\Gamma_0(N)$]\label{thm:modforms:valence-gamma0}
Let $f\in M_k(\Gamma_0(N))$, $f\ne0$. For $P\in X_0(N)$ let $v_P(f)$ be the order of $f$ at a representative $z\in\HH$ (in the variable $z$) if $P\in Y_0(N)$, and the order in $q_h$ of the
expansion of $f$ at the cusp $P$ if $P$ is a cusp; let $e_P\in\{1,2,3\}$ be the order of the stabiliser of $z$ in $\bar\Gamma_0(N)$, and $e_P=1$ at the cusps. Then
\[
\sum_{P\in X_0(N)}\frac{v_P(f)}{e_P}=\frac{k\,\mu(N)}{12}.
\]
In particular $v_\infty(f)\le k\mu(N)/12$, so $f$ is determined by $a_0(f),\dots,a_{B_k(N)}(f)$.
\end{theorem}

\begin{proof}[Sketch of proof]
Let $\SL_2(\Z)=\bigsqcup_{i=1}^{\mu}\Gamma_0(N)\gamma_i$ and $F=\prod_if|_k\gamma_i$. Replacing $\gamma_i$ by $\delta\gamma_i$ ($\delta\in\Gamma_0(N)$) does not change $f|_k\gamma_i$, and right
multiplication by $\gamma\in\SL_2(\Z)$ permutes the cosets, so $F|_{k\mu}\gamma=F$; $F$ is holomorphic at $\infty$ because each $f|_k\gamma_i$ is, so $F\in M_{k\mu}(\SL_2(\Z))$ and
\Cref{thm:modforms:valence-level-one} applies to it. Its order at a point $Q$ of $X(1)$ is the sum of the orders of the $f|_k\gamma_i$ there, and grouping the $\gamma_i$ according to
the points $P$ of $X_0(N)$ over $Q$, each $P$ contributes $v_P(f)$ times $e_Q/e_P$ (at an elliptic point) or $v_P(f)$ (at a cusp, where the order in $q$ of $f|_k\sigma$ is $v_P(f)/h$
and the width $h$ is the number of cosets over $P$). Comparing with the weights $\frac12$, $\frac13$, $1$ of the level-one formula gives the result. For $k=2$ one can instead use
\Cref{thm:modforms:weight-two-differentials} and $\deg\operatorname{div}(\omega)=2g-2$ (\Cref{ch:curves}) with the corrections at the cusps and elliptic points.
See \cite[Chapter 3, exercises of Section 3.1]{DiamondShurman}, \cite[Section 2.4]{Shimura71} and \cite[Chapter 9]{Stein} (locations not verified).
\end{proof}

For example $f_{11}=\eta(z)^2\eta(11z)^2$ (\Cref{thm:modforms:level-eleven}) has simple zeros at the cusps $\infty$ and $0$ and no other zeros: $1+1=2\cdot12/12$.

\begin{theorem}[weight-two forms are differentials]\label{thm:modforms:weight-two-differentials}
For $f\colon\HH\to\C$ holomorphic put $\omega_f=f(z)\,dz$.
\begin{enumerate}[label=(\roman*)]
\item $f\mapsto\omega_f$ is an isomorphism from $S_2(\Gamma_0(N))$ onto the space $H^0(X_0(N),\Omega^1)$ of holomorphic differentials on $X_0(N)$.
\item It maps $M_2(\Gamma_0(N))$ isomorphically onto the space of meromorphic differentials on $X_0(N)$ that are holomorphic on $Y_0(N)$ and have at most simple poles at the
cusps; the residue at the cusp $s=\sigma_s\infty$ of width $h_s$ is $\operatorname{res}_s\omega_f=\frac{h_s}{2\pi i}\,b_0(f|_2\sigma_s)$.
\item Hence $\dim S_2(\Gamma_0(N))=g(X_0(N))$, $\dim M_2(\Gamma_0(N))=g(X_0(N))+\nu_\infty(N)-1$ (which is $0$ for $N=1$), and
$\sum_sh_s\,b_0(f|_2\sigma_s)=0$ for every $f\in M_2(\Gamma_0(N))$, the sum over the cusps.
\end{enumerate}
\end{theorem}

\begin{proof}
A meromorphic differential on $X_0(N)$ restricts to a $\Gamma_0(N)$-invariant meromorphic differential $g(z)\,dz$ on $\HH$ minus the elliptic points, and by
\Cref{lem:modforms:slash}(iii), $\gamma^*(g\,dz)=g\,dz$ for all $\gamma\in\Gamma_0(N)$ iff $g|_2\gamma=g$. It remains to compare the local conditions.
\emph{Elliptic points.} Let $z_0$ have stabiliser of order $e$ in $\bar\Gamma_0(N)$, let $w$ be the Cayley coordinate at $z_0$, in which the stabiliser acts by $w\mapsto\zeta w$,
$\zeta=e^{2\pi i/e}$, and $t=w^e$ the chart (\Cref{thm:modforms:y0-riemann-surface}). An invariant differential $F(w)\,dw$ with $F$ holomorphic at $0$ satisfies $F(\zeta w)\zeta=F(w)$;
writing $F=\sum c_nw^n$ gives $c_n\zeta^{n+1}=c_n$, so only $n\equiv-1\pmod e$ occur and $F(w)=w^{e-1}G(w^e)$ with $G$ holomorphic. Then $F(w)\,dw=w^{e-1}G(t)\,dw=\frac1eG(t)\,dt$ is
holomorphic in $t$. So in weight $2$ there is no condition at the elliptic points: a holomorphic $f$ gives a differential holomorphic there.
\emph{Cusps.} At a cusp $s=\sigma\infty$ of width $h$, pull back by $\sigma$: $\sigma^*(\omega_f)=(f|_2\sigma)(z)\,dz$ and $dz=\frac h{2\pi i}\frac{dq_h}{q_h}$. If $f|_2\sigma=\sum_{n\ge n_0}b_nq_h^n$ then
$\omega_f=\frac h{2\pi i}\sum b_nq_h^{n-1}\,dq_h$ near $s$. This is holomorphic iff $b_n=0$ for $n\le0$, that is, iff $f$ vanishes at $s$, and has at most a simple pole iff $b_n=0$ for $n<0$,
that is, iff $f$ is holomorphic at $s$; the residue is $\frac h{2\pi i}b_0$.
Together: holomorphic differentials correspond exactly to cusp forms, and differentials holomorphic on $Y_0(N)$ with at most simple poles at the cusps to modular forms. This is
(i) and (ii). For (iii), $\dim H^0(X,\Omega^1)=g$ (\Cref{ch:curves}). Let $D$ be the sum of the $\nu_\infty\ge1$ cusps and $K$ a canonical divisor. By Riemann-Roch (\Cref{ch:curves}) the space of
meromorphic differentials $\omega$ with $\operatorname{div}\omega\ge-D$ has dimension $\deg(K+D)-g+1+\ell(-D)=(2g-2+\nu_\infty)-g+1+0=g+\nu_\infty-1$, as $\deg(-D)<0$. The relation between the
constant terms is the residue theorem $\sum_s\operatorname{res}_s\omega_f=0$ (\Cref{ch:curves}).
\end{proof}

For instance $S_2(\SL_2(\Z))=0$ and $g(X(1))=0$; $\dim S_2(\Gamma_0(11))=1$, $\dim S_2(\Gamma_0(37))=2$, $\dim S_2(\Gamma_0(121))=6$, $\dim S_2(\Gamma_0(539))=49$; PARI agrees.

\begin{theorem}[dimension formulas]\label{thm:modforms:dimension-formulas}
Let $g=g(X_0(N))$, $\nu_2=\nu_2(N)$, $\nu_3=\nu_3(N)$, $\nu_\infty=\nu_\infty(N)$. For even $k\ge4$,
\[
\dim S_k(\Gamma_0(N))=(k-1)(g-1)+\Bigl(\frac k2-1\Bigr)\nu_\infty+\Bigl\lfloor\frac k4\Bigr\rfloor\nu_2+\Bigl\lfloor\frac k3\Bigr\rfloor\nu_3,
\]
\[
\dim M_k(\Gamma_0(N))=\dim S_k(\Gamma_0(N))+\nu_\infty.
\]
For $k=2$: $\dim S_2=g$ and $\dim M_2=g+\nu_\infty-1$. Moreover $M_0(\Gamma_0(N))=\C$, $S_0(\Gamma_0(N))=0$, and $M_k(\Gamma_0(N))=0$ for $k<0$ or $k$ odd.
\end{theorem}

\begin{proof}
For $k=2$ this is \Cref{thm:modforms:weight-two-differentials}(iii). For $k$ odd use \Cref{prop:modforms:mf-basic}(b). For $k\le0$: a nonzero $f\in M_k$ has all $v_P\ge0$, so by
\Cref{thm:modforms:valence-gamma0} $k\ge0$, and for $k=0$ it has no zeros, so $f-f(z_0)$ vanishes identically. For even $k\ge4$ we give only a sketch; see
\cite[Theorems 3.5.1 and 3.6.1]{DiamondShurman} (numbers not verified) and \cite[Section 2.6]{Shimura71}. The map $f\mapsto f\,(dz)^{k/2}$ identifies $M_k(\Gamma_0(N))$ with the sections of a
line bundle on $X_0(N)$: the $\frac k2$-fold differentials whose orders are at least $-\frac k2$ at each cusp and at least $-\lfloor\frac k2(1-\frac1e)\rfloor$ at an elliptic point of order $e$
(the local computation of \Cref{thm:modforms:weight-two-differentials} in the $\frac k2$-th tensor power). This line bundle has degree $\frac k2(2g-2)+\frac k2\nu_\infty+\lfloor\frac k4\rfloor\nu_2+\lfloor\frac k3\rfloor\nu_3$,
which exceeds $2g-2$, so Riemann-Roch gives its dimension $\deg-g+1$, and similarly for cusp forms with $\frac k2-1$ in place of $\frac k2$ at the cusps.
\end{proof}

The formulas were checked against PARI's \texttt{mfdim([N,k],1)} and \texttt{mfdim([N,k],4)} for all $N\le60$ and $N\in\{77,99,121,539\}$, $k\in\{2,4,\dots,14\}$.
Two instances by hand. At level one ($g=0$, $\nu_2=\nu_3=\nu_\infty=1$) and $k=12$: $11\cdot(-1)+5+3+4=1$, the line $\C\Delta$ of \Cref{cor:modforms:level-one-structure}. At level $11$
($g=1$, $\nu_2=\nu_3=0$, $\nu_\infty=2$) and $k=4$: $3\cdot0+1\cdot2=2$, so $\dim S_4(\Gamma_0(11))=2$ and $\dim M_4(\Gamma_0(11))=4$. The floors $\lfloor\frac k4\rfloor$ and
$\lfloor\frac k3\rfloor$ are the elliptic corrections: a form of weight $k$ is forced to vanish at an elliptic point of order $e$ to an order that depends on $k$ modulo $2e$, as
the case $E_6(i)=0$ of \Cref{prop:modforms:delta-j} illustrates.

\begin{proposition}[the weight-two Eisenstein series of level $N$]\label{prop:modforms:eisenstein-weight-two}
For $N>1$ the function $E_2^{(N)}(z)=E_2(z)-NE_2(Nz)$ lies in $M_2(\Gamma_0(N))$, is not a cusp form, and
\[
E_2^{(N)}=(1-N)-24\sum_{n\ge1}\bigl(\sigma(n)-N\sigma(n/N)\bigr)q^n,\qquad\sigma=\sigma_1,\ \sigma(n/N)=0\text{ if }N\nmid n .
\]
For $N=p$ prime, $M_2(\Gamma_0(p))=\C E_2^{(p)}\oplus S_2(\Gamma_0(p))$.
\end{proposition}

\begin{proof}
With $E_2^*=E_2-\frac3{\pi y}$ (\Cref{prop:modforms:eisenstein-level-one}), $E_2^*(Nz)=E_2(Nz)-\frac3{\pi Ny}$, so $E_2^{(N)}=E_2^*-NE_2^*(Nz)$, and $NE_2^*(Nz)=(E_2^*|_2\alpha)(z)$ for
$\alpha=\sm N001$ (\Cref{def:modforms:slash}: $(\det\alpha)^{1}\cdot1^{-2}=N$). For $\gamma\in\Gamma_0(N)$, $\alpha\gamma\alpha^{-1}=\sm a{Nb}{c/N}d\in\SL_2(\Z)$, so $E_2^*|_2\alpha|_2\gamma=E_2^*|_2\alpha$ by the
invariance of $E_2^*$; hence $E_2^{(N)}|_2\gamma=E_2^{(N)}$. At a cusp, for $\sigma\in\SL_2(\Z)$ write $\alpha\sigma=\sigma'\beta$ with $\sigma'\in\SL_2(\Z)$ and $\beta=\sm{a'}{b'}0{d'}$, $a'd'=N$, as in
\Cref{prop:modforms:mf-basic}(e). Then $E_2^{(N)}|_2\sigma=E_2^*-E_2^*|_2\beta$, and
$(E_2^*|_2\beta)(z)=\frac{a'}{d'}E_2\bigl(\frac{a'z+b'}{d'}\bigr)-\frac{a'}{d'}\cdot\frac{3d'}{\pi a'y}$, so the non-holomorphic terms cancel and
$E_2^{(N)}|_2\sigma=E_2(z)-\frac{a'}{d'}E_2\bigl(\frac{a'z+b'}{d'}\bigr)$, a power series in $q^{1/d'}$ without negative powers, bounded on $\{\operatorname{Im}z\ge1\}$. This is the step where
$E_2^*$ earns its place: $E_2^*$ is modular but not holomorphic, $E_2$ is holomorphic but not modular, and the combination $E_2^*-E_2^*|_2\alpha$ is both, because the two
correction terms $\frac3{\pi y}$ and $N\cdot\frac3{\pi Ny}$ cancel. The expansion comes from substituting $Nz$ for $z$ in that of $E_2$:
$NE_2(Nz)=N-24N\sum_{m\ge1}\sigma(m)q^{Nm}$, and the term $q^{Nm}$ is $q^n$ with $m=n/N$. So $a_0=1-N\ne0$. For $N=p$, $\nu_\infty=2$, so $\dim M_2=g+1$ (\Cref{thm:modforms:weight-two-differentials}), and $E_2^{(p)}\notin S_2$.
\end{proof}

At the two cusps of $\Gamma_0(p)$ the constant terms are $1-p$ at $\infty$ (width $1$) and $1-\frac1p$ at $0$ (width $p$; here $\sigma=S$, $\beta=\sm100p$), and $1\cdot(1-p)+p(1-\frac1p)=0$,
as \Cref{thm:modforms:weight-two-differentials}(iii) requires.

\begin{example}[the Eisenstein congruence for \curve{11a1}]\label{ex:modforms:eisenstein-congruence}
Let $f_{11}$ be the normalised generator of $S_2(\Gamma_0(11))$ (\Cref{thm:modforms:level-eleven}). Then for every $n\ge1$,
\[
a_n(f_{11})\equiv\sigma(n)-11\sigma(n/11)\pmod5,
\]
and in particular $a_\ell(\curve{11a1})\equiv1+\ell\pmod5$ for every prime $\ell\ne11$, and $a_{11}=1$. Indeed
$h=12\bigl(f_{11}+\frac1{24}E_2^{(11)}\bigr)=-5+12\sum_{n\ge1}\bigl(a_n(f_{11})-\sigma(n)+11\sigma(n/11)\bigr)q^n$ lies in $M_2(\Gamma_0(11);\Z)$: the constant term is
$\frac{12}{24}(1-11)=-5$, and the $n$-th coefficient is $12a_n(f_{11})+\frac{12}{24}\cdot(-24)(\sigma(n)-11\sigma(n/11))$. The multiple $\frac1{24}E_2^{(11)}$ is chosen so that
the first coefficients cancel ($a_1=1-1=0$), and the factor $12$ makes the constant term $\frac{1-11}{24}=-\frac5{12}$ an integer. Its coefficients at $n=0,1,2$ are $-5$,
$12(1-1)=0$ and $12(-2-3)=-60$, all divisible by $5$. The Sturm bound is $B_2(11)=\lfloor12/6\rfloor=2$, so by \Cref{thm:modforms:sturm} every coefficient of $h$ is divisible by $5$, and
$12\bigl(a_n-\sigma(n)+11\sigma(n/11)\bigr)\equiv0$ gives the claim since $5\nmid12$. (Checked directly for $n\le5000$; the identification of $f_{11}$ with the form of \curve{11a1} at
every $n$ is \Cref{thm:modforms:level-eleven}(iv).) This congruence is the modular face of the rational point of order $5$ on \curve{11a1} (\Cref{ch:ellcurves}): if
$E(\Q)$ has a point of order $5$, then $5\mid\#\tilde E(\F_\ell)=\ell+1-a_\ell$ for every good $\ell\ne5$, because reduction is injective on torsion of order prime to $\ell$.
\end{example}

\begin{paperlink}\label{pl:modforms:eisenstein}
The paper's Conjecture conj:eis and Section sec:eis treat the Eisenstein case, where $\bar\rho$ is reducible and the newform $f$ is congruent modulo $\ell$ to an Eisenstein series.
\Cref{ex:modforms:eisenstein-congruence} is the simplest instance: for \curve{11a1} and $\ell=5$, $\bar\rho_{E,5}$ is reducible, $\deg\varphi_E=1$, the component group at $11$ (type
$\kod{I}_5$) has order $5$, and the deficit $D=1$ of Conjecture conj:eis is matched by $v_5(\#E(\Q)_{\mathrm{tors}})=1$. Open question 4 of the paper asks to identify the
deficit with an invariant of the Eisenstein ideal, for prime level the ideal of $\T$ generated by the $T_p-1-p$, $p\nmid N$, and $1+w_N$ (\cite{Mazur77}, \Cref{ch:moddeg}).
\end{paperlink}

\section{The Dedekind \texorpdfstring{$\eta$}{eta}-function and the level \texorpdfstring{$11$}{11} example}\label{sec:modforms:eta}

The running example of the paper is the curve \curve{11a1}, and its newform is the most explicit modular form there is: $\eta(z)^2\eta(11z)^2$. This section proves
that this product spans $S_2(\Gamma_0(11))$, computes its Atkin-Lehner sign, and lists its first coefficients next to the numbers $a_p$ of \curve{11a1}. The same
form reappears in the paper's Example of Section 3 (the congruences with $f_{\curve{11a}}\otimes\chi_{-11}$ and $f_{\curve{11a}}\otimes\chi_{-7}$), in the
Eisenstein congruence of \Cref{ex:modforms:eisenstein-congruence}, and in the Petersson computations of \Cref{ex:modforms:petersson-numerics}.

\begin{proposition}[the Dedekind $\eta$-function]\label{prop:modforms:eta}
The function $\eta(z)=q^{1/24}\prod_{n\ge1}(1-q^n)$, $q^{1/24}=e^{2\pi iz/24}$, is holomorphic and has no zeros on $\HH$, and
\[
\eta(z+1)=e^{\pi i/12}\eta(z),\qquad\eta(-1/z)=\sqrt{z/i}\,\eta(z),
\]
with the branch of the square root that is positive on the imaginary axis. Its logarithmic derivative is $\frac{\eta'}\eta=\frac{\pi i}{12}E_2$, and $\eta^{24}=\Delta$.
\end{proposition}

\begin{proof}
The product converges absolutely and locally uniformly since $\sum_n\abs q^n<\infty$, and no factor vanishes, so $\eta$ is holomorphic without zeros; the first law is
clear. The second we do not prove; see \cite[Chapter 3]{Apostol90}, \cite[Section III.2]{Koblitz} (the usual proof applies the residue theorem to a suitable
cotangent product, or uses the modularity of $E_2$ proved independently). The logarithmic derivative is
$\frac{2\pi i}{24}-\sum_{n\ge1}\frac{2\pi inq^n}{1-q^n}=\frac{2\pi i}{24}\bigl(1-24\sum_{n\ge1}\sum_{m\ge1}nq^{nm}\bigr)=\frac{\pi i}{12}E_2$. The identity $\eta^{24}=\Delta$ is
\Cref{prop:modforms:delta-j}(iii).
\end{proof}

\begin{theorem}[$S_2(\Gamma_0(11))$]\label{thm:modforms:level-eleven}
Let $f_{11}(z)=\eta(z)^2\eta(11z)^2=q\prod_{n\ge1}(1-q^n)^2(1-q^{11n})^2$.
\begin{enumerate}[label=(\roman*)]
\item $f_{11}\in S_2(\Gamma_0(11))$; it has simple zeros at the cusps $\infty$ and $0$ and no zeros on $\HH$.
\item $S_2(\Gamma_0(11))=\C f_{11}$, and $f_{11}$ is its normalised generator.
\item $f_{11}|_2W_{11}=-f_{11}$.
\item The coefficients of $f_{11}$ are the numbers $a_n$ of the curve \curve{11a1}: $y^2+y=x^3-x^2-10x-20$, that is, $a_p(f_{11})=p+1-\#\tilde E(\F_p)$ for $p\ne11$ and
$a_{11}(f_{11})=1$, and $\sum_na_n(f_{11})n^{-s}=L(\curve{11a1},s)$.
\end{enumerate}
\end{theorem}

\begin{proof}
(ii) and the zeros, for a generator $h$ of $S_2(\Gamma_0(11))$. By \Cref{thm:modforms:genus-formula} and \Cref{thm:modforms:weight-two-differentials}, $\dim S_2(\Gamma_0(11))=g(X_0(11))=1$. Let
$h\ne0$ span it. The holomorphic differential $\omega_h$ on the curve $X_0(11)$ of genus $1$ has a divisor of degree $2g-2=0$ (\Cref{ch:curves}) which is effective, so it is $0$.
Since $\nu_2(11)=\nu_3(11)=0$, the proof of \Cref{thm:modforms:weight-two-differentials} shows that the order of $\omega_h$ at a point of $Y_0(11)$ is the order of $h$, and at a
cusp it is the order of $h$ in $q_h$ minus $1$. So $h$ has no zeros on $\HH$, and simple zeros at the two cusps.

(i) Put $F=f_{11}^{12}=\eta(z)^{24}\eta(11z)^{24}=\Delta(z)\Delta(11z)$ (\Cref{prop:modforms:eta}). By \Cref{prop:modforms:mf-basic}(d), (e), $F\in S_{24}(\Gamma_0(11))$. Its order at $\infty$ is
$1+11=12$. At the cusp $0$: with $\alpha=\sm{11}001$, $\Delta(11z)=11^{-6}(\Delta|_{12}\alpha)(z)$ and $\alpha S=S\sm100{11}$, so $\Delta(11z)|_{12}S=11^{-6}\Delta|_{12}\sm100{11}=11^{-12}\Delta(z/11)$ and
$F|_{24}S=11^{-12}\Delta(z)\Delta(z/11)$, of order $11+1=12$ in $q_{11}=e^{2\pi iz/11}$. So the weight-$0$ function $\phi=h^{12}/F$ is $\Gamma_0(11)$-invariant, holomorphic on $\HH$ ($\eta$
has no zeros), and holomorphic and nonzero at both cusps (orders $12-12=0$). It is a holomorphic function on the compact connected Riemann surface $X_0(11)$, hence a nonzero
constant $c$. Then $(h/f_{11})^{12}=c$ on $\HH$, and the continuous function $h/f_{11}$ on the connected set $\HH$ takes values in the finite set of twelfth roots of $c$, so it is
constant. Hence $f_{11}$ is a nonzero multiple of $h$: it lies in $S_2(\Gamma_0(11))$ and has the zeros of $h$. As $a_1(f_{11})=1$, it is the normalised generator, which is (ii).

(iii) By \Cref{prop:modforms:eta}, $f_{11}(-1/(11z))=\eta(-1/(11z))^2\eta(-1/z)^2=\frac{11z}i\eta(11z)^2\cdot\frac zi\eta(z)^2=-11z^2f_{11}(z)$. With $W_{11}=\sm0{-1}{11}0$,
$(f_{11}|_2W_{11})(z)=11\cdot(11z)^{-2}f_{11}(-1/(11z))=-f_{11}(z)$.

(iv) We do not prove this; it is the Eichler-Shimura theorem for $X_0(11)$, which is itself an elliptic curve isomorphic to \curve{11a1} (\Cref{ch:moddeg},
\cite[Section 8.8]{DiamondShurman}, section number not verified). We checked $a_n(f_{11})=a_n(\curve{11a1})$ for all $n\le2000$ with PARI's \texttt{ellan}.
\end{proof}

Two remarks on (iii). The fixed point of $w_{11}$ is $z_0=i/\sqrt{11}$, and there the identity $f|_2W_{11}=\varepsilon_{11}(f)f$ reads $-f(z_0)=\varepsilon_{11}(f)f(z_0)$; so $\varepsilon_{11}(f)=+1$ would force $f(z_0)=0$.
Numerically $f_{11}(i/\sqrt{11})=0.10287776677\ldots\ne0$, which confirms $\varepsilon_{11}(f)=-1$ independently of the $\eta$-transformation law; PARI's \texttt{mfatkineigenvalues} also gives
$-1$. Second, the sign agrees with \Cref{ex:modforms:atkin-lehner-q-expansions}(c): $a_{11}(f_{11})=1=-\varepsilon_{11}(f)$.

The first coefficients are
\[
\begin{aligned}
f_{11}={}&q-2q^2-q^3+2q^4+q^5+2q^6-2q^7-2q^9-2q^{10}+q^{11}\\
&-2q^{12}+4q^{13}+4q^{14}-q^{15}-4q^{16}-2q^{17}+4q^{18}+2q^{20}+\cdots,
\end{aligned}
\]
and the $a_p$ for $p=2,3,5,7,11,13,17,19,23,29,31,37$ are $-2,-1,1,-2,1,4,-2,0,-1,0,7,3$. Reducing $y^2+y=x^3-x^2-10x-20$ modulo $p$ and counting points gives $p+1-a_p$
(\Cref{ch:ellcurves}); for instance $\#\tilde E(\F_2)=5$ and $\#\tilde E(\F_{13})=10$. The rational point of order $5$ of \curve{11a1} shows up as $5\mid p+1-a_p$, the congruence of
\Cref{ex:modforms:eisenstein-congruence}. Physically $f_{11}$ is a Fourier series on the half cylinder (\Cref{sec:modforms:forms}) whose modes are dictated by the arithmetic of
a cubic curve.

The computation in the proof of (i) also gives the example promised after \Cref{def:modforms:modular-form}, of why holomorphy must be checked at every cusp. Let
$g(z)=\Delta(11z)^2/\Delta(z)$. It is holomorphic on $\HH$ ($\Delta$ has no zeros there), and $g|_{12}\gamma=g$ for $\gamma\in\Gamma_0(11)$, since $\Delta(11z)^2$ has weight $24$, $\Delta$ has
weight $12$, and the slash operator is multiplicative in the sense $(F/G)|_{a-b}\gamma=(F|_a\gamma)/(G|_b\gamma)$. At $\infty$ its expansion is $q^{22-1}(1+\cdots)$, so $g$ is bounded on
$\{\operatorname{Im}z\ge1\}$ and even vanishes to order $21$. But at the cusp $0$, by the formula $\Delta(11z)|_{12}S=11^{-12}\Delta(z/11)$ of the proof,
$g|_{12}S=11^{-24}\Delta(z/11)^2/\Delta(z)$, and in the local parameter $q_{11}=e^{2\pi iz/11}$ the numerator has order $2$ and the denominator $\Delta(z)=q_{11}^{11}+\cdots$ order $11$: $g$ has a pole
of order $9$ at $0$, and $g\notin M_{12}(\Gamma_0(11))$. The valence count is consistent: $21-9=12=\frac{12\cdot12}{12}$ (\Cref{thm:modforms:valence-gamma0}; its sketch of proof applies to meromorphic forms, since \Cref{thm:modforms:valence-level-one} does).

\begin{example}[small levels]\label{ex:modforms:small-levels}
\begin{enumerate}[label=(\alph*)]
\item For the twelve levels $N=11$, $14$, $15$, $17$, $19$, $20$, $21$, $24$, $27$, $32$, $36$, $49$ of genus $1$, $S_2(\Gamma_0(N))$ is one-dimensional, spanned by the newform of the optimal curve of conductor $N$ (\Cref{ch:moddeg}). Eight of these are $\eta$-products, each checked against Cremona's curve for $n\le300$:
\begin{center}
\begin{tabular}{@{}llll@{}}
\toprule
$N$ & form & $N$ & form\\
\midrule
$11$ & $\eta(z)^2\eta(11z)^2$ & $24$ & $\eta(2z)\eta(4z)\eta(6z)\eta(12z)$\\
$14$ & $\eta(z)\eta(2z)\eta(7z)\eta(14z)$ & $27$ & $\eta(3z)^2\eta(9z)^2$\\
$15$ & $\eta(z)\eta(3z)\eta(5z)\eta(15z)$ & $32$ & $\eta(4z)^2\eta(8z)^2$\\
$20$ & $\eta(2z)^2\eta(10z)^2$ & $36$ & $\eta(6z)^4$\\
\bottomrule
\end{tabular}
\end{center}
(curves \curve{11a1}, \curve{14a1}, \curve{15a1}, \curve{20a1}, \curve{24a1}, \curve{27a1}, \curve{32a1}, \curve{36a1}). That each product is a modular form for $\Gamma_0(N)$ follows from
Ligozat's criterion for $\eta$-quotients \cite{Ligozat75} (stated, not proved here), or, as in \Cref{thm:modforms:level-eleven}, from the orders of its twelfth power at the cusps
(Exercise \ref{ex:modforms:ex-eta20}).
\item $\dim S_2(\Gamma_0(37))=2$, spanned by the newforms of \curve{37a1} ($q-2q^2-3q^3+2q^4-2q^5+\cdots$) and \curve{37b1} ($q+q^3-2q^4+\cdots$), with $w_{37}$-eigenvalues $+1$
and $-1$ (PARI); \curve{37a1} has rank $1$ and root number $-1$, \curve{37b1} rank $0$ and root number $+1$.
\item $\dim S_2(\Gamma_0(121))=6$: four newforms (the classes \curve{121a} to \curve{121d}) and the two forms $f_{11}(z)$, $f_{11}(11z)$ coming from level $11$ (\Cref{ch:hecke}).
\end{enumerate}
\end{example}

\section{The Petersson inner product}\label{sec:modforms:petersson}

The Petersson product is the inner product on cusp forms. It enters the paper in two ways. Qualitatively, through orthogonality: the congruence number $r_E$ is
defined with forms ``orthogonal to $f$'', Lemma lem:crit uses ``the orthogonal complement $f^\perp$'', and Lemma lem:old uses forms ``orthogonal to the newform $f$''.
Quantitatively, through the Petersson norm: ``Up to explicit factors $\degphi$ is $\langle f,f\rangle$ divided by the area of the period lattice of $E$'' (the paper's
Introduction), made exact in Remark rem:degree. This section defines the product, fixes its normalisation once and for all, proves its basic properties, and checks the
normalisation numerically against the modular degrees of five curves.

\begin{definition}[the Petersson product]\label{def:modforms:petersson}
For $f\in S_k(\Gamma_0(N))$ and $g\in M_k(\Gamma_0(N))$,
\[
\langle f,g\rangle=\langle f,g\rangle_{\Gamma_0(N)}=\int_{\Gamma_0(N)\backslash\HH}f(z)\overline{g(z)}\,y^k\,\frac{dx\,dy}{y^2},
\]
the integral over any fundamental domain, for instance $\bigcup_i\gamma_i\mathcal F$ (\Cref{prop:modforms:coset-domains}). For a subgroup $\Gamma\subseteq\Gamma_0(N)$ of finite index the
same integral over $\Gamma\backslash\HH$ is written $\langle f,g\rangle_\Gamma$.
\end{definition}

\begin{proposition}[basic properties]\label{prop:modforms:petersson}
\begin{enumerate}[label=(\roman*)]
\item The integrand $\phi(z)=f(z)\overline{g(z)}\,y^k$ is invariant under $\Gamma_0(N)$, so the integral does not depend on the fundamental domain.
\item The integral converges absolutely, and $\langle f,g\rangle=\sum_i\int_{\mathcal F}(f|_k\gamma_i)\overline{(g|_k\gamma_i)}\,y^k\,d\mu$.
\item On $S_k(\Gamma_0(N))$, $\langle\ ,\ \rangle$ is a positive definite Hermitian inner product, linear in the first variable.
\item In weight $2$, $\langle f,g\rangle=\iint_{\Gamma_0(N)\backslash\HH}f\bar g\,dx\,dy=\frac i2\int_{X_0(N)}\omega_f\wedge\overline{\omega_g}$ for $f,g\in S_2(\Gamma_0(N))$: the Petersson product is the Hodge inner
product of the holomorphic differentials of \Cref{thm:modforms:weight-two-differentials}.
\end{enumerate}
\end{proposition}

\begin{proof}
(i) $\phi(\gamma z)=\operatorname{Im}(\gamma z)^kf(\gamma z)\overline{g(\gamma z)}=y^k(f|_k\gamma)(z)\overline{(g|_k\gamma)(z)}$ for $\gamma\in\SL_2(\Z)$, by the computation of \Cref{lem:modforms:slash}(iv)
without absolute values; for $\gamma\in\Gamma_0(N)$ this is $\phi(z)$. With the invariance of $d\mu$, two fundamental domains, which differ by cutting into pieces and moving the pieces
by elements of $\Gamma_0(N)$, give the same integral.
(ii) The same identity with $\gamma=\gamma_i$ and the change of variables $z=\gamma_iw$ give the formula. On $\mathcal F$, $\operatorname{Im}w\ge\sqrt3/2$, $\abs{(f|_k\gamma_i)(w)}\le C_ie^{-2\pi\operatorname{Im}w/h_i}$
(no constant term, as in \Cref{prop:modforms:hecke-bound}) and $\abs{(g|_k\gamma_i)(w)}\le C_i'$ (holomorphy at the cusp); and $\int_{\mathcal F}e^{-2\pi y/h}y^{k-2}\,dx\,dy<\infty$.
(iii) Sesquilinearity and $\langle g,f\rangle=\overline{\langle f,g\rangle}$ are clear; $\langle f,f\rangle=\int\abs f^2y^k\,d\mu>0$ unless $f$ vanishes on an open set, hence everywhere.
(iv) $y^2\cdot y^{-2}=1$, and $\frac i2\,dz\wedge d\bar z=dx\wedge dy$, because $dz\wedge d\bar z=(dx+i\,dy)\wedge(dx-i\,dy)=-2i\,dx\wedge dy$; also
$\omega_f\wedge\overline{\omega_g}=f\bar g\,dz\wedge d\bar z$. The integral over $X_0(N)$ is the integral over a fundamental domain, the cusps and elliptic points having measure
zero (\Cref{ch:curves}).
\end{proof}

The Petersson product is the weighted $L^2$ inner product $\int f\bar g\,w\,d\mu$ of functions on a surface of finite area, with the weight $w=y^k$, just as Sturm-Liouville
theory uses weighted inner products. \Cref{ch:hecke} shows that the Hecke operators $T_n$, $(n,N)=1$, are self-adjoint for it, so that their eigenforms are orthogonal, as
the eigenfunctions of a self-adjoint operator are.

\begin{convention}[normalisations of the Petersson product]\label{conv:modforms:petersson-normalisations}
\leavevmode
\begin{enumerate}[label=(\alph*)]
\item The book's product is the \emph{unnormalised} one of \Cref{def:modforms:petersson} (\Cref{ch:notation}, B-C8).
\item PARI's \texttt{mfpetersson} returns the book's product divided by $\mu(N)$.
\item The ``normalised'' product of some books is the book's divided by $\vol(\Gamma_0(N)\backslash\HH)=\pi\mu(N)/3$; it does not depend on the level at which $f$ and $g$ are viewed
(\Cref{prop:modforms:petersson-change}).
\item The formula $\degphi\cdot\vol(\Lambda_E)=4\pi^2c_E^2\langle f,f\rangle$ (\Cref{ch:moddeg}, the paper's Remark rem:degree) holds with the book's product and with the Manin constant
$c_E>0$ defined by $\varphi_E^*\omega_E=c_E\cdot2\pi i\,f(z)\,dz$ (\Cref{ch:curves,ch:moddeg}).
\end{enumerate}
A numerical value of a Petersson norm is never quoted without its normalisation.
\end{convention}

\begin{proposition}[change of group and invariance]\label{prop:modforms:petersson-change}
\begin{enumerate}[label=(\roman*)]
\item If $\Gamma\subseteq\Gamma_0(N)$ has finite index, then $\langle f,g\rangle_\Gamma=[\bar\Gamma_0(N):\bar\Gamma]\,\langle f,g\rangle_{\Gamma_0(N)}$. In particular, for $f\in S_k(\Gamma_0(N))$, $g\in M_k(\Gamma_0(N))$
viewed at level $NM$, $\langle f,g\rangle_{\Gamma_0(NM)}=\frac{\mu(NM)}{\mu(N)}\langle f,g\rangle_{\Gamma_0(N)}$.
\item If $\alpha\in\GL_2^+(\Q)$ normalises $\Gamma_0(N)$, then $\langle f|_k\alpha,g|_k\alpha\rangle=\langle f,g\rangle$ for $f\in S_k(\Gamma_0(N))$, $g\in M_k(\Gamma_0(N))$, provided $f|_k\alpha$ and $g|_k\alpha$
again lie in these spaces.
\item (Groups commensurable with $\SL_2(\Z)$.) Let $\Gamma\subset\SL_2(\Q)$ be a subgroup that contains some $\Gamma(M)$ as a subgroup of finite index, and write $\bar\Gamma$ for its image in
$\PSL_2(\R)$. Define $M_k(\Gamma)\supseteq S_k(\Gamma)$ by the conditions of \Cref{def:modforms:modular-form} (invariance under $\Gamma$, behaviour under every $\sigma\in\SL_2(\Z)$). For a congruence
subgroup $\Gamma_1\subseteq\SL_2(\Z)$ write $\langle f,g\rangle_{\Gamma_1}=\int_{\mathcal F_{\Gamma_1}}f\bar gy^k\,d\mu$, with $\mathcal F_{\Gamma_1}$ the domain of \Cref{prop:modforms:coset-domains} for
$\pm\Gamma_1$. For $f\in S_k(\Gamma)$, $g\in M_k(\Gamma)$ the number
\[
\langle f,g\rangle_\Gamma=[\bar\Gamma:\bar\Gamma_1]^{-1}\langle f,g\rangle_{\Gamma_1},
\]
for $\Gamma_1\subseteq\Gamma\cap\SL_2(\Z)$ a congruence subgroup of finite index in $\Gamma$, does not depend on $\Gamma_1$, and for $\Gamma\subseteq\Gamma_0(N)$ of finite index it is the product of \Cref{def:modforms:petersson}.
\item (Change of variables.) In the situation of (iii) let $\alpha\in\GL_2^+(\Q)$. Then $\alpha^{-1}\Gamma\alpha$ is again a group as in (iii), $f\mapsto f|_k\alpha$ maps $M_k(\Gamma)$ onto
$M_k(\alpha^{-1}\Gamma\alpha)$ and $S_k(\Gamma)$ onto $S_k(\alpha^{-1}\Gamma\alpha)$, and
\[
\langle f|_k\alpha,g|_k\alpha\rangle_{\alpha^{-1}\Gamma\alpha}=\langle f,g\rangle_\Gamma\qquad(f\in S_k(\Gamma),\ g\in M_k(\Gamma)).
\]
\end{enumerate}
\end{proposition}

\begin{proof}
(i) A fundamental domain of $\Gamma$ is the union of $[\bar\Gamma_0(N):\bar\Gamma]$ translates of one of $\Gamma_0(N)$ (the proof of \Cref{prop:modforms:coset-domains}, with $\Gamma_0(N)$ in place of
$\SL_2(\Z)$), and the integrand is $\Gamma_0(N)$-invariant. Both $\Gamma_0(NM)$ and $\Gamma_0(N)$ contain $-I$, so the index is $\mu(NM)/\mu(N)$.
(ii) Let $D$ be a fundamental domain. The identity of (i) of \Cref{prop:modforms:petersson} holds for $\alpha\in\GL_2^+(\R)$ (the factors $\det\alpha$ cancel against
$\operatorname{Im}(\alpha z)=\det\alpha\,y/\abs{cz+d}^2$), so $\int_D(f|_k\alpha)\overline{(g|_k\alpha)}y^k\,d\mu=\int_D\phi(\alpha z)\,d\mu(z)=\int_{\alpha D}\phi\,d\mu$, and $\alpha D$ is a fundamental domain for
$\alpha\Gamma_0(N)\alpha^{-1}=\Gamma_0(N)$.

For (iii) and (iv) we use one measure-theoretic fact. Call a measurable set $D\subseteq\HH$ a \emph{fundamental set} for a congruence subgroup $\Gamma_1$ if the translates $\gamma D$,
$\gamma\in\bar\Gamma_1$, cover $\HH$ and any two distinct ones meet in a null set; $\mathcal F_{\Gamma_1}$ is one, by \Cref{prop:modforms:coset-domains} (translates of $\mathcal F$ meet only in
boundary arcs). If $D,D'$ are fundamental sets for $\Gamma_1$ and $\phi\ge0$ is measurable and $\Gamma_1$-invariant, then, by the invariance of $\phi\,d\mu$,
\[
\int_D\phi\,d\mu=\sum_{\gamma\in\bar\Gamma_1}\int_{D\cap\gamma D'}\phi\,d\mu=\sum_{\gamma\in\bar\Gamma_1}\int_{\gamma^{-1}D\cap D'}\phi\,d\mu=\int_{D'}\phi\,d\mu ;
\]
applied to $\abs\phi$ and then to the positive and negative parts of $\operatorname{Re}\phi$ and $\operatorname{Im}\phi$, the same holds for every $\Gamma_1$-invariant $\phi$ integrable over $D'$.
(iii) Since $f\bar g\,y^k$ is invariant under $-I$, it is invariant under $\pm\Gamma_1$, and the integral over $\mathcal F_{\Gamma_1}$ converges by the argument of \Cref{prop:modforms:petersson}(ii) (the
expansions of $f|_k\sigma$ in $q_h$ at every cusp have no constant term). If $\Gamma_2\subseteq\Gamma_1$ are two congruence subgroups as in (iii) and $\delta_1,\dots,\delta_r$ represent
$\bar\Gamma_2\backslash\bar\Gamma_1$, then $\bigcup_j\delta_j\mathcal F_{\Gamma_1}$ is a fundamental set for $\Gamma_2$, so $\langle f,g\rangle_{\Gamma_2}=r\langle f,g\rangle_{\Gamma_1}$, with
$r=[\bar\Gamma_1:\bar\Gamma_2]$; as indices multiply, $[\bar\Gamma:\bar\Gamma_2]^{-1}\langle f,g\rangle_{\Gamma_2}=[\bar\Gamma:\bar\Gamma_1]^{-1}\langle f,g\rangle_{\Gamma_1}$. Two choices $\Gamma_1$, $\Gamma_1'$ are compared through
$\Gamma_1\cap\Gamma_1'$, which is again a congruence subgroup of finite index in $\Gamma$. For $\Gamma\subseteq\Gamma_0(N)$ take $\Gamma_1=\Gamma$ (and $\Gamma_0(N)$ itself in \Cref{def:modforms:petersson}).
(iv) Replacing $\alpha$ by a positive multiple changes neither $\alpha^{-1}\Gamma\alpha$ nor $f|_k\alpha$ (the factors $\lambda^k$ and $\lambda^{-k}$ cancel), so we may assume $\alpha\in M_2(\Z)$; put
$m=\det\alpha$ and $\alpha'=m\alpha^{-1}\in M_2(\Z)$. For every $L\ge1$ and $X\in M_2(\Z)$,
\[
\alpha(1+mLX)\alpha^{-1}=1+L\,\alpha X\alpha',\qquad\alpha^{-1}(1+mLX)\alpha=1+L\,\alpha'X\alpha,
\]
so $\alpha\Gamma(mL)\alpha^{-1}\subseteq\Gamma(L)$ and $\alpha^{-1}\Gamma(mL)\alpha\subseteq\Gamma(L)$. Hence $\alpha^{-1}\Gamma\alpha\supseteq\alpha^{-1}\Gamma(M)\alpha\supseteq\Gamma(mM)$, and
$[\alpha^{-1}\Gamma\alpha:\Gamma(mM)]\le[\alpha^{-1}\Gamma\alpha:\alpha^{-1}\Gamma(m^2M)\alpha]=[\Gamma:\Gamma(m^2M)]<\infty$ because $\alpha^{-1}\Gamma(m^2M)\alpha\subseteq\Gamma(mM)$. So $\alpha^{-1}\Gamma\alpha$ is as in (iii).
For $\gamma\in\Gamma$, $(f|_k\alpha)|_k(\alpha^{-1}\gamma\alpha)=f|_k\gamma|_k\alpha=f|_k\alpha$ by \Cref{lem:modforms:slash}. For $\sigma\in\SL_2(\Z)$ write $\alpha\sigma=\sigma'\beta$ with $\sigma'\in\SL_2(\Z)$ and
$\beta$ upper triangular with integer entries, as in the proof of \Cref{prop:modforms:mf-basic}(e); then $(f|_k\alpha)|_k\sigma=(f|_k\sigma')|_k\beta$ is bounded on $\{\operatorname{Im}z\ge1\}$,
and tends to $0$ there if $f$ is a cusp form, by the argument given there. So $f|_k\alpha\in M_k(\alpha^{-1}\Gamma\alpha)$, in $S_k$ if $f\in S_k(\Gamma)$; $\alpha^{-1}$ gives the inverse map.
For the product take $\Gamma_1=\Gamma(m^2M)\subseteq\Gamma$. Then $\Delta=\alpha^{-1}\Gamma_1\alpha$ satisfies $\Gamma(m^3M)\subseteq\Delta\subseteq\Gamma(mM)$ (the first inclusion because
$\alpha\Gamma(m^3M)\alpha^{-1}\subseteq\Gamma(m^2M)$), so it is a congruence subgroup of finite index in $\alpha^{-1}\Gamma\alpha$, and conjugation by $\alpha$ gives
$[\overline{\alpha^{-1}\Gamma\alpha}:\bar\Delta]=[\bar\Gamma:\bar\Gamma_1]$. It therefore suffices to show $\langle f|_k\alpha,g|_k\alpha\rangle_\Delta=\langle f,g\rangle_{\Gamma_1}$. With $\phi=f\bar g\,y^k$, the
identity $\phi(\alpha z)=(f|_k\alpha)(z)\overline{(g|_k\alpha)(z)}\,y^k$ (as in (ii)) and the invariance of $d\mu$ give
\[
\langle f|_k\alpha,g|_k\alpha\rangle_\Delta=\int_{\mathcal F_\Delta}\phi(\alpha z)\,d\mu(z)=\int_{\alpha\mathcal F_\Delta}\phi\,d\mu .
\]
The set $\alpha\mathcal F_\Delta$ is a fundamental set for $\Gamma_1=\alpha\Delta\alpha^{-1}$: its translates are $\gamma\alpha\mathcal F_\Delta=\alpha(\alpha^{-1}\gamma\alpha)\mathcal F_\Delta$, and the diffeomorphism
$z\mapsto\alpha z$ preserves null sets. By the fact above (first for $\abs\phi$, which shows integrability), the integral equals $\int_{\mathcal F_{\Gamma_1}}\phi\,d\mu=\langle f,g\rangle_{\Gamma_1}$.
\end{proof}

\Cref{ch:hecke} uses (iv) in two ways. For the adjoints of double coset operators, $\Gamma=\Gamma_0(N)\cap\alpha\Gamma_0(N)\alpha^{-1}$ for a matrix $\alpha$ of a
double coset; then $\Gamma$ and $\alpha^{-1}\Gamma\alpha=\alpha^{-1}\Gamma_0(N)\alpha\cap\Gamma_0(N)$ are both congruence subgroups of $\SL_2(\Z)$, and (iv) is needed
only because $\alpha$ normalises neither. For the norms of oldforms, $\Gamma=\Gamma_0(M)$ and $\alpha=\sm p001$; then $\alpha^{-1}\Gamma\alpha$ does not lie in
$\SL_2(\Z)$, and (iii) is needed to define the product over it.

The factor in (i) is why the normalisation matters as soon as levels change: the paper passes from level $M$ to $M^\Sigma$ in Section 4.1, and Theorem thm:quant works at level
$N$ with forms of level $N_S$. The normalised product of \Cref{conv:modforms:petersson-normalisations}(c) is level-independent; congruence questions do not see the factor at all,
since they only use orthogonality.

\begin{proposition}[orthogonal complements]\label{prop:modforms:orthogonal-complement}
For $0\ne f\in S_k(\Gamma_0(N))$ let $f^\perp=\{g\in S_k(\Gamma_0(N)):\langle g,f\rangle=0\}$. Then $S_k(\Gamma_0(N))=\C f\oplus f^\perp$, and the projection onto $\C f$ along $f^\perp$ is
$g\mapsto\frac{\langle g,f\rangle}{\langle f,f\rangle}f$.
\end{proposition}

\begin{proof}
$\langle f,f\rangle>0$, so $g-\frac{\langle g,f\rangle}{\langle f,f\rangle}f\in f^\perp$: its product with $f$ is $\langle g,f\rangle-\frac{\langle g,f\rangle}{\langle f,f\rangle}\langle f,f\rangle=0$. So
$g=\frac{\langle g,f\rangle}{\langle f,f\rangle}f+\bigl(g-\frac{\langle g,f\rangle}{\langle f,f\rangle}f\bigr)$ is a decomposition, and it is unique because $\C f\cap f^\perp=0$ (if $cf\in f^\perp$ then
$c\langle f,f\rangle=0$, so $c=0$).
\end{proof}

Equivalently, $f^\perp$ is the kernel of the nonzero linear functional $g\mapsto\langle g,f\rangle$, so it has codimension $1$. For $N=22$, with $f=f_{11}$ and $S_2(\Gamma_0(22))$
spanned by $f_{11}(z)$ and $f_{11}(2z)$ (\Cref{ex:modforms:atkin-lehner-q-expansions}(b)), the line $f^\perp$ is spanned by $f_{11}(2z)-\frac{\langle f_{11}(2z),f_{11}\rangle_{\Gamma_0(22)}}{\langle f_{11},f_{11}\rangle_{\Gamma_0(22)}}f_{11}(z)$; with the
values of Exercise \ref{ex:modforms:ex-level22} (where $\langle f_{11}(z),f_{11}(2z)\rangle_{\Gamma_0(22)}=-\langle f_{11},f_{11}\rangle_{\Gamma_0(11)}$ and $\langle f_{11},f_{11}\rangle_{\Gamma_0(22)}=3\langle f_{11},f_{11}\rangle_{\Gamma_0(11)}$),
the ratio is $-\frac13$ and $f^\perp=\C\bigl(f_{11}(2z)+\frac13f_{11}(z)\bigr)$. Its $q$-expansion has non-integral coefficients, and $3f_{11}(2z)+f_{11}(z)$ spans $f^\perp\cap S_2(\Gamma_0(22);\Z)$:
this is the kind of lattice on which the paper's congruence questions are posed (here $f_{11}$ is not a newform of level $22$, so the example only illustrates the definition).

By itself $f^\perp$ is defined transcendentally. \Cref{ch:hecke} shows, for a newform $f$ with rational coefficients, that $f^\perp$ is stable under all Hecke operators,
$U_p$ included (the adjoint of $U_p$ is $w_NU_pw_N$, and $f$ is an eigenvector of $w_N$), that it is spanned by forms with rational coefficients, and that
$f^\perp\cap S_2(\Gamma_0(N);\Z)$ is a saturated sublattice of rank $\dim S_2(\Gamma_0(N))-1$. None of this is proved here.

\begin{example}[Petersson norms and modular degrees]\label{ex:modforms:petersson-numerics}
For five optimal curves of the paper's world, with $f$ the newform, $\Lambda_E$ the period lattice of the N\'eron differential of the minimal model (\Cref{ch:ellcurves}) and
$\degphi$ the modular degree (Cremona's tables, \Cref{ch:moddeg}), PARI gives (unnormalised product, $\langle f,f\rangle=\mu(N)\cdot\texttt{mfpetersson}$):
\begin{center}
\begin{tabular}{@{}lrrllcc@{}}
\toprule
curve & $N$ & $\mu(N)$ & $\langle f,f\rangle$ & $\vol(\Lambda_E)$ & $4\pi^2\langle f,f\rangle/\vol(\Lambda_E)$ & $\degphi$\\
\midrule
\curve{11a1} & $11$ & $12$ & $0.0469001478734952$ & $1.85154362345596$ & $1.0000000$ & $1$\\
\curve{20a1} & $20$ & $36$ & $0.0813494568253817$ & $3.21154782844007$ & $1.0000000$ & $1$\\
\curve{37a1} & $37$ & $38$ & $0.3717541475106961$ & $7.33813274078958$ & $2.0000000$ & $2$\\
\curve{99a1} & $99$ & $144$ & $0.5385195687108607$ & $5.31497510542145$ & $4.0000000$ & $4$\\
\curve{121a1} & $121$ & $132$ & $0.4282393600119572$ & $2.81770204819580$ & $6.0000000$ & $6$\\
\bottomrule
\end{tabular}
\end{center}
(the ratio agrees with $\degphi$ to $40$ digits; for \curve{11a1}, \texttt{mfpetersson} returns $0.0039083456561246=0.0469001478734952/12$). So
$\degphi\cdot\vol(\Lambda_E)=4\pi^2\langle f,f\rangle$ for these curves, whose Manin constant is $1$. The derivation is one line, given \Cref{ch:moddeg}: if $\varphi_E^*\omega_E=c_E\cdot2\pi i\,f(z)\,dz$,
then pulling back the area form,
\[
\degphi\cdot\frac i2\int_E\omega_E\wedge\bar\omega_E=\frac i2\int_{X_0(N)}\varphi_E^*\omega_E\wedge\overline{\varphi_E^*\omega_E}=4\pi^2c_E^2\cdot\frac i2\int_{X_0(N)}\omega_f\wedge\overline{\omega_f},
\]
that is $\degphi\cdot\vol(\Lambda_E)=4\pi^2c_E^2\langle f,f\rangle$ by \Cref{prop:modforms:petersson}(iv) and \Cref{ch:curves} (degree of a pullback of a $2$-form). The Manin constant stands in the
numerator: for the degree-$5$ parametrisation $X_0(11)\to\curve{11a3}$ one finds $\vol(\Lambda_{\curve{11a3}})=5\vol(\Lambda_{\curve{11a1}})$, so $5\cdot5\vol(\Lambda_{\curve{11a1}})=4\pi^2c^2\langle f,f\rangle$ forces the
constant of that parametrisation to be $5$ (Exercise \ref{ex:modforms:ex-petersson-pari}).
\end{example}

\begin{remark}[what other chapters add]\label{rem:modforms:petersson-forward}
\begin{enumerate}[label=(\alph*)]
\item (\Cref{ch:hecke}.) For $(n,N)=1$ the Hecke operator $T_n$ is self-adjoint for $\langle\ ,\ \rangle$; for $p\mid N$, $U_p$ is not, its adjoint being $w_NU_pw_N$; newforms with different
systems of eigenvalues are orthogonal, and oldforms are orthogonal to newforms.
\item (\Cref{ch:lfunctions}.) Rankin-Selberg: for a newform $f$ of weight $2$ and level $N$, $\langle f,f\rangle=N\,L^{\nv}(\Sym f,2)/(8\pi^3)$ with the naive symmetric square
$L$-function (\cite[Lemma 2.12]{DFG04}, \cite{Shimura76}, \cite{Rankin39}); we checked this to $40$ digits for \curve{11a1} and \curve{37a1} with PARI's \texttt{lfunsympow} (at these squarefree levels the naive and the true $L$-function agree, so the check does not see the difference).
At \curve{121a1}, of additive reduction at $11$ with $a_{11}=0$, it does: $\langle f,f\rangle=0.4282393600\ldots=121\,L^{\nv}(\Sym f,2)/(8\pi^3)$, while the true $L$-function gives
$121\,L(\Sym f,2)/(8\pi^3)=0.3925527466\ldots$; the ratio is $\frac{12}{11}$ to $40$ digits, the inverse of the Euler factor $(1+11^{1-s})^{-1}$ at $s=2$ that the true $L$-function has at $11$ and
the naive one omits (PARI's coefficients of \texttt{lfunsympow} at $11$ and $121$ are $-11$ and $121$). The Rankin
series $\sum a_n(f)^2n^{-s}$ is not $L^{\nv}(\Sym f,s)$: it equals $\zeta^{(N)}(s-1)L^{\nv}(\Sym f,s)/\zeta^{(N)}(2s-2)$, and at a good prime it carries an extra factor $1+p^{1-s}$
(\Cref{ch:notation}, B-C7).
\item (\Cref{ch:moddeg}.) The formula of Zagier and Watkins $\degphi\cdot\vol(\Lambda_E)=4\pi^2c_E^2\langle f,f\rangle$ in general.
\end{enumerate}
\end{remark}

\begin{paperlink}\label{pl:modforms:petersson}
The Introduction of the paper says: ``The classical route to the degree is the Petersson norm. Up to explicit factors $\degphi$ is $\langle f,f\rangle$ divided by the area of the
period lattice of $E$ (Zagier \cite{Zagier85}, Watkins \cite{Watkins02})''. Remark rem:degree makes this exact: ``$\degphi\cdot\operatorname{vol}(\Lambda_E)=4\pi^2c_E^2\langle f,f\rangle$, with the
unnormalised Petersson product and the Manin constant $c_E>0$ in the numerator''. Both the normalisation and the place of $c_E$ are those checked in
\Cref{ex:modforms:petersson-numerics}. The orthogonality in the definition of $r_E$, in Lemma lem:crit ($f^\perp$) and in Lemma lem:old is that of
\Cref{prop:modforms:orthogonal-complement}; in the proof of Theorem thm:twist the paper uses instead ``the $\T$-stable complement of the $g$-line'', which coincides with the Petersson
complement for a newform $g$ (\Cref{ch:hecke}), and in Step 1 of the proof of Theorem thm:quant it insists on the $\T$-stable complement ``(not the Petersson complement, which is not
$\T$-stable when the $U_p$ are involved)'', because there the form $g^S$ is not a newform.
\end{paperlink}

\section{Integral forms, the Sturm bound and congruences}\label{sec:modforms:integral}

The paper's congruences live in $S_2(\Gamma_0(N);\Z)$ and $S_2(\Gamma_0(N);\Z_\ell)$: the congruence number is defined with an ``integral cusp form'', Lemma lem:crit uses the pairing
``$S_2(\Gamma_0(N);\Z)\times\T\to\Z$'' and its version over $\Z_\ell$, Theorem thm:twist needs a twist map that ``preserves integrality'', and Step 1 of the proof of Theorem thm:quant
projects a congruence inside $S_2(\Gamma_0(N);\Z_\ell)$. This section defines these modules, states Shimura's theorem that they are lattices of full rank, proves Sturm's bound in
the form needed, and deduces that a congruence between two forms of the same level can be decided on finitely many coefficients. It then proves the congruences of the
paper's Example for all $n$. The perfect pairing with the Hecke algebra is \Cref{ch:hecke}'s.

\begin{definition}[integral forms]\label{def:modforms:integral-forms}
For a subring $A\subseteq\C$ let $M_k(\Gamma_0(N);A)=\{f\in M_k(\Gamma_0(N)):a_n(f)\in A\text{ for all }n\ge0\}$ and $S_k(\Gamma_0(N);A)=S_k(\Gamma_0(N))\cap M_k(\Gamma_0(N);A)$, which are
$A$-modules. For an arbitrary commutative ring $R$ (such as $\Z_\ell$, $\F_\ell$ or $\mathcal O_\lambda$, which are not subrings of $\C$ in a preferred way) put
\[
S_k(\Gamma_0(N);R)=S_k(\Gamma_0(N);\Z)\otimes_\Z R,
\]
with the \emph{$q$-expansion map} $S_k(\Gamma_0(N);R)\to R[[q]]$, $f\otimes r\mapsto\sum_nra_n(f)q^n$.
\end{definition}

Three remarks. (a) $S_k(\Gamma_0(N);\Z)$ is \emph{saturated} in $\Z[[q]]$ by definition: if $g\in\Z[[q]]$ and $mg\in S_k(\Gamma_0(N);\Z)$ for an integer $m\ge1$, then $g\in S_k(\Gamma_0(N))$ and so
$g\in S_k(\Gamma_0(N);\Z)$. (b) Integrality refers to the expansion at $\infty$ only. Integrality at the other cusps is a different matter (for instance $f_{11}|_2S=-\frac1{11}f_{11}(z/11)$ by
\Cref{thm:modforms:level-eleven}(iii), which has non-integral coefficients in $q_{11}$); the paper never needs it. (c) For $R=\Z_\ell$ this is the paper's $S_2(\Gamma_0(N);\Z_\ell)$;
\Cref{cor:modforms:integral-lattice}(iii) shows that it is the module of forms in $S_2(\Gamma_0(N);\Q)\otimes\Q_\ell$ with $\ell$-integral $q$-expansion.

\begin{theorem}[Shimura; integral bases]\label{thm:modforms:integral-basis}
$S_k(\Gamma_0(N))$ has a basis over $\C$ consisting of forms with integer coefficients; the same holds for $M_k(\Gamma_0(N))$.
\end{theorem}

We do not prove this; see \cite[Theorem 3.52]{Shimura71} (number not verified) and \cite[Section 6.5]{DiamondShurman}. There are three standard proofs: through the
rationality of the Hecke eigenvalues and the Galois action on newforms (Shimura), through the algebraic theory of modular forms and its $q$-expansion principle
(Deligne-Rapoport and Katz, \cite{DeligneRapoport73}), and for $k=2$ through $H^1(X_0(N),\Z)$ and the Eichler-Shimura isomorphism (\Cref{ch:moddeg}). For $M_k(\SL_2(\Z))$ it is
\Cref{cor:modforms:level-one-structure}(iv): the monomials $E_4^aE_6^b$ have integer coefficients.

\begin{theorem}[Sturm's bound]\label{thm:modforms:sturm}
Let $\mathcal O$ be the ring of integers of a number field $K\subset\C$, $\lambda$ a nonzero prime ideal of $\mathcal O$, and $f\in M_k(\Gamma_0(N);\mathcal O)$. If $a_n(f)\in\lambda$ for all
$n\le\frac{k\mu(N)}{12}$, then $a_n(f)\in\lambda$ for all $n$. For $k=2$ the condition is on $0\le n\le B_2(N)=\lfloor\mu(N)/6\rfloor$.
\end{theorem}

\begin{proof}[Sketch of proof]
\emph{Level one.} Let $d$ be the dimension of $M_k(\SL_2(\Z))$. By \Cref{cor:modforms:level-one-structure} the module $M_k(\SL_2(\Z);\Z)$ has a basis $g_0,\dots,g_{d-1}$ with $a_i(g_j)=\delta_{ij}$ for
$0\le i,j<d$ (row-reduce the monomials $E_4^aE_6^b\Delta^c$, $4a+6b+12c=k$, whose expansions are triangular with leading coefficients $1$), and the valence formula gives
$v_\infty(f)\le d-1\le k/12$ for every nonzero $f\in M_k$. For $f\in M_k(\SL_2(\Z);\mathcal O)$, $f-\sum_{j<d}a_j(f)g_j$ vanishes to order at least $d$ at $\infty$, so it is $0$; hence if
$a_j(f)\in\lambda$ for $j<d$, every coefficient of $f$ is in $\lambda$. \emph{Level $N$.} Replace $f$ by the norm $F=\prod_if|_k\gamma_i\in M_{k\mu}(\SL_2(\Z))$ of the proof of
\Cref{thm:modforms:valence-gamma0}. The coefficients of the conjugates $f|_k\gamma_i$ are algebraic, and the delicate point, which rests on the integrality statements behind
\Cref{thm:modforms:integral-basis}, is that after scaling by a suitable constant they become $\lambda$-integral and the $\lambda$-adic order of vanishing of $F$ at $\infty$ is at least
that of $f$; then the level-one case in weight $k\mu$, with bound $k\mu/12$, applies to $F$. See \cite{Sturm87} and \cite[Chapter 9]{Stein} (theorem number not verified).
\end{proof}

The values $B_2(N)$ at the levels of the paper's congruences are $B_2(11)=2$, $B_2(20)=6$, $B_2(37)=6$, $B_2(52)=14$, $B_2(99)=24$, $B_2(121)=22$, $B_2(539)=112$,
$B_2(980)=336$; they agree with PARI's \texttt{mfsturm}. The paper's ``Sturm's bound $\frac16N\prod_{p\mid N}(1+\frac1p)$'' is $\mu(N)/6$, whose integer part is $B_2(N)$.

\begin{corollary}[the integral lattice and its $q$-expansions]\label{cor:modforms:integral-lattice}
Let $d=\dim S_2(\Gamma_0(N))$, $B=B_2(N)$ and $L=S_2(\Gamma_0(N);\Z)$.
\begin{enumerate}[label=(\roman*)]
\item $L$ is free of rank $d$, the map $L\otimes_\Z\C\to S_2(\Gamma_0(N))$ is an isomorphism, and $S_2(\Gamma_0(N);\Q)=L\otimes_\Z\Q$.
\item The truncation $\tau_B\colon L\to\Z^B$, $f\mapsto(a_1(f),\dots,a_B(f))$, is injective with torsion-free cokernel; so $\tau_B(L)$ is a direct summand of $\Z^B$.
\item For every commutative ring $R$ the $q$-expansion map $S_2(\Gamma_0(N);R)\to R[[q]]$ is injective, and so is its truncation $S_2(\Gamma_0(N);R)\to R^B$. For $R=\Z_\ell$ its image is
$(S_2(\Gamma_0(N);\Q)\otimes\Q_\ell)\cap\Z_\ell[[q]]$; for $R=\F_\ell$, an element is $0$ iff its first $B$ coefficients are.
\item For $f,g\in L$ and an integer $m\ge1$, $f\equiv g\pmod m$ coefficientwise iff $a_n(f)\equiv a_n(g)\pmod m$ for $1\le n\le B$. For $f,g\in S_2(\Gamma_0(N);\mathcal O)$ and a prime $\lambda$
of $\mathcal O$, $f\equiv g\pmod\lambda$ coefficientwise iff $a_n(f)\equiv a_n(g)\pmod\lambda$ for $n\le B$.
\item Forms of levels $M_1,M_2$ are compared in $S_2(\Gamma_0(N))$ with $N=\operatorname{lcm}(M_1,M_2)$ (\Cref{prop:modforms:mf-basic}(c)), with the bound $B_2(N)$. A congruence at the primes
$p\notin\Sigma$ only is not of this kind: it must first be turned into a congruence at all $n$ (by depletion, \Cref{ch:hecke}), which raises the level and the bound.
\end{enumerate}
\end{corollary}

\begin{proof}
(i) By \Cref{thm:modforms:integral-basis} there are $f_1,\dots,f_d\in L$ forming a basis of $S_2(\Gamma_0(N))$. The vectors $\tau_B(f_j)\in\Z^B$ are linearly independent: otherwise, since they are
integer vectors, there are integers $c_j$, not all $0$, with $\sum c_j\tau_B(f_j)=0$, and $h=\sum c_jf_j\in L$ is a nonzero form (the $f_j$ are independent) with $a_n(h)=0$ for $n\le B$;
by \Cref{thm:modforms:sturm} with $\mathcal O=\Z$ every coefficient of $h$ is divisible by every prime, so $h=0$, a contradiction. Hence $\tau_B$ is injective on $S_2(\Gamma_0(N))$ (which
the $f_j$ span), so $L$ embeds in $\Z^B$ and is free of finite rank $r$. Elements of $L$ that are independent over $\Z$ have independent images in $\Q^B$, hence in $\C^B$, so $r\le d$; and
$r\ge d$ because of the $f_j$. So a $\Z$-basis of $L$ is a $\C$-basis of $S_2(\Gamma_0(N))$. If $f$ has rational coefficients, write $f=\sum c_je_j$ in a $\Z$-basis $(e_j)$ of $L$; the $c_j$
solve a linear system $\sum c_j\tau_B(e_j)=\tau_B(f)$ with rational data and a unique solution, so $c_j\in\Q$.
(ii) Injectivity was shown in (i). If $g\in L$ and $\tau_B(g)\in m\Z^B$ with $m>1$, let $p$ be a prime dividing $m$. Then $a_n(g)\equiv0\pmod p$ for $n\le B$ ($a_0=0$), so for all $n$ by
\Cref{thm:modforms:sturm}, and $g/p\in\Z[[q]]\cap S_2(\Gamma_0(N))=L$ with $\tau_B(g/p)\in\frac mp\Z^B$. By induction $g/m\in L$. So $\Z^B/\tau_B(L)$ is torsion-free, hence free, and the sequence
$0\to L\to\Z^B\to\Z^B/\tau_B(L)\to0$ splits (\Cref{ch:modules}).
(iii) Tensoring a split injection with $R$ gives a split injection $L\otimes R\to R^B$; it is the composite of the $q$-expansion map with truncation, so both are injective. For $R=\Z_\ell$:
an element of $L\otimes\Z_\ell$ has $q$-expansion in $\Z_\ell[[q]]$ and lies in $L\otimes\Q_\ell=S_2(\Gamma_0(N);\Q)\otimes\Q_\ell$. Conversely let $F=\sum c_je_j$, $c_j\in\Q_\ell$, have coefficients in
$\Z_\ell$. Since $\tau_B(e_1),\dots,\tau_B(e_d)$ extend to a basis of $\Z^B$ by (ii), the coordinates of $\tau_B(F)\in\Z_\ell^B$ in that basis are in $\Z_\ell$, and they are the $c_j$; so
$F\in L\otimes\Z_\ell$. The statement for $\F_\ell$ is injectivity of the truncation.
(iv) The first statement is the argument of (ii) applied to $h=f-g$: if $\tau_B(h)\in m\Z^B$ then $h/m\in L$. The second is \Cref{thm:modforms:sturm} applied to $f-g\in M_2(\Gamma_0(N);\mathcal O)$.
(v) is a remark; for instance the depletion of a form of level $121$ at $3$ has level $121\cdot9=1089$ and $B_2(1089)=264$.
\end{proof}

\begin{paperlink}\label{pl:modforms:crit}
Lemma lem:crit ends: ``Every step holds verbatim over $\Z_\ell$, the pairing being perfect after tensoring with $\Z_\ell$.'' On the side of forms, the module there is
$S_2(\Gamma_0(N);\Z_\ell)=S_2(\Gamma_0(N);\Z)\otimes\Z_\ell$ (\Cref{def:modforms:integral-forms}); \Cref{cor:modforms:integral-lattice}(iii) says that it embeds in $\Z_\ell[[q]]$, so that congruences modulo $\ell^t$
coefficientwise make sense for its elements, and that it consists exactly of the $\ell$-integral forms. Step 1 of the proof of Theorem thm:quant uses the same module: it speaks of
``$g^S\equiv h\pmod{\ell^t}$ coefficientwise with $h\in S_2(\Gamma_0(N);\Z_\ell)$'', and continues ``projecting
$h$ to the $\mathfrak m$-component (an idempotent of $\T_\ell$, so the projection stays in the complement and the congruence survives)''. Here the idempotent $e_{\mathfrak m}\in\T_\ell$
acts on $S_2(\Gamma_0(N);\Z_\ell)$ (the Hecke operators preserve $S_2(\Gamma_0(N);\Z)$, \Cref{ch:hecke}), and if $g^S-h\in\ell^tS_2(\Gamma_0(N);\Z_\ell)$ then $e_{\mathfrak m}(g^S-h)=g^S-e_{\mathfrak m}h\in\ell^tS_2(\Gamma_0(N);\Z_\ell)$ because $e_{\mathfrak m}g^S=g^S$; by
the injectivity of (iii) this is a congruence of $q$-expansions. In Theorem thm:twist the twist map $h\mapsto h\otimes\chi_S=\sum\chi_S(n)a_n(h)q^n$ ``maps $S_2(\Gamma_0(N_S))$ into
$S_2(\Gamma_0(\operatorname{lcm}(N_S,D^2)))\subset S_2(\Gamma_0(N))$, preserves integrality'': the inclusion of spaces is \Cref{ch:hecke}'s, and integrality is clear on $q$-expansions because $\chi_S(n)\in\{0,\pm1\}$.
\end{paperlink}

\begin{example}[the paper's congruences hold for all $n$]\label{ex:modforms:paper-congruences}
The paper's Example after the proof of Theorem thm:qual lists four congruences. In each case the two forms lie in the same space $S_2(\Gamma_0(N);\Z)$ (for the twisted forms this
uses the level of a twist, \Cref{ch:hecke}), so by \Cref{cor:modforms:integral-lattice}(iv) a check of $a_n$ for $n\le B_2(N)$ proves the congruence for every $n$.
\begin{enumerate}[label=(\alph*)]
\item $f_{\curve{52a}}\equiv f_{\curve{26b}}(z)-f_{\curve{26b}}(2z)\pmod3$ in $S_2(\Gamma_0(52);\Z)$ (the second form by \Cref{prop:modforms:mf-basic}(e)), $B_2(52)=14$; here $a_2(\curve{26b1})=1$.
\item $f_{\curve{121a}}\equiv f_{\curve{11a}}\otimes\chi_{-11}=f_{\curve{121d}}\pmod3$ in $S_2(\Gamma_0(121);\Z)$, $B_2(121)=22$. At the primes $3$ and $11$, excluded in the original version of
the paper: $a_3=2$ and $-1$, $a_{11}=0$ and $0$.
\item $f_{\curve{539c}}\equiv f_{\curve{11a}}\otimes\chi_{-7}=f_{\curve{539d}}\pmod3$ in $S_2(\Gamma_0(539);\Z)$, $B_2(539)=112$. Here $a_3=-2$ and $1$, $a_7=0$ and $0$, $a_{11}=1$ and $1$.
\item $f_{\curve{980i}}\equiv f_{\curve{20a}}\otimes\chi_{-7}=f_{\curve{980g}}\pmod5$ in $S_2(\Gamma_0(980);\Z)$, $B_2(980)=336$. Here $a_2=0$ and $0$, $a_5=1$ and $1$, $a_7=0$ and $0$.
\end{enumerate}
We checked each congruence for all $n\le3000$, far beyond the bounds, and the equalities $f_{\curve{11a}}\otimes\chi_{-11}=f_{\curve{121d}}$, $f_{\curve{11a}}\otimes\chi_{-7}=f_{\curve{539d}}$,
$f_{\curve{20a}}\otimes\chi_{-7}=f_{\curve{980g}}$ of twisted and tabulated forms for all $n\le3000$ (each of these is again a statement in one space $S_2(\Gamma_0(N);\Z)$, so the check up to
$B_2(N)$ proves it once the twist is known to lie there). The Eisenstein congruence of \Cref{ex:modforms:eisenstein-congruence} is the same argument in $M_2(\Gamma_0(11);\Z)$.
\end{example}

\begin{paperlink}\label{pl:modforms:example}
The revised paper says, at the start of its Example: ``in each case the two forms lie in the same space $S_2(\Gamma_0(N))$ and have integer coefficients, so a check up to Sturm's bound
\cite{Sturm87}, $\frac16N\prod_{p\mid N}(1+\frac1p)$, proves the congruence for all $n$.'' This is \Cref{cor:modforms:integral-lattice}(iv), and \Cref{ex:modforms:paper-congruences}
supplies the numbers. The version of 10 October 2026 stated the congruences only at the primes away from a finite set, with the checks ``for $p<500$'' and ``for $p<1000$''; such
checks already contain a proof once the excluded $a_p$ are compared as well, because the $a_n$ with $n\le B_2(N)$ are determined by the $a_p$ with $p\le B_2(N)$ through the Hecke recursion
(\Cref{ch:hecke}). In Lemma lem:old the paper passes from the $a_p$ to all $a_n$ by that recursion, not by Sturm's bound.
\end{paperlink}

\section{The Atkin-Lehner involutions}\label{sec:modforms:atkin-lehner}

The Atkin-Lehner involutions are the symmetries of $X_0(N)$ that do not come from $\SL_2(\Z)$. The paper meets the most important one, $w_N$, in Section 4.1: the pairing $\hat\delta_!$
of Diamond, Flach and Guo is the cup product composed with ``the operator $[UwU]\omega$, $w=\begin{psmallmatrix}0&-1\\M^\Sigma&0\end{psmallmatrix}$ \cite[(13), (14)]{DFG04},
which for weight $2$ and trivial character is the Atkin-Lehner involution $w_{M^\Sigma}$ with no scalar factor''. (In \cite[Section 1.5.3]{DFG04} the operator is written $[UwU]_\omega$, $w$ being the element with these entries of $\prod_{p\mid N}\GL_2(\Z_p)$
at level $N$, and $[UwU]_\omega$ ``coincides with $\psi(-1_\infty)N^{k-2}[Uw^{-1}U]_\omega$''; for $k=2$ and trivial character that scalar is $1$, as the paper says; checked.) Through $\hat\delta_!$ the involution enters Proposition prop:pullback,
Lemma lem:hecke and Step 2 of the proof of Theorem thm:quant. This section defines $W_d$ for every exact divisor $d$ of $N$, proves that the $w_d$ are commuting involutions of
$X_0(N)$ and of the spaces of forms, that they are isometries for the Petersson product and the cup product, computes their action on cusps, and explains why their action
on $q$-expansions is not simple. The identity ``the adjoint of $T_n$ for the cup product is $w_NT_nw_N$'' and the pairing $\hat\delta_!$ itself are \Cref{ch:moddeg}'s.

\begin{definition}[Atkin-Lehner matrices]\label{def:modforms:atkin-lehner}
Let $d\,\|\,N$, that is, $d\mid N$ and $\gcd(d,N/d)=1$. An \emph{Atkin-Lehner matrix} for $d$ is an integer matrix
\[
W_d=\begin{pmatrix}dx&y\\Nz&dw\end{pmatrix},\qquad\det W_d=d^2xw-Nyz=d .
\]
Such matrices exist: dividing by $d$, the condition is $dxw-(N/d)yz=1$, which has solutions since $\gcd(d,N/d)=1$. For $d=N$ we take $W_N=\sm0{-1}N0$ ($x=w=0$, $y=-1$, $z=1$), and for
$d=1$, $W_1=I$. We write $w_d$ for the induced map of $X_0(N)$ and for the operator $f\mapsto f|_kW_d$ on forms.
\end{definition}

\begin{proposition}[the Atkin-Lehner involutions]\label{prop:modforms:atkin-lehner}
Let $d,d'\,\|\,N$, and let $\mathcal W_d$ be the set of all Atkin-Lehner matrices for $d$.
\begin{enumerate}[label=(\roman*)]
\item $\Gamma_0(N)\mathcal W_d=\mathcal W_d=\mathcal W_d\Gamma_0(N)$, and any two elements of $\mathcal W_d$ differ by left (or right) multiplication by an element of $\Gamma_0(N)$. In particular $W_d$
normalises $\Gamma_0(N)$.
\item $W_d^2\in d\cdot\Gamma_0(N)$; explicitly $W_d^2=d\begin{psmallmatrix}dx^2+(N/d)yz&(x+w)y\\Nz(x+w)&(N/d)yz+dw^2\end{psmallmatrix}$.
\item $W_dW_{d'}\in\gcd(d,d')\cdot\mathcal W_{d*d'}$ with $d*d'=dd'/\gcd(d,d')^2$.
\item $z\mapsto W_dz$ induces a holomorphic involution $w_d$ of $X_0(N)$, independent of the choice of $W_d$; the $w_d$ commute, $w_dw_{d'}=w_{d*d'}$, and they form a group of automorphisms
of $X_0(N)$ that is a quotient of $(\Z/2)^{\omega(N)}$ ($\omega(N)$ the number of primes dividing $N$). The involution $w_N$ is $z\mapsto-1/(Nz)$, and in the moduli description
of \Cref{thm:modforms:moduli} it is $(E,C)\mapsto(E/C,E[N]/C)$.
\item For even $k$, $f\mapsto f|_kW_d$ is a well-defined involution of $M_k(\Gamma_0(N))$ preserving $S_k(\Gamma_0(N))$; these operators commute, and $w_N=\prod_{p^e\|N}w_{p^e}$.
\item In weight $2$, $w_d^*(\omega_f)=\omega_{f|_2W_d}$ for $f\in M_2(\Gamma_0(N))$.
\end{enumerate}
\end{proposition}

\begin{proof}
(i) For $\gamma=\sm ab{Nc}e\in\Gamma_0(N)$,
\[
\gamma W_d=\begin{pmatrix}adx+bNz&ay+bdw\\Ncdx+eNz&Ncy+edw\end{pmatrix},
\]
whose upper left and lower right entries are divisible by $d$ (as $d\mid N$), whose lower left entry is divisible by $N$, and whose determinant is $d$; so $\gamma W_d\in\mathcal W_d$, and likewise
$W_d\gamma\in\mathcal W_d$. For $W,W'\in\mathcal W_d$ with entries $(dx,y,Nz,dw)$ and $(dx',y',Nz',dw')$, the matrix $W'W^{-1}=\frac1dW'\sm{dw}{-y}{-Nz}{dx}$ has entries
$dx'w-\frac Ndy'z$, $xy'-x'y$, $N(z'w-w'z)$, $dxw'-\frac Ndyz'$: integers, with the lower left divisible by $N$, and determinant $1$. So $W'\in\Gamma_0(N)W$, and similarly
$W'\in W\Gamma_0(N)$. Then $W_d\gamma W_d^{-1}\in\mathcal W_dW_d^{-1}\subseteq\Gamma_0(N)$.
(ii) Multiply out: $W_d^2=\sm{d^2x^2+Nyz}{d(x+w)y}{dN(x+w)z}{Nyz+d^2w^2}$, which is $d$ times the displayed matrix; that matrix is integral with lower left entry divisible by $N$ and
determinant $\det(W_d)^2/d^2=1$.
(iii) If $\gcd(d,d')=1$ then $dd'\,\|\,N$, and with $W_{d'}=\sm{d'x'}{y'}{Nz'}{d'w'}$,
\[
W_dW_{d'}=\begin{pmatrix}dd'xx'+Nyz'&dxy'+d'yw'\\N(d'zx'+dwz')&Nzy'+dd'ww'\end{pmatrix}.
\]
The upper left and lower right entries are divisible by $dd'$ because $dd'\mid N$, the lower left entry is divisible by $N$, and the determinant is $dd'$: so $W_dW_{d'}\in\mathcal W_{dd'}$. In general write $g=\gcd(d,d')$, $d=ge$, $d'=ge'$. For a prime $p\mid g$, $v_p(d)=v_p(N)=v_p(d')$ (both are
exact divisors), so $v_p(g)=v_p(d)$ and $g$, $e$, $e'$ are pairwise coprime exact divisors of $N$. By the coprime case and (i), $\mathcal W_d=\mathcal W_g\mathcal W_e=\mathcal W_e\mathcal W_g$ and $\mathcal W_{d'}=\mathcal W_g\mathcal W_{e'}$, so by
(ii) $W_dW_{d'}\in\mathcal W_e\mathcal W_g\mathcal W_g\mathcal W_{e'}=g\,\mathcal W_e\Gamma_0(N)\mathcal W_{e'}=g\,\mathcal W_{ee'}$, and $ee'=d*d'$.
(iv) A matrix of $\GL_2^+(\Q)$ that normalises $\Gamma_0(N)$ induces a bijection of $\Gamma_0(N)\backslash\HH^*$ (it maps $\PP^1(\Q)$ to itself), holomorphic in the charts of
\Cref{thm:modforms:y0-riemann-surface,thm:modforms:x0-compact} because it conjugates stabilisers to stabilisers. By (i) the map does not depend on the choice of $W_d$, by (ii) its square
is induced by $d\gamma$, which acts as $\gamma\in\Gamma_0(N)$, that is, trivially; and by (iii) $w_dw_{d'}=w_{d*d'}=w_{d'}w_d$. The map $S\mapsto w_{\prod_{p\in S}p^{v_p(N)}}$ from subsets of the primes
dividing $N$, with symmetric difference as group law, is then a surjective homomorphism from $(\Z/2)^{\omega(N)}$. For $W_N$, $W_Nz=-1/(Nz)$, and the moduli statement is
\Cref{thm:modforms:moduli}(iii).
(v) By (i) two choices of $W_d$ differ by $\gamma\in\Gamma_0(N)$, which acts trivially on $M_k(\Gamma_0(N))$; for $\gamma\in\Gamma_0(N)$, $f|_kW_d|_k\gamma=f|_k\gamma'|_kW_d=f|_kW_d$ with
$\gamma'=W_d\gamma W_d^{-1}\in\Gamma_0(N)$; and holomorphy at the cusps follows as in \Cref{prop:modforms:mf-basic}(e), writing $W_d\sigma=\sigma'\beta$ with $\sigma'\in\SL_2(\Z)$ and $\beta$ upper triangular.
By (ii) and \Cref{lem:modforms:slash}(ii), $f|_kW_d^2=f|_k(d\gamma)=f$ for $k$ even (with Diamond-Shurman's normalisation one would get $d^{k-2}f$). By (iii) the operators commute, and
$w_N=\prod w_{p^e}$ since $N$ is the $*$-product of the $p^e$.
(vi) \Cref{lem:modforms:slash}(iii).
\end{proof}

The three normalisations of $W_N$ in use, $\sm0{-1}N0$, $\sm01{-N}0$ and $\sm{0}{1/\sqrt N}{-\sqrt N}0$, differ by the scalars $-1$ and $\sqrt N$, which act trivially on $\HH$ and, with the
slash operator of \Cref{def:modforms:slash}, on forms of even weight. So the paper's $w$ in Section 4.1 is exactly $w_{M^\Sigma}$, with no hidden scalar: in weight $2$ the factor
$(\det)^{k/2}=(\det)^{k-1}=\det$ is the same in all conventions.

\begin{proposition}[Atkin-Lehner and the Petersson and cup products]\label{prop:modforms:atkin-lehner-petersson}
Let $k$ be even and $d\,\|\,N$.
\begin{enumerate}[label=(\roman*)]
\item For $f,g\in S_k(\Gamma_0(N))$: $\langle f|_kW_d,g|_kW_d\rangle=\langle f,g\rangle$ and $\langle f|_kW_d,g\rangle=\langle f,g|_kW_d\rangle$. So $w_d$ is a self-adjoint isometric involution of $S_k(\Gamma_0(N))$, which is the
orthogonal direct sum of its $\pm1$-eigenspaces $S_k(\Gamma_0(N))^\pm$, and the $w_d$ are simultaneously diagonalisable.
\item For $k=2$, $\dim S_2(\Gamma_0(N))^+$ (for $w_N$) is the genus of the quotient curve $X_0(N)/w_N$.
\item $w_N$ is an involutive automorphism of the compact Riemann surface $X_0(N)$, so $w_N^*$ preserves the cup product on $H^1(X_0(N),\Z)$ and $(w_N^*)^2=1$. If
$f\in S_2(\Gamma_0(N))$ and $f|_2W_N=\varepsilon_N(f)f$ with $\varepsilon_N(f)=\pm1$, then $w_N^*$ acts by $\varepsilon_N(f)$ on $\omega_f$, on $\overline{\omega_f}$, and on the real plane of $H^1(X_0(N),\R)$ spanned by
$\operatorname{Re}\omega_f$ and $\operatorname{Im}\omega_f$.
\end{enumerate}
\end{proposition}

\begin{proof}
(i) The first identity is \Cref{prop:modforms:petersson-change}(ii). Applying it to $f$ and $g|_kW_d$ and using $g|_kW_d|_kW_d=g$ gives the second. A self-adjoint involution of a finite-dimensional
Hermitian space is diagonalisable with eigenvalues $\pm1$ and orthogonal eigenspaces, and commuting diagonalisable operators are simultaneously diagonalisable.
(ii) The quotient $X^+=X_0(N)/\langle w_N\rangle$ of a compact Riemann surface by a finite group of automorphisms is a compact Riemann surface, and $\pi^+\colon X_0(N)\to X^+$ is holomorphic of
degree $2$ (stated; \cite[Chapter III, Section 3]{Miranda}, section not verified). Pullback identifies $H^0(X^+,\Omega^1)$ with the $w_N$-invariant holomorphic differentials on $X_0(N)$: it is
injective, and an invariant holomorphic differential descends, by the local computation in the proof of \Cref{thm:modforms:weight-two-differentials} with $e=2$ at the fixed points
of $w_N$. By (vi) of \Cref{prop:modforms:atkin-lehner} the invariant differentials are the $\omega_f$ with $f\in S_2(\Gamma_0(N))^+$, and $\dim H^0(X^+,\Omega^1)=g(X^+)$ (\Cref{ch:curves}).
(iii) An automorphism of a compact Riemann surface preserves the orientation, so it preserves the integral of a $2$-form and hence the cup product (\Cref{ch:curves}). By
\Cref{prop:modforms:atkin-lehner}(vi), $w_N^*\omega_f=\omega_{f|_2W_N}=\varepsilon_N(f)\omega_f$; conjugating, $w_N^*\overline{\omega_f}=\varepsilon_N(f)\overline{\omega_f}$ as $\varepsilon_N(f)$ is real; and the real and imaginary parts are
combinations of $\omega_f$ and $\overline{\omega_f}$ with real coefficients. By \Cref{ch:curves}, $H^1(X_0(N),\C)=H^0(\Omega^1)\oplus\overline{H^0(\Omega^1)}$, so these classes live in $H^1$.
\end{proof}

Check of (ii) by Riemann-Hurwitz: if $w_N$ has $r$ fixed points then $2g(X_0(N))-2=2(2g(X^+)-2)+r$. The number of fixed points is given by class numbers (for $N>4$ it is
$h(-4N)+h(-N)$ if $N\equiv3\pmod4$ and $h(-4N)$ otherwise, $h$ the class number of the order of that discriminant; stated, \cite{Ogg73}). For $N=11$: $r=h(-44)+h(-11)=3+1=4$, so
$g(X^+)=0$ and $S_2(\Gamma_0(11))^+=0$, consistent with $f_{11}|_2W_{11}=-f_{11}$. For $N=23$: $r=3+3=6$, $g(X^+)=0$. For $N=37$: $r=h(-148)=2$, $g(X^+)=1$, and indeed $\dim S_2(\Gamma_0(37))^+=1$
(the form of \curve{37a1}). For $N=43$: $r=h(-172)+h(-43)=3+1=4$, $g(X^+)=1=\dim S_2(\Gamma_0(43))^+$. The dimensions of the $+$-spaces were computed from PARI's Atkin-Lehner matrices.

\begin{paperlink}\label{pl:modforms:wN}
The pairing of Section 4.1 is $\hat\delta_!=\frac12\langle x,w_{M^\Sigma}^*y\rangle$ on classes from $X_0(M^\Sigma)$, ``perfect and alternating'', and ``the Hecke operators, $U_p$ included, are
self-adjoint for it''. What this chapter supplies is: the paper's $w$ is the matrix $W_{M^\Sigma}$ of \Cref{def:modforms:atkin-lehner}, with no scalar (the remark after
\Cref{prop:modforms:atkin-lehner}); $w_N$ is an involutive automorphism of $X_0(N)$ (\Cref{prop:modforms:atkin-lehner}(iv)), so $w_N^*=(w_N^{-1})^*$ and $w_N^*$ preserves the cup product
(\Cref{prop:modforms:atkin-lehner-petersson}(iii)), which is the input needed for the pairing $\langle x,w_N^*y\rangle$ to be perfect on each local component (\Cref{ch:curves}); and on the
$f$-part, $w_N^*$ acts by the sign $\varepsilon_N(f)$, on $f\,dz$ and on its conjugate (\Cref{prop:modforms:atkin-lehner-petersson}(iii)), which is the step ``on which it acts by $\varepsilon_N(f)$
($f|w_N=\varepsilon_N(f)f$ and $\varepsilon_N(f)$ is real)'' of the proof of Proposition prop:pullback. That a newform is an eigenvector of $w_N$ is \Cref{ch:hecke}'s; for $f_{11}$ it is
\Cref{thm:modforms:level-eleven}(iii). The identity $T_n^\dagger=w_NT_nw_N$ for the cup product, and $\hat\delta_!$ itself, are \Cref{ch:moddeg}'s.
\end{paperlink}

\begin{proposition}[Atkin-Lehner involutions and cusps]\label{prop:modforms:atkin-lehner-cusps}
\begin{enumerate}[label=(\roman*)]
\item $w_N$ exchanges $\infty$ and $0$, and maps the cusps with $\gcd(\text{denominator},N)=c$ to those with $\gcd(\text{denominator},N)=N/c$ (\Cref{prop:modforms:cusps-gamma0}(i)).
\item If $N$ is squarefree and $e\mid N$, then $w_e$ maps the cusp $1/c$ ($c\mid N$) to the cusp $1/\bigl(ce/\gcd(c,e)^2\bigr)$. The group $\{w_e:e\mid N\}$ is isomorphic to $(\Z/2)^{\omega(N)}$ and
acts simply transitively on the $2^{\omega(N)}$ cusps.
\end{enumerate}
\end{proposition}

\begin{proof}
(i) $W_N\infty=0$ and $W_N0=\infty$. Let $a/c'$ be a cusp in lowest terms with $\gcd(c',N)=c$. Then $W_N(a/c')=-c'/(Na)$. Write $c'=cc''$ with $\gcd(c'',N/c)=1$ and $\gcd(a,c')=1$; then
$-c'/(Na)=-c''/((N/c)a)$, and since $\gcd(c'',(N/c)a)=1$ this is in lowest terms with denominator $(N/c)a$, and $\gcd((N/c)a,N)=(N/c)\gcd(a,c)=N/c$ (a prime dividing $a$ and $c$ would divide
$a$ and $c'$).
(ii) For a prime $p\mid N$ and a cusp $u/v$ in lowest terms let $\delta_p(u/v)=1$ if $p\mid v$ and $0$ otherwise; by \Cref{prop:modforms:cusps-gamma0}(i), (iii) the vector $(\delta_p)_{p\mid N}$ is a
well-defined invariant of the cusp which separates the $2^{\omega(N)}$ cusps, the cusp $1/c$ having $\delta_p=1$ exactly for $p\mid c$. We show that $w_e$ changes $\delta_p$ exactly for $p\mid e$.
Let $W_e=\sm{ex}y{Nz}{ew}$ and $p\mid N$. If $p\nmid e$: $p$ divides the lower left entry $Nz$ and not the determinant $e$, so it divides neither the upper left nor the lower right entry.
For $u/v$ in lowest terms the image is $(exu+yv)/(Nzu+ewv)$; if $p\mid v$ then $p\nmid u$, the numerator is $\equiv exu\not\equiv0$ and the denominator $\equiv0\pmod p$, so $p$ divides the reduced
denominator; if $p\nmid v$ the denominator is $\equiv ewv\not\equiv0$, so $p$ does not divide the reduced denominator. So $\delta_p$ is unchanged. If $p\mid e$ (so $p\nmid N/e$): from
$exw-(N/e)yz=1$, $p\nmid y$. If $p\nmid v$, the numerator $exu+yv\equiv yv\not\equiv0$ and the denominator $\equiv0\pmod p$, so the image has $\delta_p=1$. So $w_e$ maps the set $X$ of
cusps with $\delta_p=0$ into the set $Y$ with $\delta_p=1$; both have $2^{\omega(N)-1}$ elements and $w_e$ is a bijection, so $w_e(X)=Y$ and $w_e(Y)=X$. Hence $w_e$ toggles $\delta_p$
for $p\mid e$ and fixes it for $p\nmid e$, which is the formula: the primes of the new denominator are those of $c$ and $e$ not common to both. The action of the $w_e$ on the vectors
$(\delta_p)$ is by symmetric difference, simply transitive; in particular $w_e\ne1$ for $e\ne1$, so the group is $(\Z/2)^{\omega(N)}$.
\end{proof}

We checked (ii) for $N=6$, $30$, $77$ with explicit matrices $W_e$. For $N$ not squarefree the action can have fixed points: on the six cusps of $X_0(20)$, $w_4$ fixes $\frac12$ and
$\frac1{10}$ (Exercise \ref{ex:modforms:ex-cusps20}).

\begin{example}[the action on $q$-expansions]\label{ex:modforms:atkin-lehner-q-expansions}
\begin{enumerate}[label=(\alph*)]
\item $f|_kW_N(z)=N^{k/2}(Nz)^{-k}f(-1/(Nz))$ is, up to the factor and the change of variable $z\mapsto Nz$, the expansion of $f$ at the cusp $0$. It is determined by $f$ through the
transformation law. Since $M_k(\Gamma_0(N))$ is finite-dimensional and $f\mapsto(a_0(f),\dots,a_B(f))$ is injective for $B=\lfloor k\mu(N)/12\rfloor$ (\Cref{thm:modforms:valence-gamma0}), each
$a_n(f|_kW_N)$ is a linear combination $\sum_{m\le B}c_{n,m}a_m(f)$ with constants $c_{n,m}$ depending on $N$, $k$ and $n$; but these constants depend on the whole space and can only be
computed through the transformation law (for instance from a basis whose expansions at $0$ are known). There is no uniform coefficientwise rule: $a_n(f|_kW_N)=\pm a_n(f)$ for all $n$ holds
only for eigenvectors of $w_N$.
\item (Oldforms; \Cref{ch:hecke}.) Let $h\in M_k(\Gamma_0(M))$, $p\nmid M$ prime, and $N=Mp$. Then $W_N=W_M\sm p001$ and $\sm p001W_N=W_M\cdot pI$, so by \Cref{prop:modforms:mf-basic}(e)
\[
(h|_kW_N)(z)=p^{k/2}(h|_kW_M)(pz),\qquad\bigl(h(pz)\bigr)\big|_kW_N=p^{-k/2}(h|_kW_M)(z).
\]
So $w_N$ exchanges the lines of $h(z)$ and $h(pz)$ when $h$ is a $w_M$-eigenvector, and it is not diagonal in the basis of $q$-expansions. For $h=f_{11}$, $k=2$, $p=2$: in the basis
$f_{11}(z),f_{11}(2z)$ of $S_2(\Gamma_0(22))$ (which they span, $g(X_0(22))=2$), $w_{22}$ has the matrix $\sm0{-1/2}{-2}0$, whose square is $I$ (checked numerically at a point of $\HH$).
\item (\Cref{ch:hecke,ch:lfunctions}; stated.) A newform $f\in S_k(\Gamma_0(N))$ is an eigenvector of every $w_d$, $d\,\|\,N$, with eigenvalue $\varepsilon_d(f)=\pm1$; for $k=2$ and $p\,\|\,N$,
$a_p(f)=-\varepsilon_p(f)$ (\cite{AtkinLehner70}, Theorem 3; part not verified); and for $k=2$ the sign of the functional equation of $L(f,s)$ is $-\varepsilon_N(f)$. Checks: \curve{11a1} has $\varepsilon_{11}=-1$,
$a_{11}=1$, root number $+1$; \curve{37a1} has $\varepsilon_{37}=+1$, $a_{37}=-1$, root number $-1$ and rank $1$.
\end{enumerate}
\end{example}

For (b): $W_M\sm p001=\sm0{-1}{Mp}0=W_N$ and $\sm p001W_N=\sm0{-p}{Mp}0=W_M\cdot pI$; then $h|_kW_N=(h|_kW_M)|_k\sm p001=p^{k/2}(h|_kW_M)(pz)$, and $h(pz)=p^{-k/2}h|_k\sm p001$ gives
$h(pz)|_kW_N=p^{-k/2}h|_kW_M|_k(pI)=p^{-k/2}h|_kW_M$.

\section{Exercises}\label{sec:modforms:exercises}

\begin{exercise}\label{ex:modforms:ex-reduce}
Find $\gamma\in\SL_2(\Z)$ with $\gamma z\in\mathcal F$ for $z=(2+i)/7$, following the proof of \Cref{thm:modforms:fundamental-domain}(i), and write $\gamma$ as a word in $S$ and $T$.
\emph{Answer:} $Sz=-\frac{14}5+\frac75i$, then $T^3$: $\gamma=T^3S=\sm3{-1}10$ and $\gamma z=\frac{1+7i}5$.
\end{exercise}

\begin{exercise}\label{ex:modforms:ex-gamma1}
Compute $[\SL_2(\Z):\Gamma_1(11)]$ and $[\SL_2(\Z):\Gamma(11)]$, and the degrees of $X_1(11)\to X(1)$ and $X(11)\to X(1)$. Explain why the degree is half the index.
\emph{Answer:} $120$ and $1320$; the degrees are $60$ and $660$, because $-I\notin\Gamma_1(11)$ and $\pm\gamma$ act in the same way (\Cref{conv:modforms:conventions}(c)).
\end{exercise}

\begin{exercise}\label{ex:modforms:ex-cusps-p3}
Let $p$ be a prime. Using \Cref{prop:modforms:cusps-gamma0}(ii), show that $\Gamma_0(p^3)$ has $2p$ cusps: one of width $p^3$, $p-1$ of width $p$, and $p$ of width $1$. Check that the
widths add up to $\mu(p^3)$. Deduce $g(X_0(27))=1$ and $g(X_0(125))=8$.
\emph{Hint:} $c=1,p,p^2,p^3$ give $\varphi(\gcd(c,p^3/c))=1,p-1,p-1,1$ cusps of widths $p^3,p,1,1$; the sum is $p^3+p^2$. For $125$: $\mu=150$, $\nu_2=2$, $\nu_3=0$, $\nu_\infty=10$.
\end{exercise}

\begin{exercise}\label{ex:modforms:ex-elliptic}
Find the elliptic points of $\Gamma_0(65)$ and of $\Gamma_0(91)$, by solving $u^2+1\equiv0$ and $u^2-u+1\equiv0$ modulo $N$, and compute $g(X_0(65))$ and $g(X_0(91))$.
\emph{Answer:} for $65$, $u\in\{8,18,47,57\}$ give $\nu_2=4$, and $\nu_3=0$ ($5\equiv2\pmod3$); $\mu=84$, $\nu_\infty=4$, $g=5$. For $91$, $\nu_2=0$ ($7\equiv3\pmod4$), and
$u\in\{10,17,75,82\}$ give $\nu_3=4$ ($7\equiv13\equiv1\pmod3$); $\mu=112$, $\nu_\infty=4$, $g=7$.
\end{exercise}

\begin{exercise}\label{ex:modforms:ex-genus-bound}
(a) Show that for a prime $p\ge5$, $g(X_0(p))=\lfloor\frac{p+1}{12}\rfloor-1$ if $p\equiv1\pmod{12}$ and $\lfloor\frac{p+1}{12}\rfloor$ otherwise, and deduce that $g(X_0(p))=0$ exactly for
$p\in\{2,3,5,7,13\}$. (b) Show $\nu_2(N),\nu_3(N)\le2^{\omega(N)}\le\nu_\infty(N)$ and hence $g(X_0(N))\ge1+\frac{\mu(N)-13\nu_\infty(N)}{12}$. (c) Show $\nu_\infty(p^e)\le2p^{\lfloor e/2\rfloor}$, so that
$\nu_\infty(N)/\mu(N)\le\prod_{p^e\|N}\frac{2p^{1+\lfloor e/2\rfloor-e}}{p+1}$, and conclude that $\mu(N)-13\nu_\infty(N)\ge12$, hence $g(X_0(N))\ge2$, for all but finitely many $N$, all of which
can be listed. \emph{Answer to the last part:} there are $64$ exceptions ($63$ with $N\ge2$), the largest being $N=169$; together with the table computed for $N\le2000$ this proves that the lists of
levels of genus $0$ and $1$ in \Cref{sec:modforms:modular-curves} are complete. (For (a): $\mu=p+1$, $\nu_\infty=2$, and $(\nu_2,\nu_3)$ is determined by $p\bmod12$.)
\end{exercise}

\begin{exercise}\label{ex:modforms:ex-sigma}
(a) Use $\dim M_8(\SL_2(\Z))=1$ to show $E_4^2=E_8$ and deduce $\sigma_7(n)=\sigma_3(n)+120\sum_{m=1}^{n-1}\sigma_3(m)\sigma_3(n-m)$. (b) Do the same with $E_4E_6=E_{10}$.
\emph{Answer to (b):} $11\sigma_9(n)=21\sigma_5(n)-10\sigma_3(n)+5040\sum_{m=1}^{n-1}\sigma_3(m)\sigma_5(n-m)$, using $E_{10}=1-264\sum\sigma_9(n)q^n$.
\end{exercise}

\begin{exercise}\label{ex:modforms:ex-j1728}
Show $j-1728=E_6^2/\Delta$. Deduce that, as functions of $z$, $j-1728$ has a double zero at $i$ and $j$ a triple zero at $\zeta_3$, and check that this is consistent with $j$ being an
isomorphism $X(1)\to\PP^1(\C)$ in the charts of \Cref{thm:modforms:y0-riemann-surface}.
\emph{Hint:} $E_4^3-1728\Delta=E_6^2$; in the chart $t=w^2$ at $i$ (respectively $t=w^3$ at $\zeta_3$) the zero is simple.
\end{exercise}

\begin{exercise}\label{ex:modforms:ex-residue}
The group $\Gamma_0(4)$ has the cusps $\infty$, $\frac12$, $0$ of widths $1,1,4$, genus $0$, and $\dim M_2(\Gamma_0(4))=2$. Compute the constant terms of $E_2^{(4)}$ at the three cusps by the method of
the proof of \Cref{prop:modforms:eisenstein-weight-two}, and verify the relation $\sum_sh_sb_0(f|_2\sigma_s)=0$ of \Cref{thm:modforms:weight-two-differentials}. Show that $E_2^{(2)}$ and $E_2^{(4)}$ form a
basis of $M_2(\Gamma_0(4))$.
\emph{Answer:} $-3$ at $\infty$; at $\frac12$, with $\sigma=\sm1021$, $\sm4001\sigma=\sm2111\sm2{-1}02$, so the constant is $1-\frac22=0$; at $0$, $1-\frac14$. Then
$1\cdot(-3)+1\cdot0+4\cdot\frac34=0$. Likewise $E_2^{(2)}$ has constant terms $-1$, $-1$, $\frac12$ (at $\frac12$: $\sm2001\sigma=\sm1011\sm2001$). The two vectors $(-1,-1,\frac12)$ and
$(-3,0,\frac34)$ are not proportional, so the forms are independent.
\end{exercise}

\begin{exercise}\label{ex:modforms:ex-eis19}
Let $f_{19}$ be the normalised generator of $S_2(\Gamma_0(19))$, $f_{19}=q-2q^3-2q^4+\cdots$. Prove that $a_n(f_{19})\equiv\sigma(n)-19\sigma(n/19)\pmod3$ for all $n$, using
$h=4\bigl(f_{19}+\frac1{24}E_2^{(19)}\bigr)\in M_2(\Gamma_0(19);\Z)$ and $B_2(19)=3$. Deduce $a_\ell(\curve{19a1})\equiv1+\ell\pmod3$ for primes $\ell\ne19$, and relate it to the rational point of
order $3$ of \curve{19a1}.
\emph{Hint:} $h=-3+4\sum(a_n-\sigma(n)+19\sigma(n/19))q^n$ has coefficients $-3,0,-12,-24$ for $n\le3$. (The identification of $f_{19}$ with the form of \curve{19a1} is \Cref{ch:moddeg}; we
checked the congruence for $n\le3000$.)
\end{exercise}

\begin{exercise}\label{ex:modforms:ex-eta20}
Prove that $f_{20}=\eta(2z)^2\eta(10z)^2\in S_2(\Gamma_0(20))$ by the method of \Cref{thm:modforms:level-eleven}(i). You may use Ligozat's formula \cite{Ligozat75}: the order of
$\prod_\delta\eta(\delta z)^{r_\delta}$ ($\delta\mid N$) at a cusp with denominator $c\mid N$, in the local parameter of $X_0(N)$, is $\frac N{24}\sum_\delta\frac{\gcd(c,\delta)^2r_\delta}{\gcd(c,N/c)\,c\,\delta}$.
\emph{Hint:} $X_0(20)$ has genus $1$ and no elliptic points, so a generator $h$ of $S_2(\Gamma_0(20))$ has a simple zero at each of the six cusps and no other zero. The form
$f_{20}^{12}=\Delta(2z)\Delta(10z)$ lies in $S_{24}(\Gamma_0(20))$ by \Cref{prop:modforms:mf-basic}(e), and the formula gives order $12$ at each of the six cusps $c=1,2,4,5,10,20$.
\end{exercise}

\begin{exercise}\label{ex:modforms:ex-petersson-pari}
(a) Compute $\langle f_{11},f_{11}\rangle$ without \texttt{mfpetersson}, by integrating over the fundamental domain $\mathcal F\cup\bigcup_{j=0}^{10}ST^j\mathcal F$ of \Cref{prop:modforms:coset-domains}. Show
first, from $W_{11}=S\sm{11}001$ and $f_{11}|_2W_{11}=-f_{11}$, that $(f_{11}|_2S)(z)=-\frac1{11}f_{11}(z/11)$, and deduce from \Cref{prop:modforms:petersson}(ii) that
$\langle f_{11},f_{11}\rangle=\int_{\mathcal F}\bigl(\abs{f_{11}(z)}^2+\frac1{121}\sum_{j=0}^{10}\abs{f_{11}(\frac{z+j}{11})}^2\bigr)dx\,dy$. Evaluate the integral numerically with the first $400$
coefficients. (b) The curves \curve{11a2} and \curve{11a3} receive parametrisations $X_0(11)\to\curve{11a1}\to\curve{11a2}$, $\curve{11a3}$ of degree $5$. Compute
$\vol(\Lambda_{\curve{11a2}})$ and $\vol(\Lambda_{\curve{11a3}})$ and use $\deg\varphi\cdot\vol(\Lambda)=4\pi^2c_\varphi^2\langle f,f\rangle$ to find the constants $c_\varphi$ of these two parametrisations.
\emph{Answer:} (a) $f_{11}|_2S=f_{11}|_2W_{11}|_2\sm{1/11}001=-\frac1{11}f_{11}(z/11)$, and $(f|_2ST^j)(z)=(f|_2S)(z+j)$; in weight $2$, $y^2\,d\mu=dx\,dy$. Numerical integration
(truncating at $y=40$) gives $0.04690014787349$, in agreement with $12\cdot\texttt{mfpetersson}=0.0469001478734952$ of \Cref{ex:modforms:petersson-numerics}. (b) The volumes are $\frac15$ and $5$
times $\vol(\Lambda_{\curve{11a1}})=1.8515\ldots$, so $c_\varphi^2=1$ and $25$: the constants are $1$ and $5$. With the Manin constant in the
denominator one would get the impossible value $\frac15$ for \curve{11a3}.
\end{exercise}

\begin{exercise}\label{ex:modforms:ex-level22}
Let $f=f_{11}$. (a) Show $\langle f,f\rangle_{\Gamma_0(22)}=3\langle f,f\rangle_{\Gamma_0(11)}$. (b) Deduce from \Cref{prop:modforms:petersson-change}(iv), with $\alpha=\sm2001$, that
$\langle f(2z),f(2z)\rangle_{\Gamma_0(22)}=\frac34\langle f,f\rangle_{\Gamma_0(11)}$.
\emph{Hint:} $\alpha\Gamma_0(22)\alpha^{-1}=\Gamma_0(11)\cap\Gamma^0(2)$, where $\Gamma^0(2)$ consists of the matrices with $b$ even; it has index $3$ in $\Gamma_0(11)$, and $f(2z)=\frac12f|_2\alpha$. (PARI gives
$-\langle f,f\rangle_{\Gamma_0(11)}$ for $\langle f(z),f(2z)\rangle_{\Gamma_0(22)}$, which is $\frac{a_2(f)}2\langle f,f\rangle_{\Gamma_0(11)}$; \Cref{ch:hecke} explains the factor.)
\end{exercise}

\begin{exercise}\label{ex:modforms:ex-primes-to-n}
Let $f,g\in S_2(\Gamma_0(N);\Z)$ be normalised eigenforms of all the Hecke operators $T_n$, $n\ge1$ (\Cref{ch:hecke}: then $a_{mn}=a_ma_n$ for $\gcd(m,n)=1$,
$a_{p^{r+1}}=a_pa_{p^r}-p\,a_{p^{r-1}}$ for $p\nmid N$, and $a_{p^r}=a_p^r$ for $p\mid N$). Show that if $a_p(f)\equiv a_p(g)\pmod\ell$ for all primes $p\le B_2(N)$, then $f\equiv g\pmod\ell$
coefficientwise.
\emph{Hint:} the recursions express $a_n$, $n\le B_2(N)$, through the $a_p$ with $p\le n$, by polynomials with integer coefficients; then use \Cref{cor:modforms:integral-lattice}(iv). The
hypothesis includes the primes $p\mid N$, whose $a_p$ are not determined by the others, and the prime $p=\ell$. For the congruence modulo $3$ at level $121$ these are the primes
excluded in the original version of the paper ($p=\ell=3$, and $p=11$ dividing $N$), which is why \Cref{ex:modforms:paper-congruences}(b) lists $a_3$ and $a_{11}$ separately.
\end{exercise}

\begin{exercise}\label{ex:modforms:ex-al-existence}
Find explicit Atkin-Lehner matrices $W_4$ and $W_5$ for $N=20$, and check $W_4^2\in4\Gamma_0(20)$ and $W_4W_5\in\mathcal W_{20}$.
\emph{Answer:} for instance $W_4=\sm{-4}1{-20}4$ ($x=-1$, $y=1$, $z=-1$, $w=1$) and $W_5=\sm51{20}5$ ($x=y=z=w=1$). Then $W_4^2=-4I=4\cdot(-I)$, and
$W_4W_5=\sm01{-20}0=-W_{20}$, an Atkin-Lehner matrix for $20$.
\end{exercise}

\begin{exercise}\label{ex:modforms:ex-cusps20}
List the six cusps of $X_0(20)$ (\Cref{ex:modforms:cusp-lists}) and compute the permutations induced by $w_4$, $w_5$, $w_{20}$. Show that the group $\{1,w_4,w_5,w_{20}\}$ has two orbits on the
cusps, so that \Cref{prop:modforms:atkin-lehner-cusps}(ii) fails for $N$ not squarefree.
\emph{Answer:} $w_4$: $0\leftrightarrow\frac14$, $\frac15\leftrightarrow\infty$, $\frac12$ and $\frac1{10}$ fixed; $w_5$: $0\leftrightarrow\frac15$, $\frac12\leftrightarrow\frac1{10}$, $\frac14\leftrightarrow\infty$; $w_{20}$:
$0\leftrightarrow\infty$, $\frac12\leftrightarrow\frac1{10}$, $\frac14\leftrightarrow\frac15$. The orbits are $\{0,\frac14,\frac15,\infty\}$ and $\{\frac12,\frac1{10}\}$.
\end{exercise}

\begin{exercise}\label{ex:modforms:ex-al-p2}
Let $h\in S_2(\Gamma_0(M))$ with $h|_2W_M=\varepsilon_M(h)h$, $p\nmid M$ prime, and $N=Mp^2$. Show that in the basis $h(z),h(pz),h(p^2z)$ the operator $w_N$ has the matrix
$\varepsilon_M(h)\begin{psmallmatrix}0&0&p^{-2}\\0&1&0\\p^2&0&0\end{psmallmatrix}$, and find its eigenvectors. For $h=f_{11}$, $p=2$, which combinations of $f_{11}(z),f_{11}(2z),f_{11}(4z)$ lie in
$S_2(\Gamma_0(44))^+$?
\emph{Hint:} $W_N=W_M\sm{p^2}001$, $\sm p001W_N=p\,W_M\sm p001$, $\sm{p^2}001W_N=p^2W_M$. The eigenvectors are $h(pz)$ (eigenvalue $\varepsilon_M(h)$) and $h\pm p^2h(p^2z)$ (eigenvalues $\pm\varepsilon_M(h)$).
For $f_{11}$ ($\varepsilon_{11}(f_{11})=-1$) only $f_{11}(z)-4f_{11}(4z)$.
\end{exercise}

\begin{exercise}\label{ex:modforms:ex-fixed-points}
(a) Show that $z_0=i/\sqrt N$ is fixed by $w_N$, and that for $E=\C/\Lambda_{z_0}$ multiplication by $i\sqrt N$ is an endomorphism of degree $N$ with kernel $C=\langle\frac1N\rangle$; so the pair
$(E,C)$ satisfies $(E/C,E[N]/C)\simeq(E,C)$, as \Cref{thm:modforms:moduli}(iii) predicts. (b) Using Riemann-Hurwitz for $X_0(N)\to X_0(N)/w_N$ and
\Cref{prop:modforms:atkin-lehner-petersson}(ii), show that $w_N$ has $2g(X_0(N))+2-4\dim S_2(\Gamma_0(N))^+$ fixed points. Given $\dim S_2(\Gamma_0(121))^+=2$ (PARI), find the number of fixed
points of $w_{121}$ and compare with the class number $h(-484)=6$.
\emph{Answer:} (a) $i\sqrt N\cdot1=N\cdot\frac i{\sqrt N}$ and $i\sqrt N\cdot\frac i{\sqrt N}=-1$, so $i\sqrt N\Lambda_{z_0}\subseteq\Lambda_{z_0}$, of index $\abs{i\sqrt N}^2=N$, and $(i\sqrt N)^{-1}\cdot\frac i{\sqrt N}=\frac1N$ generates the
kernel. (b) $2\cdot6+2-4\cdot2=6$.
\end{exercise}

\section{Notes and references}\label{sec:modforms:notes}

The chapter follows Diamond and Shurman \cite{DiamondShurman} (Chapters 1 to 3 for the modular curves, the genus and the dimensions, Chapter 5 for the Petersson product and
the Atkin-Lehner operators, Section 6.5 for integrality), Serre's \emph{Course in Arithmetic} \cite[Chapter VII]{SerreCourse} for level one (fundamental domain, valence formula,
the structure of $M_*$, the bijectivity of $j$), and Shimura \cite{Shimura71} (Chapters 1 to 3) for the general theory. Other good sources are Miyake \cite{Miyake} (Section 4.6 for
the Atkin-Lehner operators), Koblitz \cite[Chapter III]{Koblitz}, Knapp \cite{Knapp} (Chapters VIII, IX and XI), Apostol \cite{Apostol90} (the $\eta$-function, Chapter 3) and Stein's
computational book \cite{Stein} (Chapter 9 for Sturm's bound). Atkin and Lehner's original paper is \cite{AtkinLehner70}; Sturm's theorem is \cite{Sturm87}; Ligozat's criterion
and his study of the curves of genus one are in \cite{Ligozat75}; the cusps of $X_0(N)$, their fields of definition and the fixed points of $w_N$ are in Ogg \cite{Ogg73}; the moduli
problem and the model of $X_0(N)$ over $\Q$ and over $\Z$ are in Deligne-Rapoport \cite{DeligneRapoport73} and Katz-Mazur \cite{KatzMazur85}. Deligne's proof of the
Ramanujan-Petersson conjecture is \cite{Deligne}. For the formula of Zagier and Watkins see \cite{Zagier85,Watkins02} and \Cref{ch:moddeg}; for the Rankin-Selberg formula
\cite{Rankin39,Shimura76} and \cite[Lemma 2.12]{DFG04}, and \Cref{ch:lfunctions}. The curves and their invariants are Cremona's \cite{Cremona97}, and the computations use
PARI/GP \cite{PARI}.

Pointers that were not checked against the printed books, and are given by chapter or section only: \cite{DiamondShurman} Sections 1.1 to 1.5, 2.1 to 2.4, 3.1, 3.5 to 3.9,
5.4, 5.5, 6.5, 7.7, 8.7, 8.8, and the theorem numbers 3.5.1, 3.6.1 quoted for the dimension formulas; \cite{Shimura71} Sections 1.5, 2.4, 2.6, 6.7 and Theorem 3.52 (integral bases);
\cite{SerreCourse} Chapter VII; \cite{Stein} Chapter 9; \cite{Miranda} Chapter III, Section 3 (quotients by finite groups); \cite{Apostol90} Chapter 3; \cite{Koblitz} Section III.2;
\cite{AtkinLehner70} Theorem 3 for $a_p=-\varepsilon_p$. Every number printed in the chapter (indices, cusps and widths, elliptic points, genera, dimensions, Sturm bounds, $q$-expansions,
Atkin-Lehner signs and matrices, Petersson norms, period areas, class numbers, congruences) was recomputed with PARI/GP 2.17 through \texttt{cypari2} and Cremona's tables for
conductors below $10{,}000$; the ranges of the checks are stated where the numbers appear.

The Petersson product has several normalisations in the literature (\Cref{conv:modforms:petersson-normalisations}); the book's is the unnormalised one, in which the
paper's Remark rem:degree holds as stated. The slash operator is normalised with $(\det\gamma)^{k/2}$, as in Shimura \cite{Shimura71} and Miyake \cite{Miyake} (not checked against
their texts); Diamond and Shurman use $(\det\gamma)^{k-1}$. Neither choice matters in weight $2$.
